% Faithful transcription — original French. Transcribed from the original first printing:
% N. H. Abel, "Recherches sur les fonctions elliptiques" (No. 12), Journal für die reine und
% angewandte Mathematik (Crelle's Journal), Zweiter Band, 2. Heft (Berlin: G. Reimer, 1827),
% pp. 101–181 (first part; the second part is in Band 3, 1828, pp. 160–190, not transcribed here).
% Scan: GDZ (Göttingen), PPN243919689_0002 (printed pp. 101–181 = image numbers 113–193).
%
% \origpage uses the PRINTED page numbers. Running heads, sheet signatures and catchwords are
% apparatus, not text, and are not transcribed. Printer's errors are reproduced as printed and
% marked with \ednote.
\documentclass{article}
\usepackage{readmasters}
\begin{document}

\origpage{101}

\section*{Recherches sur les fonctions elliptiques.}

(Par M. \emph{N. H. Abel}.)

Depuis longtems les fonctions logarithmiques, et les fonctions exponentielles et circulaires ont été les seules fonctions transcendantes, qui ont attiré l'attention des géomètres. Ce n'est que dans les derniers tems, qu'on a commencé à en considérer quelques autres. Parmi celles-ci il faut distinguer les fonctions, nommées elliptiques, tant pour leurs belles propriétés analytiques, que pour leur application dans les diverses branches des mathématiques. La première idée de ces fonctions a été donnée par l'immortel \emph{Euler}, en démontrant, que l'équation séparée
\[ \frac{\partial x}{\sqrt{(\alpha+\beta x+\gamma x^{2}+\delta x^{3}+\varepsilon x^{4})}}+\frac{\partial y}{\sqrt{(\alpha+\beta y+\gamma y^{2}+\delta y^{3}+\varepsilon y^{4})}}=0 \tag{1.} \]
est intégrable algébriquement. Après \emph{Euler}, \emph{Lagrange} y a ajouté quelque chose, en donnant son élégante théorie de la transformation de l'intégrale $\displaystyle\int\frac{R.\partial x}{\sqrt{[(1-p^{2}x^{2})(1-q^{2}x^{2})]}}$, ou $R$ est une fonction rationelle de $x$. Mais le premier et, si je ne me trompe, le seul, qui ait approfondi la nature de ces fonctions, est M. \emph{Legendre}, qui, d'abord dans un mémoire sur les fonctions elliptiques, et ensuite dans ses excellents exercices de mathématiques, a développé nombre de poropriétés\ednote{Printed poropriétés instead of propriétés; an apparent misprint.} élégantes de ces fonctions, et a montré leur application. Lors de la publication de cet ouvrage, rien n'a été ajouté à la théorie de M. \emph{Legendre}. Je crois, qu'on ne verra pas ici sans plaisir des recherches ultérieures sur ces fonctions.

En général on comprend sous la dénomination de fonctions elliptiques, toute fonction, comprise dans l'intégrale
\[ \int\frac{R\,\partial x}{\sqrt{(\alpha+\beta x+\gamma x^{2}+\delta x^{3}+\varepsilon x^{4})}}, \]
où $R$ est une fonction rationelle et $\alpha$, $\beta$, $\gamma$, $\delta$, $\varepsilon$ sont des quantités constantes et \emph{réelles}. M. \emph{Legendre} a démontré, que par des substitutions convenables on peut toujours ramener cette intégrale à la forme

\origpage{102}

\[ \int\frac{P\,\partial y}{\sqrt{(a+by^{2}+cy^{4})}}, \]
où $P$ est une fonction rationelle de $y^{2}$. Par les réductions convenables, cette intégrale peut être ensuite ramenée à la forme
\[ \int\frac{A+By^{2}}{C+Dy^{2}}\cdot\frac{\partial y}{\sqrt{(a+by^{2}+cy^{4})}}, \]
et celle-ci à:
\[ \int\frac{A+B\sin^{2}\theta}{C+D\sin^{2}\theta}\cdot\frac{\partial\theta}{\sqrt{(1-c^{2}\sin^{2}\theta)}}, \]
où $c$ est réel, et moindre que l'unité.

De là suit, que toute fonction elliptique peut être réduite à l'une des trois formes:
\[ \int\frac{\partial\theta}{\sqrt{(1-c^{2}\sin^{2}\theta)}};\ \int\partial\theta\sqrt{(1-c^{2}\sin^{2}\theta)};\ \int\frac{\partial\theta}{(1+n\sin^{2}\theta)\sqrt{(1-c^{2}\sin^{2}\theta)}}, \]
auxquelles M. \emph{Legendre} donne les noms de fonctions elliptiques de la première, seconde et troisième espèce. Ce sont ces trois fonctions, que M. \emph{Legendre} a considérées; surtout la première, qui a les propriétés les plus remarquables et les plus simples.

Je me propose, dans ce mémoire, de considérer la fonction inverse, c'est-à-dire la fonction $\Phi\alpha$, déterminée par les équations
\[ \begin{gathered} \alpha=\int\frac{\partial\theta}{\sqrt{(1-c^{2}\sin^{2}\theta)}}\ \text{et}\\ \sin\theta=\varphi(\alpha)=x. \end{gathered} \]
La dernière équation donne
\[ \partial\theta.\sqrt{(1-\sin^{2}\theta)}=\partial.\Phi\alpha=\partial x, \]
donc
\[ \alpha=\int_{0}\frac{\partial x}{\sqrt{[(1-x^{2})(1-c^{2}x^{2})]}}. \]
M. \emph{Legendre} suppose $c^{2}$ positif, mais j'ai remarqué, que les formules deviennent plus simples, en supposant $c^{2}$ négatif, $=-e^{2}$. De même j'écris pour plus de symmétrie $1-c^{2}x^{2}$ au lieu de $1-x^{2}$. En sorte que la fonction $\Phi\alpha=x$ sera donnée par l'équation
\[ \alpha=\int_{0}\frac{\partial x}{\sqrt{[(1-c^{2}x^{2})(1+e^{2}x^{2})]}}, \]
ou bien
\[ \Phi'\alpha=\partial\alpha.\sqrt{[(1-c^{2}\Phi^{2}\alpha)(1+e^{2}\Phi^{2}\alpha)]}.\]
\ednote{Printed $\partial\alpha$ before the radical; the factor is spurious, since the derivative is the square root alone.}
Pour abréger, j'introduis deux autres fonctions de $\alpha$, savoir:
\[ f\alpha=\sqrt{(1-c^{2}\Phi^{2}\alpha)};\ F\alpha=\sqrt{(1+e^{2}\Phi^{2}\alpha)}. \]

\origpage{103}

Plusieurs propriétés de ces fonctions se déduisent immédiatement des propriétés connues de la fonction elliptique de la première espèce, mais d'autres sont plus cachées. Par ex.\ on démontre, que les équations $\Phi\alpha=0$, $f\alpha=0$, $F\alpha=0$ ont un nombre infini de racines, qu'on peut trouver toutes. Une des propriétés les plus remarquables est, qu'on peut exprimer rationellement $\Phi(m\alpha)$, $f(m\alpha)$, $F(m\alpha)$ ($m$ étant un nombre entier) en $\Phi\alpha$, $f\alpha$, $F\alpha$. Aussi rien n'est plus facile, que de trouver $\Phi(m\alpha)$, $f(m\alpha)$, $F(m\alpha)$, lorsqu'on connoît $\Phi\alpha$, $f\alpha$, $F\alpha$; mais le problème inverse, savoir de déterminer $\Phi\alpha$, $f\alpha$, $F\alpha$ en $\Phi(m\alpha)$, $f(m\alpha)$, $F(m\alpha)$ est plus difficile, parcequ'il dépend d'une équation d'un degré plus élevé (savoir du degré $m^{2}$).

La solution de cette équation est l'objet principal de ce mémoire. D'abord on fera voir, comment ont\ednote{Printed ont instead of on; an apparent misprint.} peut trouver toutes les racines, au moyen des fonctions $\Phi,f,F$. On traitera ensuite de la solution algébrique de l'equation en question, et on parviendra à ce résultat remarquable, que $\varphi\left(\frac{\alpha}{m}\right)$; $f\left(\frac{\alpha}{m}\right)$; $F\left(\frac{\alpha}{m}\right)$ peuvent être exprimés en $\Phi\alpha$, $f\alpha$, $F\alpha$, au moyen d'une fonction, qui, par rapport à $\alpha$, ne contient d'autres irrationalités que des radicaux. Cela produit une classe très générale d'équations, qui sont resolubles algébriquement. Il est à remarquer, que les expressions des racines contiennent des quantités constantes, qui, en général, ne sont pas exprimables par des quantités algébriques. Ces quantités constantes dépendent d'une équation du degré $\mathrm{m}^{2}-1$. On fera voir comment, au moyen de fonctions algébriques, on peut en ramener la solution à celle d'une équation du degré $\mathrm{m}+1$. On donnera plusieurs expressions des fonctions $\Phi(2n+1)\alpha$; $f(2n+1)\alpha$; $F(2n+1)\alpha$ en fonctions de $\Phi\alpha$, $f\alpha$, $F\alpha$. On en déduira ensuite les valeurs de $\Phi\alpha$, $f\alpha$, $F\alpha$ en fonctions de $\alpha$. On démontrera, que ces fontions\ednote{Printed fontions instead of fonctions; an apparent misprint.} peuvent être décomposées en un nombre infini de facteurs, et même en une infinité de fractions partielles.

\section*{§. I.}

\emph{Propriétés fondamentales des fonctions} $\Phi\alpha$; $f\alpha$; $F\alpha$.

\subsection*{1.}

En supposant, que $\Phi\alpha=x$ (1.), on aura en vertu de ce qui précède:
\[ \alpha=\int_{0}^{x}\frac{\partial x}{\sqrt{[(1-c^{2}x^{2})(1+e^{2}x^{2})]}}. \tag{2.} \]

\origpage{104}
Par là on voit, que $\alpha$, considéré comme fonction de $x$, est positif depuis $x=0$, jusqu'à $x=\frac{1}{c}$. En faisant donc
\[ \frac{\omega}{2}=\int_{0}^{\frac{1}{c}}\frac{\partial x}{\sqrt{[(1-c^{2}x^{2})(1+e^{2}x^{2})]}}, \tag{3.} \]
il est évident, que $\Phi\alpha$ et\ednote{Printed et instead of est; an apparent misprint.} positif et va en augmentant depuis $\alpha=0$ jusqu'à $\alpha=\frac{\omega}{2}$, et qu'on aura
\[ \varphi(0)=0,\ \varphi\left(\frac{\omega}{2}\right)=\frac{1}{c}. \tag{4.} \]
Parceque $\alpha$ change de signe, lorsqu'on écrit $-x$ à la place de $x$; il en est de même de la fonction $\Phi\alpha$ par rapport à $\alpha$, et par conséquent on aura l'équation
\[ \varphi(-\alpha)=-\varphi(\alpha). \tag{5.} \]
En mettant dans (1.) $xi$ au lieu de $x$ (ou $i$, pour abréger, représente la quantité imaginaire $\sqrt{-1}$) et désignant la valeur de $\alpha$ par $\beta i$, il viendra
\[ xi=\Phi(\beta i)\ \text{et}\ \beta=\int_{0}^{x}\frac{\partial x}{\sqrt{[(1+c^{2}x^{2})(1-e^{2}x^{2})]}}. \tag{6.} \]
$\beta$ est réel et positif depuis $x=0$ jusqu'à $x=\frac{1}{e}$, donc en faisant
\[ \frac{\varpi}{2}=\int_{0}^{\frac{1}{e}}\frac{\partial x}{\sqrt{[(1-e^{2}x^{2})(1+c^{2}x^{2})]}}, \tag{7.} \]
$x$ sera positif, depuis $\beta=$ zéro jusqu'à $\beta=\frac{\varpi}{2}$, c'est-à-dire, la fonction $\frac{1}{i}\Phi(\beta i)$ sera positive entre les mêmes limites. En faisant $\beta=\alpha$ et $y=\frac{\varphi(\alpha i)}{i}$, on a
\[ \alpha=\int_{0}^{y}\frac{\partial y}{\sqrt{[(1-e^{2}y^{2})(1+c^{2}y^{2})]}}, \]
donc on voit, qu'en supposant $c$ au lieu de $e$ et $e$ au lieu de \uncertain{$\varepsilon$},\ednote{Printed as a curled e (apparently $\varepsilon$) at the last place; the sense requires $c$ (interchange of $c$ and $e$).}
\[ \frac{\varphi(\alpha i)}{i}\ \text{se changera en}\ \Phi\alpha. \]

Et parceque
\[ \begin{gathered} f\alpha=\sqrt{(1-c^{2}\Phi^{2}\alpha)},\\ F\alpha=\sqrt{(1+e^{2}\Phi^{2}\alpha)}, \end{gathered} \]
on voit, que par le changement de $c$ en $e$ et $e$ en $c$, $f(\alpha i)$ et $F(\alpha i)$ rentreront respectivement en $F\alpha$ et $f\alpha$. Enfin les équations (2.) et (7.) font voir, que par la même transformation, $\omega$ et $\varpi$ rentreront respectivement en $\varpi$ et $\omega$.

\origpage{105}

Suivant (7.) on aura $x=\frac{1}{e}$ pour $\beta=\frac{\varpi}{2}$, donc en vertu de l'équation $xi=\varphi(\beta i)$, il viendra
\[ \varphi\left(\frac{\varpi i}{2}\right)=i.\frac{1}{e}. \tag{8.} \]

\subsection*{2.}

En vertu de ce qui précède, on aura les valeurs de $\Phi\alpha$ pour toute valeur réelle de $\alpha$, comprise entre $-\frac{\omega}{2}$ et $+\frac{\omega}{2}$, et pour toute valeur imaginaire de la forme $\beta i$ de cette quantité, si $\beta$ est une quantité continue entre les limites $-\frac{\varpi}{2}$ et $+\frac{\varpi}{2}$. Il s'agit maintenant de trouver la valeur de cette fonction pour une valeur quelconque, réelle ou imaginaire, de la variable. Pour y parvenir, nous allons d'abord établir les propriétés fondamentales des fonctions $\Phi$, $f$ et $F$.

Parcequ'on a
\[ \begin{gathered} f^{2}\alpha=1-c^{2}\Phi^{2}\alpha,\\ F^{2}\alpha=1+e^{2}\Phi^{2}\alpha, \end{gathered} \]
on aura, en différentiant:
\[ \begin{gathered} f\alpha.f'\alpha=-c^{2}\Phi\alpha.\varphi'\alpha,\\ F\alpha.F'\alpha=e^{2}\Phi\alpha\,\varphi'\alpha, \end{gathered} \]
Or d'après (2.) on a
\[ \varphi'\alpha=\sqrt{[(1-c^{2}\Phi^{2}\alpha)(1+e^{2}\Phi^{2}\alpha)]}=f\alpha.F\alpha, \]
donc, en substituant cette valeur de $\Phi'\alpha$ dans les deux équations précédentes, on trouvera, que les fonctions $\Phi\alpha$, $f\alpha$, $F\alpha$ sont liées entre elles par les équations
\[ \left\{\begin{aligned} \varphi'\alpha&=f\alpha.F\alpha,\\ f'\alpha&=-c^{2}\Phi\alpha.F\alpha,\\ F'\alpha&=e^{2}\Phi\alpha.f\alpha. \end{aligned}\right. \tag{9.} \]
Cela posé, je dis, qu'en désignant par $\alpha$ et $\beta$ deux indéterminées, on aura
\[ \left\{\begin{aligned} \Phi(\alpha+\beta)&=\frac{\varphi\alpha.f\beta.F\beta+\varphi\beta.f\alpha.F\alpha}{1+e^{2}c^{2}.\varphi^{2}\alpha.\varphi^{2}\beta},\\ f(\alpha+\beta)&=\frac{f\alpha.f\beta-c^{2}.\varphi\alpha.\varphi\beta.F\alpha.F\beta}{1+e^{2}c^{2}\varphi^{2}\alpha.\varphi^{2}\beta},\\ F(\alpha+\beta)&=\frac{F\alpha.F\beta+e^{2}.\varphi\alpha.\varphi\beta.f\alpha.f\beta}{1+e^{2}c^{2}\varphi^{2}\alpha.\varphi^{2}\beta}. \end{aligned}\right. \tag{10.} \]

Ces formules peuvent être déduites sur le champ des propriétés connues des fonctions elliptiques${}^{*)}$; mais on peut aussi les vérifier aisément de la manière suivante.

\textbf{*)} \emph{Legendre} Exercices de calcul intégral.

\origpage{106}

En désignant par $r$ le second membre de la première des équations (10.), on aura, en differentiant par rapport à $\alpha$:
\[ \begin{aligned} \left(\frac{\partial r}{\partial\alpha}\right)&=\frac{\{\varphi'\alpha.f\beta.F\beta+\varphi\beta.F\alpha.f'\alpha+\varphi\beta.f\alpha.F'\alpha\}}{1+e^{2}c^{2}.\varphi^{2}\alpha.\varphi^{2}\beta}\\ &\quad-\frac{(\varphi\alpha.f\beta F\beta+\varphi\beta.f\alpha\ F\alpha).2e^{2}c^{2}\varphi\alpha.\varphi^{2}\beta.\varphi'\alpha}{(1+e^{2}c^{2}\varphi^{2}\alpha.\varphi^{2}\beta)^{2}}. \end{aligned} \]
En substituant pour $\Phi'\alpha$, $f'\alpha$, $F'\alpha$ leurs valeurs données par les équations (9.), il viendra
\[ \begin{aligned} \left(\frac{\partial r}{\partial\alpha}\right)&=\frac{f\alpha.F\alpha.f\beta.F\beta}{1+e^{2}c^{2}\varphi^{2}\alpha.\varphi^{2}\beta}-\frac{2e^{2}c^{2}\varphi^{2}\alpha.\varphi^{2}\beta.f\alpha.f\beta.F\alpha.F\beta}{(1+e^{2}c^{2}\varphi^{2}\alpha.\varphi^{2}\beta)^{2}}\\ &\quad+\frac{\varphi\alpha.\varphi\beta.\{1+e^{2}c^{2}\varphi^{2}\alpha\varphi^{2}\beta\}\{-c^{2}F^{2}\alpha+e^{2}f^{2}\alpha\}-2e^{2}c^{2}\varphi\alpha\,\varphi\beta.\varphi^{2}\beta.f^{2}\alpha.F^{2}\alpha}{(1+e^{2}c^{2}\varphi^{2}\alpha\varphi^{2}\beta)^{2}}, \end{aligned} \]
d'où, en substituant pour $f^{2}\alpha$ et $F^{2}\alpha$ leurs valeurs: $1-c^{2}\Phi^{2}\alpha$, $1+e^{2}\Phi^{2}\alpha$, et en réduisant, on tire
\[ \left(\frac{\partial r}{\partial\alpha}\right)=\frac{(1-e^{2}c^{2}\varphi^{2}\alpha.\varphi^{2}\beta)\{(e^{2}-c^{2})+\varphi\alpha.\varphi\beta+f\alpha.f\beta.F\alpha.F\beta\}-2e^{2}c^{2}\varphi\alpha\,\varphi\beta(\varphi^{2}\alpha+\varphi^{2}\beta)}{(1+e^{2}c^{2}\varphi^{2}\alpha.\varphi^{2}\beta)^{2}}. \]
\ednote{Printed with a plus sign (and a stray hairline before it) between $(e^{2}-c^{2})$ and $\varphi\alpha.\varphi\beta$; a product $(e^{2}-c^{2})\varphi\alpha.\varphi\beta$ is expected.}
Maintenant $\alpha$ et $\beta$ entrent symétriquement dans l'expression de $r$; donc on aura la valeur de $\left(\frac{\partial r}{\partial\beta}\right)$, en permutant $\alpha$ et $\beta$ dans la valeur de $\left(\frac{\partial r}{\partial\alpha}\right)$. Or par cela l'expression de $\left(\frac{\partial r}{\partial\alpha}\right)$ ne change pas de valeur, donc on aura
\[ \left(\frac{\partial r}{\partial\alpha}\right)=\left(\frac{\partial r}{\partial\beta}\right). \]

Cette équation aux différentielles partielles fait voir que $r$ est fonction de $\alpha+\beta$; donc on aura
\[ r=\psi(\alpha+\beta). \]
La forme de la fonction $\psi$ se trouvera, en donnant à $\beta$ une valeur particulière. En supposant par ex.\ $\beta=0$, et remarquant que $\Phi(0)=0$, $f(0)=1$, $F(0)=1$, les deux valeurs de $r$ deviendront
\[ r=\Phi(\alpha)\ \text{et}\ r=\psi(\alpha), \]
donc
\[ \psi(\alpha)=\Phi(\alpha), \]
et de là
\[ r=\psi(\alpha+\beta)=\Phi(\alpha+\beta). \]
La première des formules (10.) a donc lieu effectivement.

De la même manière on verifira les deux autres formules.

\origpage{107}

\subsection*{3.}

Des formules (10.) on peut en déduire une foule d'autres. Je vais rapporter quelques-unes des plus remarquables: Pour abréger je fais
\[ 1+e^{2}c^{2}\Phi^{2}\alpha.\Phi^{2}\beta=R. \tag{11.} \]

En changeant d'abord le signe de $\beta$, on obtiendra
\[ \left\{\begin{aligned} \Phi(\alpha+\beta)+\Phi(\alpha-\beta)&=\frac{2\varphi\alpha.f\beta.F\beta}{R},\\ \Phi(\alpha+\beta)-\Phi(\alpha-\beta)&=\frac{2\varphi\beta.f\alpha.F\alpha}{R},\\ f(\alpha+\beta)+f(\alpha-\beta)&=\frac{2f\alpha.f\beta}{R},\\ f(\alpha+\beta)-f(\alpha-\beta)&=\frac{-2c^{2}.\varphi\alpha.\varphi\beta.F\alpha.F\beta}{R},\\ F(\alpha+\beta)+F(\alpha-\beta)&=\frac{2F\alpha.F\beta}{R},\\ F(\alpha+\beta)-F(\alpha-\beta)&=\frac{2e^{2}.\varphi\alpha.\varphi\beta.f\alpha.f\beta}{R}. \end{aligned}\right. \tag{12.} \]

En formant le produit de $\Phi(\alpha+\beta)$ et $\Phi(\alpha-\beta)$, on trouvera
\[ \begin{aligned} \Phi(\alpha+\beta).\Phi(\alpha-\beta)&=\frac{\varphi\alpha.f\beta.F\beta+\varphi\beta.f\alpha.F\alpha}{R}.\frac{\varphi\alpha.f\beta.F\beta-\varphi\beta.f\alpha.F\alpha}{R}\\ &=\frac{\varphi^{2}\alpha.f^{2}\beta.F^{2}\beta-\varphi^{2}\beta.f^{2}\alpha.F^{2}\alpha}{R^{2}}, \end{aligned} \]
ou, en substituant les valeurs de $f^{2}\beta$, $F^{2}\beta$, $f^{2}\alpha$, $F^{2}\alpha$ dans $\Phi\beta$ et $\Phi\alpha$:
\[ \begin{aligned} \Phi(\alpha+\beta).\Phi(\alpha-\beta)&=\frac{\varphi^{2}\alpha-\varphi^{2}\beta-e^{2}c^{2}\varphi^{2}\alpha.\varphi^{4}\beta+e^{2}c^{2}\varphi^{2}\beta.\varphi^{4}\alpha}{R^{2}}\\ &=\frac{(\varphi^{2}\alpha-\varphi^{2}\beta)(1+e^{2}c^{2}\varphi^{2}\alpha.\varphi^{2}\beta)}{R^{2}}; \end{aligned} \]
or $R=1+e^{2}c^{2}\Phi^{2}\alpha.\Phi^{2}\beta$, donc
\[ \Phi(\alpha+\beta).\Phi(\alpha-\beta)=\frac{\varphi^{2}\alpha-\varphi^{2}\beta}{R}. \tag{13.} \]
On trouvera de même
\[ \left\{\begin{aligned} f(\alpha+\beta).f(\alpha-\beta)&=\frac{f^{2}\alpha-c^{2}\varphi^{2}\beta.F^{2}\alpha}{R}=\frac{f^{2}\beta-c^{2}\varphi^{2}\alpha.F^{2}\beta}{R}\\ &=\frac{1-c^{2}\varphi^{2}\alpha-c^{2}\varphi^{2}\beta-c^{2}e^{2}\varphi^{2}\alpha.\varphi^{2}\beta}{R}=\frac{f^{2}\alpha.f^{2}\beta-c^{2}(c^{2}+e^{2})\varphi^{2}\alpha.\varphi^{2}\beta}{R},\\ F(\alpha+\beta).F(\alpha-\beta)&=\frac{F^{2}\alpha+e^{2}\varphi^{2}\beta.f^{2}\alpha}{R}=\frac{F^{2}\beta+e^{2}\varphi^{2}\alpha.f^{2}\beta}{R}\\ &=\frac{1+e^{2}\varphi^{2}\alpha+e^{2}\varphi^{2}\beta-e^{2}c^{2}\varphi^{2}\alpha.\varphi^{2}\beta}{R}=\frac{F^{2}\alpha.F^{2}\beta-e^{2}(c^{2}+e^{2})\varphi^{2}\alpha.\varphi^{2}\beta}{R}. \end{aligned}\right. \tag{14.} \]

\origpage{108}

\subsection*{4.}

En faisant dans les formules (10.) $\beta=\pm\frac{\omega}{2}$, $\beta=\pm\frac{\varpi}{2}i$, et remarquant, que $f\left(\pm\frac{\omega}{2}\right)=0$, $F\left(\pm\frac{\varpi}{2}i\right)=0$, on aura
\[ \left\{\begin{aligned} &\varphi\left(\alpha\pm\frac{\omega}{2}\right)=\pm\varphi\frac{\omega}{2}.\frac{f\alpha}{F\alpha};\ f\left(\alpha\pm\frac{\omega}{2}\right)=\mp\frac{F\frac{\omega}{2}}{\varphi\frac{\omega}{2}}.\frac{\varphi\alpha}{F\alpha};\\ &F\left(\alpha\pm\frac{\omega}{2}\right)=\frac{F\frac{\omega}{2}}{F\alpha};\\ &\varphi\left(\alpha\pm\frac{\varpi}{2}i\right)=\pm\varphi\frac{\varpi}{2}i.\frac{F\alpha}{f\alpha};\ F\left(\alpha\pm\frac{\varpi}{2}i\right)=\mp\frac{f\frac{\varpi}{2}i}{\varphi\frac{\varpi}{2}i}.\frac{\varphi\alpha}{f\alpha};\\ &f\left(\alpha\pm\frac{\varpi}{2}i\right)=\frac{f\frac{\varpi}{2}i}{f\alpha}; \end{aligned}\right. \tag{15.} \]
ou bien:
\[ \left\{\begin{aligned} &\varphi\left(\alpha\pm\frac{\omega}{2}\right)=\pm\frac{1}{c}.\frac{f\alpha}{F\alpha};\ f\left(\alpha\pm\frac{\omega}{2}\right)=\mp\sqrt{(e^{2}+c^{2})}.\frac{\varphi\alpha}{F\alpha};\\ &F\left(\alpha\pm\frac{\omega}{2}\right)=\frac{\sqrt{(e^{2}+c^{2})}}{c}.\frac{1}{F\alpha};\\ &\varphi\left(\alpha\pm\frac{\varpi}{2}i\right)=\pm\frac{i}{e}.\frac{F\alpha}{f\alpha};\ F\left(\alpha\pm\frac{\varpi}{2}i\right)=\pm i\sqrt{(e^{2}+c^{2})}.\frac{\varphi\alpha}{f\alpha};\\ &f\left(\alpha\pm\frac{\varpi}{2}i\right)=\frac{\sqrt{(e^{2}+c^{2})}}{e}.\frac{1}{f\alpha}. \end{aligned}\right. \tag{16.} \]

De là on tire sur le champ:
\[ \left\{\begin{aligned} &\varphi\left(\frac{\omega}{2}+\alpha\right)=\Phi\left(\frac{\omega}{2}-\alpha\right);\ f\left(\frac{\omega}{2}+\alpha\right)=-f\left(\frac{\omega}{2}-\alpha\right);\\ &F\left(\frac{\omega}{2}+\alpha\right)=F\left(\frac{\omega}{2}-\alpha\right);\\ &\varphi\left(\frac{\varpi}{2}i+\alpha\right)=\Phi\left(\frac{\varpi}{2}i-\alpha\right);\ F\left(\frac{\varpi}{2}i+\alpha\right)=-F\left(\frac{\varpi}{2}i-\alpha\right);\\ &f\left(\frac{\varpi}{2}i+\alpha\right)=f\left(\frac{\varpi}{2}i-\alpha\right); \end{aligned}\right. \tag{17.} \]
\[ \varphi\left(\alpha\pm\frac{\omega}{2}\right).\varphi\left(\alpha+\frac{\varpi}{2}i\right)=\pm\frac{i}{ce};\ F\left(\alpha\pm\frac{\omega}{2}\right).F\alpha=\frac{b}{c};\ f\left(\alpha\pm\frac{\varpi}{2}i\right).f\alpha=\frac{b}{e}. \tag{18.} \]
En faisant $\alpha=\frac{\omega}{2}$ et $\frac{\varpi}{2}i$, on trouve de là
\[ \Phi\left(\frac{\omega}{2}+\frac{\varpi}{2}i\right)=\frac{1}{0},\ f\left(\frac{\omega}{2}+\frac{\varpi}{2}i\right)=\frac{1}{0},\ F\left(\frac{\omega}{2}+\frac{\varpi}{2}i\right)=\frac{1}{0}. \]

\origpage{109}
Ensuite les équations (17.), en y mettant dans les trois premières $\alpha+\frac{\omega}{2}$ au lieu de $\alpha$, et dans les trois dernières $\alpha+\frac{\varpi}{2}i$ au lieu de $\alpha$, donnent les suivantes
\[ \left\{\begin{aligned} &\varphi(\alpha+\omega)=-\Phi\alpha;\ f(\alpha+\omega)=-f\alpha;\ F(\alpha+\omega)=F\alpha;\\ &\varphi(\alpha+\varpi i)=-\Phi\alpha;\ f(\alpha+\varpi i)=f\alpha;\ F(\alpha+\varpi i)=-F\alpha; \end{aligned}\right. \tag{19.} \]
et en mettant $\alpha+\omega$ et $\alpha+\varpi i$ au lieu de $\alpha$:
\[ \left\{\begin{aligned} &\Phi(2\omega+\alpha)=\Phi\alpha;\ \Phi(2\varpi i+\alpha)=\Phi\alpha;\ \Phi(\omega+\varpi i+\alpha)=\Phi\alpha;\\ &f(2\omega+\alpha)=f\alpha;\ f(\varpi i+\alpha)=f\alpha;\\ &F(\omega+\alpha)=F\alpha;\ F(2\varpi i+\alpha)=F\alpha. \end{aligned}\right. \tag{20.} \]

Ces équations font voir, que les fonctions $\Phi\alpha$, $f\alpha$, $F\alpha$ sont des fonctions \emph{périodiques}. On en déduira sans peine les suivantes, où $m$ et $n$ sont deux nombres entiers positifs ou négatifs:
\[ \left\{\begin{aligned} &\varphi((m+n)\omega+(m-n)\varpi i+\alpha)=\varphi\alpha;\ \varphi((m+n)\omega+(m-n+1)\varpi i+\alpha)=-\varphi\alpha;\\ &f(2m\omega+n\varpi i+\alpha)=f\alpha;\ f((2m+1)\omega+n\varpi i+\alpha)=-f\alpha;\\ &F(m\omega+2n\varpi i+\alpha)=+F\alpha;\ F(m\omega+(2n+1)\varpi i+\alpha)=-F\alpha. \end{aligned}\right. \tag{21.} \]
Ces formules peuvent aussi s'écrire comme il suit:
\[ \left\{\begin{aligned} &\Phi(m\omega+n\varpi i\pm\alpha)=\pm(-1)^{m+n}.\Phi\alpha,\\ &f(m\omega+n\varpi i\pm\alpha)=(-1)^{m}.f\alpha,\\ &F(m\omega+n\varpi i\pm\alpha)=(-)^{n}.F\alpha. \end{aligned}\right. \tag{22.} \]
\ednote{In the last line of this system the base of the exponent is printed as a bare minus sign, without the digit 1; an apparent misprint for $(-1)^{n}$.}
On peut remarquer comme cas particuliers:
\[ \left\{\begin{aligned} &\Phi(m\omega\pm\alpha)=\pm(-1)^{m}.\Phi\alpha;\ \Phi(n\varpi i\pm\alpha)=\pm(-1)^{n}.\Phi\alpha;\\ &f(m\omega\pm\alpha)=(-1)^{m}.f\alpha;\ f(n\varpi i\pm\alpha)=f\alpha;\\ &F(m\omega\pm\alpha)=F\alpha;\ F(n\varpi i\pm\alpha)=(-1)^{n}.F\alpha. \end{aligned}\right. \tag{22$'$.} \]

\subsection*{5.}

Les formules qu'on vient d'établir, font voir qu'on aura les valeurs des fonctions $\Phi\alpha$; $f\alpha$; $F\alpha$ pour toutes les valeurs réelles ou imaginaires de la variable, lorsqu'on les connoît pour les valeurs réelles de cette quantité, comprises entre $\frac{\omega}{2}$ et $-\frac{\omega}{2}$ et pour les valeurs imaginaires de la forme $\beta i$, où $\beta$ est comprise entre $\frac{\varpi}{2}$ et $-\frac{\varpi}{2}$.

En effet, supposons qu'on demande la valeur des fonctions $\Phi(\alpha+\beta i)$, $f(\alpha+\beta i)$, $F(\alpha+\beta i)$, où $\alpha$ et $\beta$ sont des quantités réelles quelconques. En mettant dans les formules (10.) $\beta i$ à la place de $\beta$, il est clair, qu'on aura les trois fonctions dont il s'agit, exprimées par les fonctions: $\Phi\alpha$; $f\alpha$; $F\alpha$; $\Phi(\beta i)$; $f(\beta i)$; $F(\beta i)$. Il ne reste donc, qu'à déterminer ces dernières.

\origpage{110}
Or, quelles que soient les valeurs de $\alpha$ et $\beta$, on peut toujours trouver deux nombres entiers $m$ et $n$, tels que $\alpha=m.\omega\pm\alpha'$, $\beta=n\varpi\pm\beta'$, où $\alpha'$ est une quantité comprise entre $0$ et $+\frac{\omega}{2}$, et $\beta'$ entre $0$ et $+\frac{\varpi}{2}$. Donc on aura, en vertu des équations (22.), en substituant les valeurs précédentes de $\alpha$ et $\beta$:
\[ \begin{aligned} \Phi(\alpha)&=\varphi(m\omega\pm\alpha')=\pm(-1)^{m}.\Phi\alpha',\\ f(\alpha)&=f(m\omega\pm\alpha')=(-1)^{m}.f\alpha',\\ F(\alpha)&=F(m\omega\pm\alpha')=F\alpha',\\ \Phi(\beta i)&=\varphi(+n\varpi i\pm\beta'i)=\pm(-1)^{n}.\varphi(\beta'i),\\ f(\beta i)&=f(n\varpi i\pm\beta'i)=f(\beta'i),\\ F(\beta i)&=F(n\varpi i\pm\beta'i)=(-1)^{n}.F(\beta'i). \end{aligned} \]
Donc les fonctions $\Phi\alpha$, $f\alpha$, $F\alpha$, $\Phi(\beta i)$, $f(\beta i)$, $F(\beta i)$ seront exprimées comme on vient de le dire et par suite aussi les fonctions $\Phi(\alpha+\beta i)$; $f(\alpha+\beta i)$; $F(\alpha+\beta i)$.

Nous avons vu précédemment, que $\Phi\alpha$ est réel depuis $\alpha=-\frac{\omega}{2}$ jusqu'à $\alpha=+\frac{\omega}{2}$, et que $\frac{\varphi(\alpha i)}{i}$ est réel depuis $\alpha=-\frac{\varpi}{2}$ jusqu'à $\alpha=+\frac{\varpi}{2}$.

Donc en vertu des équations (22.) il est clair:

1) que $\Phi(\alpha)$ et $\frac{\varphi(\alpha i)}{i}$ sont réels pour toute valeur réelle de $\alpha$; $\Phi\alpha$ est compris entre $-\frac{1}{c}$ et $+\frac{1}{c}$, et $\frac{\varphi(\alpha i)}{i}$ entre $-\frac{1}{e}$ et $+\frac{1}{e}$;

2) que $\Phi(\alpha)$ s'évanouit pour $\alpha=m\omega$, et $\frac{\varphi(\alpha i)}{i}$ pour $\alpha=m\varpi$, $m$ étant un nombre entier positif ou négatif; mais $\Phi(\alpha)$ n'est pas nul pour aucune autre valeur réelle de $\alpha$.

En remarquant, que $f\alpha=\sqrt{(1-c^{2}\Phi^{2}\alpha)}$, $F\alpha=\sqrt{(1+e^{2}\Phi^{2}\alpha)}$, il suit de ce que nous venons de dire:

1) que les fonctions $f(\alpha)$; $F(\alpha)$; $f(\alpha i)$; $F(\alpha i)$ sont réelles pour toute valeur de $\alpha$;

2) que $f(\alpha)$ est compris entre les limites $-1$ et $+1$ et $F(\alpha)$ entre les limites $+1$ et $+\sqrt{\left(1+\frac{e^{2}}{c^{2}}\right)}$, ensorte que $F\alpha$ est positif pour toute valeur réelle de $\alpha$;

3) que $f(\alpha i)$ est positif et compris entre les limites $+1$ et $\sqrt{\left(1+\frac{c^{2}}{e^{2}}\right)}$ et $F(\alpha i)$ entre les limites $-1$ et $+1$ pour toute valeur réelle de $\alpha$;

\origpage{111}

4) que $f(\alpha)$ s'évanouit pour $\alpha=(m+\frac{1}{2})\omega$ et $F(i\alpha)$ pour $\alpha=(m+\frac{1}{2})\varpi$; mais pour nulle autre valeur de $\alpha$.

On remarquera ce qui suit comme corollaires qu'on tire des formules (22.):

1) Soit $\alpha=0$. Dans ce cas, en remarquant, que $\Phi(0)=0$, $f(0)=1$, $F(0)=1$, on aura
\[ \left\{\begin{aligned} \varphi(m\omega+n\varpi i)&=0,\\ f(m\omega+n\varpi i)&=(-1)^{m},\\ F(m\omega+n\varpi i)&=-1. \end{aligned}\right. \tag{23.} \]
\ednote{The last line is printed as $-1$ over a blotted impression; by formula (22.) one would expect $(-1)^{n}$.}

2) Soit $\alpha=\frac{\omega}{2}$. En vertu des équations:
\[ \Phi\left(\frac{\omega}{2}\right)=\frac{1}{c};\ f\left(\frac{\omega}{2}\right)=0;\ F\left(\frac{\omega}{2}\right)=\frac{b}{c}; \]
on aura
\[ \left\{\begin{aligned} \Phi((m+\tfrac{1}{2})\omega+n\varpi i)&=(-1)^{m+n}.\frac{1}{c},\\ f((m+\tfrac{1}{2})\omega+n\varpi i)&=0,\\ F((m+\tfrac{1}{2})\omega+n\varpi i)&=(-1)^{n}.\frac{b}{c}. \end{aligned}\right. \tag{24.} \]

3) Soit $\alpha=\frac{\varpi}{2}i$. En vertu des équations
\[ \Phi\left(\frac{\varpi}{2}i\right)=\frac{i}{e};\ f\left(\frac{\varpi}{2}i\right)=\frac{b}{e};\ F\left(\frac{\varpi}{2}i\right)=0; \]
on aura:
\[ \left\{\begin{aligned} \Phi(m\omega+(n+\tfrac{1}{2})\varpi i)&=(-1)^{m+n}\frac{i}{e},\\ f(m\omega+(n+\tfrac{1}{2})\varpi i)&=(-1)^{m}.\frac{b}{e},\\ F(m\omega+(n+\tfrac{1}{2})\varpi i)&=0. \end{aligned}\right. \tag{25.} \]

4) Soit $\alpha=\frac{\omega}{2}+\frac{\varpi}{2}i$. En vertu des équations ci dessus on aura
\[ \left\{\begin{aligned} \Phi((m+\tfrac{1}{2})\omega+(n+\tfrac{1}{2})\varpi i)&=\frac{1}{0},\\ f((m+\tfrac{1}{2})\omega+(n+\tfrac{1}{2})\varpi i)&=\frac{1}{0},\\ F((m+\tfrac{1}{2})\omega+(n+\tfrac{1}{2})\varpi i)&=\frac{1}{0}. \end{aligned}\right. \tag{26.} \]

\subsection*{6.}

Les équations (23.), (24.), (25.) font voir, que la fonction $\Phi(\alpha)$ s'évanouit toutes les fois que $\alpha$ est de la forme $\alpha=m\omega+n\varpi i$; que $f\alpha$ s'évanouit toutes les fois que $\alpha$ est de la forme $\alpha=(m+\frac{1}{2})\omega+n\varpi i$, et que $F\alpha$ s'évanouit toutes les fois que $\alpha$ est de la forme $\alpha=m\omega+(n+\frac{1}{2})\varpi i$.

\origpage{112}
Or je dis que pour toute autre valeur de $\alpha$, les fonctions $\Phi\alpha$, $f\alpha$, $F\alpha$ auront nécessairement une valeur différente de zéro.

Supposons en effet, qu'on ait
\[ \Phi(\alpha+\beta i)=0, \]
$\alpha$ et $\beta$ étant des quantités réelles. En vertu de la première des formules (10.), cette équation peut s'écrire comme il suit:
\[ \frac{\varphi\alpha.f(\beta i).F(\beta i)+\varphi(\beta i).f\alpha.F\alpha}{1+e^{2}c^{2}\varphi^{2}\alpha.\varphi^{2}(\beta i)}=0. \]

Maintenant les quantités $\Phi\alpha$, $f(\beta i)$, $F(\beta i)$ sont réelles et $\Phi(\beta i)$ est de la forme $i.A$, où $A$ est réel; donc cette équation ne peut pas subsister à moins qu'on n'ait séparément:
\[ \Phi(\alpha).f(\beta i).F(\beta i)=0;\ \Phi(\beta i).f\alpha.F\alpha=0. \]
Ces équations ne peuvent être satisfaites, que de deux manières, savoir en faisant
\[ \Phi(\alpha)=0,\ \Phi(\beta i)=0, \]
ou
\[ f(\beta i).F(\beta i)=0,\ f\alpha.F\alpha=0. \]

Les deux premières équations donnent $\alpha=m\omega$; $\beta=n\varpi$.

Les deux dernières, en remarquant que $F\alpha$ et $f(\beta i)$ ne peuvent jamais s'évanouir, donnent
\[ f\alpha=0,\ F(\beta i)=0, \]
d'où
\[ \alpha=(m+\tfrac{1}{2})\omega,\ \beta=(n+\tfrac{1}{2})\varpi. \]
Mais pour ces valeurs de $\alpha$ et $\beta$ la valeur de $\Phi(\alpha+\beta i)$ deviendra infinie; donc les seules valeurs de $\alpha$ et $\beta$ sont $\alpha=m\omega$ et $\beta=n\varpi$, et par conséquent toutes les racines de l'équation
\[ \Phi(x)=0, \]
peuvent être représentées par
\[ x=m\omega+n\varpi i. \tag{27.} \]
De la même manière on trouvera, que toutes les racines de l'équation
\[ f(x)=0, \]
peuvent être représentées par
\[ x=(m+\tfrac{1}{2})\omega+n\varpi i, \tag{28.} \]
et celles de l'équation
\[ F(x)=0, \]
par
\[ x=m\omega+(n+\tfrac{1}{2})\varpi i. \tag{29.} \]

\origpage{113}

\subsection*{7.}

Les formules (26.) font voir, qu'on satisfait aux trois équations
\[ \Phi(x)=\frac{1}{0},\ f(x)=\frac{1}{0},\ F(x)=\frac{1}{0}, \]
en donnant à $x$ une des valeurs de la forme
\[ x=(m+\tfrac{1}{2})\omega+(n+\tfrac{1}{2})\varpi i. \tag{30.} \]

Or on peut démontrer que les équations en question n'ont pas d'autres racines.

En effet, ayant
\[ \begin{gathered} F(x)=\frac{b}{c}.\frac{1}{F\left(x-\frac{\omega}{2}\right)};\ \Phi(x)=\frac{-i}{ec}.\frac{1}{\varphi\left(x-\frac{\omega}{2}-\frac{\varpi}{2}i\right)};\\ f(x)=\frac{b}{e}.\frac{1}{f\left(x-\frac{\varpi}{2}i\right)}; \end{gathered} \]
les équations en question entraîneront celles-ci:
\[ \varphi\left(x-\frac{\omega}{2}-\frac{\varpi}{2}i\right)=0;\ f\left(x-\frac{\varpi}{2}i\right)=0;\ F\left(x-\frac{\omega}{2}\right)=0, \]
mais en vertu de ce qu'on vient de voir dans le numéro précédent, les équations donnent respectivement:
\[ \begin{gathered} x-\frac{\omega}{2}-\frac{\varpi}{2}i=m\omega+n\varpi i;\ x-\frac{\varpi}{2}i=(m+\tfrac{1}{2})\omega+n\varpi i;\\ x-\frac{\omega}{2}=m\omega+(n+\tfrac{1}{2})\varpi i; \end{gathered} \]
c'est-à-dire: on aura pour les trois équations:
\[ x=(m+\tfrac{1}{2})\omega+(n+\tfrac{1}{2})\varpi i, \]
c.\ q.\ f.\ d.

\subsection*{8.}

Ayant trouvé comme ci dessus toutes les racines des équations
\[ \begin{gathered} \Phi(x)=0;\ f(x)=0;\ F(x)=0;\\ \Phi(x)=\frac{1}{0};\ f(x)=\frac{1}{0};\ F(x)=\frac{1}{0}; \end{gathered} \]
je vais maintenant chercher les racines des équations plus générales.
\[ \Phi(x)=\Phi a,\ f(x)=fa,\ F(x)=Fa, \]
où $a$ est une quantité quelconque réelle ou imaginaire.

Considérons d'abord l'équation
\[ \Phi(x)-\Phi a=0. \]
En faisant dans la seconde des formules (12.)
\[ \alpha=\frac{x+a}{2},\ \beta=\frac{x-a}{2}, \]
on trouvera
\[ \Phi x-\Phi a=\frac{2\varphi\left(\frac{x-a}{2}\right).f\left(\frac{x+a}{2}\right).F\left(\frac{x+a}{2}\right)}{1+e^{2}c^{2}.\varphi^{2}\left(\frac{x+a}{2}\right).\varphi^{2}\left(\frac{x-a}{2}\right)}=0. \]

\origpage{114}
Cette équation ne peut subsister que dans l'un des cinq cas suivans:

1. si $\Phi\left(\frac{x-a}{2}\right)=0$, d'où $x=a+2m\omega+2n\varpi i$,

2. si $f\left(\frac{x+a}{2}\right)=0$, d'où $x=-a+(2m+1)\omega+2n\varpi i$,

3. si $F\left(\frac{x+a}{2}\right)=0$, d'où $x=-a+2m\omega+(2n+1)\varpi i$,

4. si $\Phi\left(\frac{x-a}{2}\right)=\frac{1}{0}$, d'où $x=a+(2m+1)\omega+(2n+1)\varpi i$,

5. si $\Phi\left(\frac{x+a}{2}\right)=\frac{1}{0}$, d'où $x=-a+(2m+1)\omega+(2n+1)\varpi i$.

La résolution de ces cinq équations est contenue dans les formules (27.), (28.), (29.).

Des valeurs trouvées de $x$ il faut rejeter celles, que donne la formule
\[ x=-a+(2m+1)\omega+(2n+1)\varpi i, \]
car une telle valeur de $x$ donne en vertu de (22.):
\[ \Phi x=-\Phi a, \]
tandis qu'on doit avoir $\Phi x=\Phi a$; mais les autres valeurs de $x$, exprimées par les quatre premières formules, peuvent être admises. Elles sont, comme on voit, contenues dans la seule formule:
\[ x=(-1)^{m+n}.a+m\omega+n\varpi i. \tag{31.} \]
Telle est donc l'expression générale de toutes les racines de l'équation
\[ \Phi x=\Phi a. \]
De la même manière on trouvera, que toutes les racines de l'équation
\[ fx=fa \]
sont représentées par la formule
\[ x=\pm a+2m\omega+n\varpi i, \tag{32.} \]
et toutes celles de l'équation
\[ Fx=Fa, \]
par la formule
\[ x=\pm a+m\omega+2n\varpi i. \tag{33.} \]

\section*{§. II.}

\emph{Formules, qui donnent les valeurs de} $\Phi(n\alpha)$, $f(n\alpha)$, $F(n\alpha)$ \emph{exprimées en fonctions rationelles de} $\Phi\alpha$, $f\alpha$, $F\alpha$.

\subsection*{9.}

Reprenons les formules (12.). En faisant dans la $1{}^{\mathrm{e}}$, $3{}^{\mathrm{e}}$ et $5{}^{\mathrm{e}}$ $\alpha=n\beta$, il viendra:

\origpage{115}
\[ \left\{\begin{aligned} \Phi(n+1)\beta&=-\Phi(n-1)\beta+\frac{2\varphi(n\beta).f\beta.F\beta}{R},\\ f(n+1)\beta&=-f(n-1)\beta+\frac{2f(n\beta).f\beta}{R},\\ F(n+1)\beta&=-F(n-1)\beta+\frac{2F(n\beta).F\beta}{R}, \end{aligned}\right. \tag{34.} \]
où $R=1+c^{2}e^{2}.\Phi^{2}(n\beta).\Phi^{2}\beta$.

Ces formules donnent la valeur de $\Phi(n+1)\beta$ en $\Phi(n-1)\beta$ et $\Phi(n\beta)$; celle de $f(n+1)\beta$ en $f(n-1)\beta$ et $f(n\beta)$, et celle de $F(n+1)\beta$ en $F(n-1)\beta$ et $F(n\beta)$. Donc en faisant successivement $n=1$, 2, 3\ldots, on trouvera successivement les valeurs des fonctions:
\[ \begin{gathered} \Phi(2\beta);\ \Phi(3\beta);\ \Phi(4\beta)\ .....\ \Phi(n\beta),\\ f(2\beta);\ f(3\beta);\ f(4\beta)\ .....\ f(n\beta),\\ F(2\beta);\ F(3\beta);\ F(4\beta)\ .....\ F(n\beta) \end{gathered} \]
exprimées en fonctions rationelles des trois quantités
\[ \Phi\beta;\ f\beta;\ F\beta. \]
En faisant p.\ ex.\ $n=1$, on aura:
\[ \left\{\begin{aligned} \Phi(2\beta)&=\frac{2.\varphi\beta.f\beta.F\beta}{1+e^{2}c^{2}\varphi^{4}\beta},\\ f(2\beta)&=-1+\frac{2f^{2}\beta}{1+e^{2}c^{2}\varphi^{4}\beta},\\ F(2\beta)&=-1+\frac{2F^{2}\beta}{1+e^{2}c^{2}\varphi^{4}\beta}. \end{aligned}\right. \tag{35.} \]

Les fonctions $\Phi(n\beta)$, $f(n\beta)$, $F(n\beta)$ étant des fonctions rationelles de $\Phi\beta$, $f\beta$, $F\beta$, on peut toujours les réduire à la forme $\frac{P}{Q}$, où $P$ et $Q$ sont des fonctions entières de $\Phi\beta$, $f\beta$, $F\beta$. De même il est clair, que le dénominateur $Q$ aura la même valeur pour les trois fonctions, que l'on considère. Soit donc
\[ \Phi(n\beta)=\frac{P_n}{Q_n},\ f(n\beta)=\frac{P'_n}{Q_n},\ F(n\beta)=\frac{P''_n}{Q_n}, \tag{35$'$.} \]
on aura également
\[ \begin{gathered} \Phi(n+1)\beta=\frac{P_{n+1}}{Q_{n+1}},\ f(n+1)\beta=\frac{P'_{n+1}}{Q_{n+1}},\ F(n+1)\beta=\frac{P''_{n+1}}{Q_{n+1}},\\ \Phi(n-1)\beta=\frac{P_{n-1}}{Q_{n-1}},\ f(n-1)\beta=\frac{P'_{n-1}}{Q_{n-1}},\ F(n-1)\beta=\frac{P''_{n-1}}{Q_{n-1}}. \end{gathered} \]
En substituant ces valeurs, la première des formules (34.) deviendra:
\[ \frac{P_{n+1}}{Q_{n+1}}=-\frac{P_{n-1}}{Q_{n-1}}+\frac{2f\beta.F\beta.\dfrac{P_n}{Q_n}}{1+c^{2}e^{2}\varphi^{2}\beta.\dfrac{P_n^{2}}{Q_n^{2}}}, \]

\origpage{116}
ou bien
\[ \frac{P_{n+1}}{Q_{n+1}}=\frac{-P_{n-1}.(Q_n^{2}+c^{2}e^{2}\varphi^{2}\beta.P_n^{2})+2P_n.Q_n.Q_{n-1}.f\beta.F\beta}{Q_{n-1}.(Q_n^{2}+e^{2}c^{2}\varphi^{2}\beta.P_n^{2})}. \]

En égalant les numérateurs et les dénominateurs de ces deux fractions, on aura
\[ P_{n+1}=-P_{n-1}(Q_n^{2}+c^{2}e^{2}\Phi^{2}\beta.P_n^{2})+2f\beta.F\beta.P_n.Q_n.Q_{n-1}, \tag{36.} \]
\[ Q_{n+1}=Q_{n-1}(Q_n^{2}+e^{2}c^{2}\Phi^{2}\beta.P_n^{2}). \tag{37.} \]
La seconde et la troisième des équations (34.) donneront de la même manière:
\[ P'_{n+1}=-P'_{n-1}.(Q_n^{2}+c^{2}e^{2}\Phi^{2}\beta.P_n^{2})+2f\beta\,P'_n.Q_n.Q_{n-1}, \tag{38.} \]
\[ P''_{n+1}=-P''_{n-1}.(Q_n^{2}+c^{2}e^{2}\Phi^{2}\beta.P_n^{2})+2F\beta.P''_n.Q_n.Q_{n-1}. \tag{39.} \]
En faisant dans ces quatre formules $n=1$, 2, 3\ ....., et remarquant, qu'on aura:
\[ \begin{gathered} Q_0=1;\ Q_1=1;\ P_0=0;\ P_1=\Phi\beta;\\ P'_0=1;\ P'_1=f\beta;\ P''_0=1;\ P''_1=F\beta; \end{gathered} \]
on trouvera successivement les fonctions entières $Q_n$, $P_n$, $P'_n$, $P''_n$, pour toutes les valeurs de $n$.

Soient pour abréger:
\[ \Phi\beta=x,\ f\beta=y,\ F\beta=z\ \text{et} \tag{40.} \]
\[ R_n=Q_n^{2}+e^{2}c^{2}x^{2}P_n^{2}, \tag{41.} \]
les formules précédentes donneront:
\[ \left\{\begin{aligned} Q_{n+1}&=Q_{n-1}.R_n,\\ P_{n+1}&=-P_{n-1}.R_n+2yz.P_n.Q_n.Q_{n-1},\\ P'_{n+1}&=-P'_{n-1}.R_n+2y.P'_n.Q_n.Q_{n-1},\\ P''_{n+1}&=-P''_{n-1}.R_n+2z.P''_n.Q_n.Q_{n-1}. \end{aligned}\right. \tag{42.} \]
En posant $n=1$, 2, on aura:
\[ \left\{\begin{aligned} R_1&=Q_1^{2}+e^{2}c^{2}x^{2}P_1^{2}=1+e^{2}c^{2}x^{4},\\ Q_2&=Q_0R_1=1+e^{2}c^{2}x^{4},\\ P_2&=-P_0R_1+2yzP_1.Q_1.Q_0=2xyz,\\ P'_2&=-P'_0R_1+2yP'_1.Q_1.Q_0=-1-e^{2}c^{2}x^{4}+2y^{2},\\ P''_2&=-P''_0R_1+2zP''_1.Q_1.Q_0=-1-e^{2}c^{2}x^{4}+2z^{2}. \end{aligned}\right. \tag{43.} \]
\[ \left\{\begin{aligned} R_2&=Q_2^{2}+e^{2}c^{2}x^{2}.P_2^{2}=(1+e^{2}c^{2}x^{4})^{2}+e^{2}c^{2}x^{2}.4x^{2}y^{2}z^{2},\\ Q_3&=Q_1R_2=R_2,\\ P_3&=-P_1R_2+2yzP_2Q_2Q_1=-xR_2+2y^{2}z^{2}x.Q_2\\ &=x\{2y^{2}z^{2}.Q_2-R_2\},\\ P'_3&=-P'_1.R_2+2y.P'_2.Q_2Q_1=-y.R_2+2yP'_2Q_2\\ &=y\{2Q_2P'_2-R_2\},\\ P''_3&=z\{2Q_2P''_2-R_2\}. \end{aligned}\right. \tag{44.} \]

\origpage{117}
En continuant de cette sorte, et en remarquant que $y^{2}=1-c^{2}x^{2}$, $z^{2}=1+e^{2}x^{2}$, on verra aisément, que les quantités:
\[ Q_n,\ \frac{P_{2n}}{xyz},\ \frac{P_{2n+1}}{x},\ P'_{2n},\ \frac{P'_{2n+1}}{y},\ P''_{2n},\ \frac{P''_{2n+1}}{z} \]
sont des fonctions entières des trois quantités $x^{2}$, $y^{2}$, $z^{2}$ et par conséquent aussi de l'une de ces quantités quelconque pour une valeur entière quelconque de $n$.

Cela fait voir que les expressions de $\Phi(n\beta)$, $f(n\beta)$, $F(n\beta)$ seront de la forme suivante:
\[ \left\{\begin{aligned} \Phi(2n\beta)&=\Phi\beta.f\beta.F\beta.T, & \Phi(2n+1)\beta&=\Phi\beta.T',\\ f(2n\beta)&=T_1, & f(2n+1)\beta&=f\beta.T'',\\ F(2n\beta)&=T_2, & F(2n+1)\beta&=F\beta.T''', \end{aligned}\right. \tag{45.} \]
où $T$ etc.\ représentent des fonctions rationelles des quantités $(\Phi\beta)^{2}$, $(f\beta)^{2}$, $(F\beta)^{2}$.

\section*{§. III.}

\emph{Résolution des équations}
\[ \Phi(n\beta)=\frac{P_n}{Q_n},\ f(n\beta)=\frac{P'_n}{Q_n},\ F(n\beta)=\frac{P''_n}{Q_n}. \]

\subsection*{10.}

Suivant ce qu'on a vu, les fonctions $\Phi(n\beta)$, $f(n\beta)$, $F(n\beta)$ s'expriment rationellement en $x$, $y$, $z$. Le cas contraire n'a pas lieu, car les équations (35.) sont en général d'un degré très élevé. Elles ont par cette raison un certain nombre de racines. Nous allons voir, comment on peut aisément exprimer toutes ces racines au moyen des fonctions $\Phi$, $f$, $F$.

A. Considérons d'abord l'équation $\Phi(n\beta)=\frac{P_n}{Q_n}$, ou $Q_n.\Phi(n\beta)=P_n$, et cherchons toutes les valeurs de $x$.

Il faut distinguer deux cas, selon que $n$ est pair ou impair:

1) \emph{Si $n$ est un nombre pair.}

D'après ce qu'on a vu dans le paragraphe précédent (45.), on aura dans ce cas
\[ \Phi(2n\beta)=\psi(x^{2}).x.y.z, \]
c'est-à-dire, en vertu des formules
\[ y=\sqrt{(1-c^{2}x^{2})},\ z=\sqrt{(1+e^{2}x^{2})}, \]
\[ \Phi(2n\beta)=x\psi(x^{2}).\sqrt{[(1-c^{2}x^{2})(1+e^{2}x^{2})]}. \]
Donc l'équation en $x$ deviendra,
\[ \Phi^{2}(2n\beta)=x^{2}(\psi(x^{2}))^{2}.(1-c^{2}x^{2})(1+e^{2}x^{2}). \]

\origpage{118}
En désignant le second nombre par $\theta(x^{2})$, on aura
\[ \Phi^{2}(2n\beta)=\theta(x^{2}). \]
$\Phi\beta$ étant une des valeurs de $x$, on aura
\[ \Phi^{2}(2n\beta)=\theta(\Phi^{2}\beta), \tag{46.} \]
équation, qui a lieu, quelle que soit la valeur de $\beta$. Pour trouver les autres valeurs de $x$, soit $x=\Phi\alpha$ une racine quelconque, on doit avoir
\[ \Phi^{2}(2n\beta)=\theta(\Phi^{2}\alpha). \]
Or, en mettant dans (46.) $\alpha$ au lieu de $\beta$, il viendra
\[ \Phi^{2}(2n\alpha)=\theta(\Phi^{2}\alpha),\ \text{donc:} \]
\[ \Phi^{2}(2n\beta)=\Phi^{2}(2n\alpha); \tag{47.} \]
équation, qui revient à ces deux-ci:
\[ \Phi(2n\alpha)=\Phi(2n\beta)\ \text{et}\ \Phi(2n\alpha)=-\Phi(2n\beta). \]
La première donne en vertu de (31.)
\[ 2n\alpha=2n\beta.(-1)^{m+\mu}+m\omega+\mu\varpi i, \]
où $m$ et $\mu$ sont deux nombres entiers quelconques, positifs ou négatifs, zéro y compris.

La seconde donne les mêmes valeurs de $2m\alpha$\ednote{Printed $2m\alpha$ instead of $2n\alpha$; an apparent misprint.}, mais de signe contraire, comme il est aisé de voir, en l'écrivant comme suit:
\[ \Phi(-2n\alpha)=\Phi(2n\beta). \]
Toute valeur de $2n\alpha$, qui satisfait à l'équation (47.), peut donc être représentée par
\[ 2n\alpha=\pm(2n\beta.(-1)^{m+\mu}+m\omega+\mu\varpi i). \]
De là on tire la valeur de $\alpha$, en divisant par $2n$, savoir:
\[ \alpha=\pm\left((-1)^{m+\mu}.\beta+\frac{m}{2n}\omega+\frac{\mu}{2n}\varpi i\right). \]
Ayant la valeur de $\alpha$, on aura
\[ \Phi\alpha=\pm\Phi\left((-1)^{m+\mu}.\beta+\frac{m}{2n}\omega+\frac{\mu}{2n}\varpi i\right)=x. \tag{48.} \]
Donc toutes les valeurs de $x$ sont contenues dans cette expression, et on les trouvera, en donnant aux nombres $m$ et $\mu$ toutes les valeurs entières depuis $-\infty$ jusqu'à $+\infty$. Or pour avoir toutes celles qui sont différentes entre elles, il suffit de donner à $m$ et $\mu$ des valeurs entières moindres que $2n$. En effet, quels que soient ces nombres, on peut toujours les supposer réduits a la forme:
\[ m=2n.k+m',\ \mu=2n.k'+\mu', \]
où $k$, $k'$ sont des nombres entiers, et $m'$, $\mu'$, des nombres entiers moindres que $2n$. En substituant ces valeurs dans l'expression de $x$, elle deviendra:

\origpage{119}
\[ x=\pm\Phi\left\{(-1)^{m'+\mu'}\beta+\frac{m'}{2n}\omega+\frac{\mu'}{2n}\varpi i+k\omega+k'\varpi i\right\}, \]
or en vertu de (22.) cette expression se réduit à
\[ x=\pm\Phi\left\{(-1)^{m'+\mu'}\beta+\frac{m'}{2n}\omega+\frac{\mu'}{2n}\varpi i\right\}. \tag{49.} \]
Cette valeur de $x$ est de la même forme que la précédente (48.), $m$ et $\mu$ seulement sont remplacés par $m'$ et $\mu'$, qui, tous les deux, sont positifs et moindres que $2n$; donc on obtiendra toutes les valeurs différentes de $x$, en donnant seulement à $m$ et $\mu$ toutes les valeurs entières depuis zéro jusqu'à $2n$ excl. Toutes ces valeurs sont nécessairement différentes entre elles. En effet, supposons par ex.\ qu'on ait
\[ \begin{gathered} \pm\Phi\left\{(-1)^{m'+\mu'}\beta+\frac{m'}{2n}\omega+\frac{\mu'}{2n}\varpi i\right\}\\ =\pm\Phi\left\{(-1)^{m+\mu}\beta+\frac{m}{2n}\omega+\frac{\mu}{2n}\varpi i\right\}, \end{gathered} \]
il s'en suivrait, d'après (31.):
\[ (-1)^{m'+\mu'}\beta+\frac{m'}{2n}\omega+\frac{\mu'}{2n}\varpi i=\pm\left\{(-1)^{m+\mu}.\beta+\frac{m}{2n}\omega+\frac{\mu}{2n}\varpi i\right\}+k\omega+k'\varpi i, \]
$k$ et $k'$ étant des entiers.

Cette équation donne:
\[ \mu'=k'.2n\pm\mu,\ m'=k.2n\pm m,\ (-1)^{m'+\mu'}=\pm(-1)^{m+\mu}. \]
Les deux premières équations ne peuvent pas subsister à moins que $k'=1$, $k=1$, $\mu'=2n-\mu$, $m'=2n-m$, et alors la dernière deviendra:
\[ (-1)^{4n-m-\mu}=-(-1)^{m+\mu}, \]
d'où l'on tire:
\[ (-1)^{2m+2\mu}=-1, \]
résultat absurde.

Donc toutes les valeurs de $x$, contenues dans la formule (48.), sont différentes entre elles, si $m$ et $\mu$ sont positifs et moindres que $m$\ednote{Printed $m$ instead of $2n$; an apparent misprint.}.

Le nombre total des valeurs de $x$ est, comme il est aisé de voir, égal à $2(2n)^{2}=8n^{2}$, or l'équation $\Phi^{2}(2n\beta)=\theta(x^{2})$ ne peut pas avoir des racines égales, car dans ce cas on auroit $\frac{\partial.\theta(x^{2})}{\partial x}=0$, ce qui donneroit pour $x$ une valeur indépendante de $\beta$. Donc le degré de l'équation $\Phi^{2}(2n\beta)=\theta(x^{2})$ est égal au nombre des racines, c'est-à-dire à $8n^{2}$. Si par ex.\ $n=1$, on aura l'équation
\[ \Phi^{2}(2\beta)=\theta(x^{2})=\frac{4x^{2}(1-c^{2}x^{2})(1+e^{2}x^{2})}{(1+e^{2}c^{2}x^{4})^{2}}, \]
ou bien
\[ (1+e^{2}c^{2}x^{4})^{2}.\Phi^{2}(2\beta)=4x^{2}(1-c^{2}x^{2})(1+e^{2}x^{2}), \]

\origpage{120}
et d'après la formule (48.) les racines de cette équation, au nombre de huit, seront:
\[ \begin{gathered} x=\pm\Phi\beta,\ x=\pm\Phi\left(-\beta+\frac{\omega}{2}\right),\\ x=\pm\Phi\left(-\beta+\frac{\varpi}{2}i\right),\ x=\pm\Phi\left(\beta+\frac{\omega}{2}+\frac{\varpi}{2}i\right). \end{gathered} \]

2) \emph{Si $n$ est un nombre impair} $=2n+1$.

Dans ce cas $\frac{P_{2n+1}}{Q_{2n+1}}$ est, comme nous l'avons vu, une fonction rationelle de $x$, et par conséquent l'équation en $x$ sera:
\[ \Phi(2n+1)\beta=\frac{P_{2n+1}}{Q_{2n+1}}. \tag{50.} \]
Précisément comme dans le cas précédent, on trouvera, que toutes les racines de cette équation peuvent être représentées par
\[ x=\Phi\left((-1)^{m+\mu}.\beta+\frac{m}{2n+1}\omega+\frac{\mu}{2n+1}\varpi i\right), \tag{51.} \]
où il faut donner à $m$ et $\mu$ toutes les valeurs entières depuis $-n$ jusqu'à $+n$ incl. Donc le nombre des racines différentes est $(2n+1)^{2}$. C'est aussi le degré de l'équation en question. On peut aussi exprimer les racines par
\[ x=(-1)^{m+\mu}.\Phi\left(\beta+\frac{m}{2n+1}\omega+\frac{\mu}{2n+1}\varpi i\right). \]
Si par ex.\ $n=1$, on aura une équation du degré $3^{2}=9$.

La formule (51.) donne pour $x$ les 9 valeurs suivantes:
\[ \begin{gathered} \Phi(\beta);\\ \Phi\left(-\beta-\frac{\omega}{3}\right);\\ \Phi\left(-\beta+\frac{\omega}{3}\right);\\ \Phi\left(-\beta-\frac{\varpi}{3}i\right);\\ \Phi\left(-\beta+\frac{\varpi}{3}i\right);\\ \Phi\left(\beta-\frac{\omega}{3}-\frac{\varpi}{3}i\right);\\ \Phi\left(\beta-\frac{\omega}{3}+\frac{\varpi}{3}i\right);\\ \Phi\left(\beta+\frac{\omega}{3}-\frac{\varpi}{3}i\right);\\ \Phi\left(\beta+\frac{\omega}{3}+\frac{\varpi}{3}i\right); \end{gathered} \]

\origpage{121}

B. Considérons maintenant l'équation
\[ f(n\beta)=\frac{P'_n}{Q_n} \tag{52.} \]
et cherchons les valeurs de $y$, qui satisfont à cette équation. La fonction $\frac{P'_n}{Q_n}$ étant, comme on a vu plus haut, rationelle en $y$: l'équation en $y$, en faisant $\frac{P'_n}{Q_n}=\psi(y)$, sera
\[ f(n\beta)=\psi(y). \]
Une des racines de cette équation est $y=f\beta$, donc, quelle que soit la leur\ednote{Printed ``leur'' instead of ``valeur''; the first syllable of the word is missing at the line break.} de $\beta$:
\[ f(n\beta)=\psi(f\beta). \tag{53.} \]
Pour trouver les autres valeurs de $y$, soit $\alpha$ une nouvelle inconnue, telle que $y=f\alpha$, on aura
\[ f(n\beta)=\psi(f\alpha); \]
or, en vertu de (53.), le second membre est égal à $f(n\alpha)$, donc pour déterminer $\alpha$, on aura l'équation:
\[ f(n\alpha)=f(n\beta). \]
En vertu de (32.) cette équation donne pour expression générale de $n\alpha$:
\[ n\alpha=\pm n\beta+2m\omega+\mu\varpi i, \]
$m$ et $\mu$ étant deux nombres entiers positifs ou négatifs, zéro y compris.

De là on tire
\[ \alpha=\pm\beta+\frac{2m}{n}\omega+\frac{\mu}{n}\varpi i, \]
et par conséquent:
\[ f\alpha=f\left(\pm\beta+\frac{2m}{n}\omega+\frac{\mu}{n}\varpi i\right)=y. \]
C'est la valeur générale de $y$. Maintenant pour avoir les valeurs différentes de $y$, je dis, qu'il suffit de prendre $\beta$ avec le signe $+$ et de donner à $m$ et $\mu$ toutes les valeurs entières, moindres que $n$. En effet, comme on a $f(+\alpha)=f(-\alpha)$, on aura d'abord:
\[ f\left(-\beta+\frac{2m}{n}\omega+\frac{\mu}{n}\varpi i\right)=f\left(\beta-\frac{2m}{n}\omega-\frac{\mu}{n}\varpi i\right). \]
Donc on peut toujours dans l'expression de $y$ prendre $\beta$ avec le signe $+$. Ainsi toutes les valeurs de $y$ sont contenues dans l'expression
\[ y=f\left(\beta+\frac{2m}{n}\omega+\frac{\mu}{n}\varpi i\right). \tag{54.} \]

Maintenant quels que soient les nombres $m$ et $\mu$, on peut toujours supposer,
\[ m=k.n+m',\ \mu'=k'.n+\mu', \]
\ednote{Printed $\mu'$ on the left of the second equation instead of $\mu$; an apparent misprint.}

\origpage{122}
où $k$, $k'$, $m'$, $\mu'$ sont des nombres entiers, les deux derniers étant en même temps positifs et moindres que $n$.

En substituant, il viendra
\[ y=f\left(\beta+\frac{2m'}{n}\omega+\frac{\mu'}{n}\varpi i+2k\omega+k'\varpi i\right). \]
Or, en vertu de (22.) le second membre de cette équation est égal à
\[ f\left(\beta+\frac{2m'}{n}\omega+\frac{\mu'}{n}\varpi i\right)=y, \tag{55.} \]
quantité de la même forme, que le second membre de (54.); seulement $m'$ et $\mu'$ sont positifs et moindres que $n$. Donc etc.

En donnant à $m$ et $\mu$ toutes les valeurs possibles, moindres que $n$, on trouvera un nombre $n^{2}$ de valeurs de $y$. Or, en général toutes ces quantités sont différentes entre elles. En effet, supposons par ex.
\[ f\left(\beta+\frac{2m}{n}\omega+\frac{\mu}{n}\varpi i\right)=f\left(\beta+\frac{2\mu'}{n}+\frac{\mu'}{n}\varpi i\right), \]
\ednote{In the right member the print has $2\mu'/n$ where $2m'\omega/n$ is expected; the $\omega$ is absent. An apparent misprint.}
on aura en vertu de (32.), en désignant par $k$, $k'$ deux nombres entiers:
\[ \beta+\frac{2m}{n}\omega+\frac{\mu}{n}\varpi i=\pm\left(\beta+\frac{2m'}{n}\omega+\frac{\mu'}{n}\varpi i\right)+2k\omega+k'\varpi i. \]
Puisque $\beta$ peut avoir une valeur irrationelle quelconque, il est clair que cette équation ne peut pas subsister à moins qu'on ne préfère dans le second membre le signe supérieur. Alors il viendra
\[ \frac{2m}{n}\omega+\frac{\mu}{n}\varpi i=\frac{2m'}{n}\omega+\frac{\mu'}{n}\varpi i+2k\omega+k'\varpi i, \]
d'où l'on tire, en égalant les parties réelles et les parties imaginaires:
\[ m=m'+kn;\ \mu=\mu'+k'n, \]
équations absurdes, en remarquant que les nombres $m$, $m'$, $\mu$ et $\mu'$ sont tous positifs et inférieurs à $n$. Donc en général l'équation
\[ f(n\beta)=\psi(y) \]
a un nombre $n^{2}$ de racines différentes entre elles et non un plus grand nombre. Or généralement toutes les racines de cette équation sont différentes entres elles. En effet, si deux d'entre elles étoient égales, on auroit à la fois:
\[ f(n\beta)=\psi(y)\ \text{et}\ 0=\psi'(y), \]
et cela est impossible, si l'on remarque que les coéfficiens de $y$ dans $\psi(y)$ ne contiennent pas $\beta$. Donc généralement l'équation (52.) est nécessairement du degré $n^{2}$.

\origpage{123}

C. L'équation
\[ F(n\beta)=\frac{P''_n}{Q_n}, \tag{56.} \]
étant traitée absolument de la même manière par rapport à $z$, que l'équation $f(n\beta)=\frac{P'_n}{Q_n}$ l'a été par rapport à $y$, donne pour expression générale des valeurs de $z$:
\[ z=F\left(\beta+\frac{m}{n}\omega+\frac{2\mu}{n}\varpi i\right), \tag{57.} \]
où $m$ et $\mu$ sont entiers, positifs et moindres que $n$. Le nombre des valeurs de $z$ est $n^{2}$, et elles sont en général toutes différentes entre elles.

Donc généralement l'équation (56.) est du degré $n^{2}$.

\subsection*{11.}

Nous avons trouvé ci-dessus toutes les racines des équations
\[ \varphi(n\beta)=\frac{P_n}{Q_n},\ f(n\beta)=\frac{P'_n}{Q_n},\ F(n\beta)=\frac{P''_n}{Q_n}, \]
racines, qui sont exprimées par les formules (48.), (51.), (54.), (57.). Toutes ces racines sont différentes entre elles, excepté les cas de valeurs particulières de $\beta$; mais pour ces valeurs, les racines différentes sont contenues dans les mêmes formules. --- Dans ce dernier cas un certain nombre des valeurs des quantités $x$, $y$, $z$ seront égales; mais il est clair, que toutes les valeurs égales ou inégales seront néanmoins les racines des équations, dont il s'agit. Cela se fait voir en faisant converger $\beta$ vers une valeur particulière, qui donne pour $x$, ou $y$, ou $z$ des valeurs égales.

En faisant dans la formule (48.) $\beta=\frac{\alpha}{2n}$, on aura l'équation
\[ \varphi^{2}\alpha=\frac{P_{2n}^{2}}{Q_{2n}^{2}},\ \text{dont les racines sont}\ x=\pm\varphi\left((-1)^{m+\mu}\frac{\alpha}{2n}+\frac{m}{2n}\omega+\frac{\mu}{2n}\varpi i\right); \tag{58.} \]
et où $m$ et $\mu$ ont toutes les valeurs entières et positives moindres que $2n$.

En faisant de même dans la formule (50.) $\beta=\frac{\alpha}{2n+1}$, on aura
\[ \Phi\alpha=\frac{P_{2n+1}}{Q_{2n+1}},\ \text{dont les racines sont} \]
\[ x=(-1)^{m+\mu}.\Phi\left(\frac{\alpha}{2n+1}+\frac{m\omega+\mu\varpi i}{2n+1}\right), \tag{59.} \]
$m$ et $\mu$ ayant pour valeurs tous les nombres entiers depuis $-n$ jusqu'à $+n$.

Enfin en faisant dans (52.), (56.) $\beta=\frac{\alpha}{n}$, on aura l'équation
\[ f\alpha=\frac{P'_n}{Q_n},\ \text{dont les racines sont} \]
\[ y=f\left(\frac{\alpha}{n}+\frac{2m}{n}\omega+\frac{\mu}{n}\varpi i\right), \tag{60.} \]

\origpage{124}

et l'équation
\[ F\alpha=\frac{P''_n}{Q_n},\ \text{dont les racines sont} \]
\[ z=F\left(\frac{\alpha}{n}+\frac{m}{n}\omega+\frac{2\mu}{n}\varpi i\right), \tag{61.} \]
où $m$ et $\mu$ sont renfermés entre les limites 0 et $n-1$ incl. Si $n$ est impair $=2n+1$, on peut aussi supposer
\[ \begin{gathered} y=(-1)^{m}.f\left(\frac{\alpha}{2n+1}+\frac{m}{2n+1}\omega+\frac{\mu}{2n+1}\varpi i\right),\\ z=(-1)^{\mu}.F\left(\frac{\alpha}{2n+1}+\frac{m}{2n+1}\omega+\frac{\mu}{2n+1}\varpi i\right), \end{gathered} \]
$m$ et $\mu$ ayant toutes les valeurs entières de $-n$ à $+n$.

Dans toutes ces équations la quantité $\alpha$ peut avoir une valeur quelconque.

Comme cas particuliers on doit remarquer les suivans:

1) En faisant dans (58.) et (59.) $\alpha=0$, on aura les équations
\[ \left\{\begin{aligned} &P_{2n}^{2}=0,\ \text{dont les racines sont}\ x=\pm\Phi\left(\frac{m}{2n}\omega+\frac{\mu}{2n}\varpi i\right)\\ &\text{(les limites de $m$ et $\mu$ étant 0 et $2n-1$),}\\ &P_{2n+1}=0,\ \text{dont les racines sont}\ x=\Phi\left(\frac{m}{2n+1}\omega+\frac{\mu}{2n+1}\varpi i\right)\\ &\text{(les limites de $m$ et $\mu$ étant $-n$ et $+n$).} \end{aligned}\right. \tag{62.} \]

2) En faisant dans (60.) $\alpha=\frac{\omega}{2}$ et dans (61.) $\alpha=\frac{\varpi}{2}i$, et remarquant, que $f\left(\frac{\omega}{2}\right)=0$, $F\left(\frac{\varpi}{2}i\right)=0$, on obtiendra les deux équations:
\[ P'_n=0,\ \text{dont les racines sont}\ y=f\left((2m+\tfrac{1}{2})\frac{\omega}{n}+\frac{\mu}{n}\varpi i\right) \tag{63.} \]
\[ P''_n=0,\ \text{dont les racines sont}\ z=F\left(\frac{m}{n}\omega+(2\mu+\tfrac{1}{2})\frac{\varpi i}{n}\right) \tag{64.} \]
(les limites de $m$ et $\mu$ étant 0 et $n-1$).

3) En faisant dans (58.) $\alpha=\frac{\omega}{2}+\frac{\varpi}{2}i$, et en remarquant, que $\varphi\left(\frac{\omega}{2}+\frac{\varpi}{2}i\right)=\frac{1}{0}$, on aura l'équation
\[ Q_{2n}^{2}=0, \]
dont les racines seront:
\[ x=\pm\Phi\left((m+\tfrac{1}{2}(-1)^{m+\mu})\frac{\omega}{2n}+(\mu+\tfrac{1}{2}(-1)^{m+\mu}).\frac{\varpi i}{2n}\right). \]

Les valeurs de $x$, doivent être égales par paires, et l'on verra aisement, que les valeurs inégales peuvent être représentées par
\[ x=\Phi\left((m+\tfrac{1}{2})\frac{\omega}{2n}+(\mu+\tfrac{1}{2})\frac{\varpi i}{2n}\right), \tag{65.} \]

\origpage{125}
en donnant à $m$ et $\mu$ toutes les valeurs entières depuis 0 à $2n-1$. Donc ce sont les racines de l'équation
\[ Q_{2n}=0,\ \text{par rapport à } x. \]

En faisant de même dans (59.) $\alpha=\frac{\omega}{2}+\frac{\varpi}{2}i$, on aura l'équation
\[ Q_{2n+1}=0, \]
dont les racines seront:
\[ \left\{\begin{aligned}
x&=(-1)^{m+\mu}.\varphi\left((m+\tfrac{1}{2})\frac{\omega}{2n+1}+(\mu+\tfrac{1}{2})\frac{\varpi i}{2n+1}\right),\\
y&=(-1)^{m}.f\left((m+\tfrac{1}{2})\frac{\omega}{2n+1}+(\mu+\tfrac{1}{2})\frac{\varpi i}{2n+1}\right),\\
z&=(-1)^{\mu}.F\left((m+\tfrac{1}{2})\frac{\omega}{2n+1}+(\mu+\tfrac{1}{2})\frac{\varpi i}{2n+1}\right),
\end{aligned}\right. \tag{66.} \]
$m$ et $\mu$ ayant pour valeurs tous les nombres entiers de $-n$ à $+n$.

Parmi les valeurs de $x$, $y$, $z$, il faut remarquer celle, qui répond à $m=n$, $\mu=n$. Alors on a
\[ \begin{aligned}
x&=(-1)^{n}.\varphi\left(\frac{\omega}{2}+\frac{\varpi}{2}i\right)=\frac{1}{0},\\
y&=(-1)^{n}.f\left(\frac{\omega}{2}+\frac{\varpi}{2}i\right)=\frac{1}{0},\\
z&=(-1)^{n}.F\left(\frac{\omega}{2}+\frac{\varpi}{2}i\right)=\frac{1}{0}.
\end{aligned} \]

Ces valeurs infinies font voir que l'équation $Q_{2n+1}=0$ est d'un degré moindre d'une unité que celui des équations, dont elle sort. En écartant ces valeurs, les restantes, en nombre de $(2n+1)^{2}-1$, seront les racines de l'équation $Q_{2n+1}=0$.

\section*{§. IV.}

\emph{Résolution algébrique des équations}
\[ \Phi\alpha=\frac{P_{2n+1}}{Q_{2n+1}},\quad f\alpha=\frac{P'_{2n+1}}{Q_{2n+1}},\quad F\alpha=\frac{P''_{2n+1}}{Q_{2n+1}}. \]

\subsection*{12.}

Nous avons vu dans le §. précédent, comment on peut exprimer aisément les racines des équations en question au moyen des fonctions $\Phi$, $f$, $F$. Nous allons maintenant en déduire la résolution de ces mêmes équations, ou la determination des fonctions $\varphi\left(\frac{\alpha}{n}\right)$, $f\left(\frac{\alpha}{n}\right)$, $F\left(\frac{\alpha}{n}\right)$, en fonctions de $\varphi\alpha$, $f\alpha$, $F\alpha$.

\origpage{126}

Comme on a
\[ \varphi\left(\frac{\alpha}{m\mu}\right)=\varphi\left(\frac{1}{m}\left(\frac{\alpha}{\mu}\right)\right), \]
on peut supposer que $n$ est un nombre premier. D'abord nous considérerons le cas, où $n=2$, et ensuite celui, où $n$ est un nombre impair.

A. \emph{Expressions des fonctions} $\varphi\left(\frac{\alpha}{2}\right)$, $f\left(\frac{\alpha}{2}\right)$, $F\left(\frac{\alpha}{2}\right)$.

\subsection*{13.}

Les valeurs de $\varphi\left(\frac{\alpha}{2}\right)$, $f\left(\frac{\alpha}{2}\right)$, $F\left(\frac{\alpha}{2}\right)$ peuvent être trouvées très facilement de la manière suivante. En supposant dans les formules (35.) $\beta=\frac{\alpha}{2}$, et faisant
\[ x=\varphi\left(\frac{\alpha}{2}\right),\ y=f\left(\frac{\alpha}{2}\right),\ z=F\left(\frac{\alpha}{2}\right), \]
il viendra:
\[ f(\alpha)=\frac{y^{2}-c^{2}x^{2}z^{2}}{1+e^{2}c^{2}x^{4}},\ F(\alpha)=\frac{z^{2}+e^{2}y^{2}x^{2}}{1+e^{2}c^{2}x^{4}}, \]
ou bien, en substituant les valeurs de $y^{2}$ et $z^{2}$ en $x^{2}$:
\[ f(\alpha)=\frac{1-2c^{2}x^{2}-c^{2}e^{2}x^{4}}{1+e^{2}c^{2}x^{4}},\ F(\alpha)=\frac{1+2e^{2}x^{2}-e^{2}c^{2}x^{4}}{1+e^{2}c^{2}x^{4}}. \]
Ces équations donnent
\[ \begin{aligned}
1+f\alpha&=\frac{2(1-c^{2}x^{2})}{1+e^{2}c^{2}x^{4}},\ 1-f\alpha=\frac{2c^{2}x^{2}(1+e^{2}x^{2})}{1+e^{2}c^{2}x^{4}},\\
F\alpha-1&=\frac{2e^{2}x^{2}(1-c^{2}x^{2})}{1+e^{2}c^{2}x^{4}},\ F\alpha+1=\frac{2(1+e^{2}x^{2})}{1+e^{2}c^{2}x^{4}},
\end{aligned} \]
\ednote{A blot resembling a 2 is printed after the minus sign in the numerator of the first fraction; the sense requires $2(1-c^{2}x^{2})$, as in the later fractions.}
d'où
\[ \frac{F\alpha-1}{1+f\alpha}=e^{2}.x^{2},\ \frac{1-f\alpha}{F\alpha+1}=c^{2}x^{2}, \]
et de là, en remarquant, que $y^{2}=1-c^{2}x^{2}$, $z^{2}=1+e^{2}x^{2}$:
\[ z^{2}=\frac{F\alpha+f\alpha}{1+f\alpha},\ y^{2}=\frac{F\alpha+f\alpha}{1+F\alpha}. \]

De ces équations on tire, en extrayant la racine carrée, et remplaçant $x$, $y$, $z$ par leurs valeurs $\varphi\left(\frac{\alpha}{2}\right)$, $f\left(\frac{\alpha}{2}\right)$, $F\left(\frac{\alpha}{2}\right)$:
\[ \left\{\begin{aligned}
\varphi\left(\frac{\alpha}{2}\right)&=\frac{1}{c}\sqrt{\left(\frac{1-f\alpha}{1+F\alpha}\right)}=\frac{1}{e}\sqrt{\left(\frac{F\alpha-1}{f\alpha+1}\right)},\\
f\left(\frac{\alpha}{2}\right)&=\sqrt{\left(\frac{F\alpha+f\alpha}{1+F\alpha}\right)},\ F\left(\frac{\alpha}{2}\right)=\sqrt{\left(\frac{F\alpha+f\alpha}{1+f\alpha}\right)}.
\end{aligned}\right. \tag{67.} \]
Telles sont les formes les plus simples qu'on peut donner aux valeurs des fonctions $\varphi\left(\frac{\alpha}{2}\right)$, $f\left(\frac{\alpha}{2}\right)$, $F\left(\frac{\alpha}{2}\right)$. De cette manière on peut exprimer

\origpage{127}
algébriquement $\varphi\left(\frac{\alpha}{2}\right)$, $f\left(\frac{\alpha}{2}\right)$, $F\left(\frac{\alpha}{2}\right)$ en $f\alpha$, $F\alpha$. De la même manière $\varphi\left(\frac{\alpha}{4}\right)$, $f\left(\frac{\alpha}{4}\right)$, $F\left(\frac{\alpha}{4}\right)$ s'exprimeront en $f\left(\frac{\alpha}{2}\right)$, $F\left(\frac{\alpha}{2}\right)$, et ainsi de suite. Donc en général les fonctions $\varphi\left(\frac{\alpha}{2^{n}}\right)$, $f\left(\frac{\alpha}{2^{n}}\right)$, $F\left(\frac{\alpha}{2^{n}}\right)$ peuvent être exprimées au moyen d'extractions de racines carrées en fonctions des trois quantités $\varphi\alpha$, $f\alpha$, $F\alpha$.

Pour appliquer les formules trouvées ci-dessus pour la \emph{bissection} à un exemple, supposons $\alpha=\frac{\omega}{2}$.

Alors on aura $f\left(\frac{\omega}{2}\right)=0$, $F\left(\frac{\omega}{2}\right)=\frac{\sqrt{(e^{2}+c^{2})}}{c}$, donc en substituant:
\[ \varphi\left(\frac{\omega}{4}\right)=\frac{1}{c}\sqrt{\left(\frac{1}{1+\frac{1}{c}\sqrt{(e^{2}+c^{2})}}\right)}=\frac{1}{e}.\sqrt{\left(\frac{1}{c}\sqrt{(e^{2}+c^{2})}-1\right)}, \]
\[ f\left(\frac{\omega}{4}\right)=\sqrt{\left(\frac{\frac{1}{c}\sqrt{(e^{2}+c^{2})}}{1+\frac{1}{c}\sqrt{(e^{2}+c^{2})}}\right)}, \]
\[ F\left(\frac{\omega}{4}\right)=\sqrt{\left(\frac{1}{c}\sqrt{(e^{2}+c^{2})}\right)}, \]
ou bien
\[ \begin{aligned}
\varphi\left(\frac{\omega}{4}\right)&=\frac{1}{\sqrt{(c^{2}+c\sqrt{(e^{2}+c^{2})})}}=\frac{\sqrt{(\sqrt{(e^{2}+c^{2})}-c)}}{e\sqrt{c}},\\
f\left(\frac{\omega}{4}\right)&=\frac{\sqrt[4]{(e^{2}+c^{2})}}{\sqrt{(c+\sqrt{(e^{2}+c^{2})})}}=\frac{1}{e}.\sqrt{(e^{2}+c^{2}-c\sqrt{(e^{2}+c^{2})})}\\
F\left(\frac{\omega}{4}\right)&=\sqrt[4]{\left(1+\frac{e^{2}}{c^{2}}\right)}=\sqrt{\left[F\left(\frac{\omega}{2}\right)\right]}.
\end{aligned} \]

B. \emph{Expressions des fonctions} $\varphi\left(\frac{\alpha}{2n+1}\right)$, $f\left(\frac{\alpha}{2n+1}\right)$, $F\left(\frac{\alpha}{2n+1}\right)$, \emph{en fonctions algébriques des quantités} $\varphi\alpha$, $f\alpha$, $F\alpha$.

\subsection*{14.}

Pour trouver les valeurs de $\varphi\left(\frac{\alpha}{2n+1}\right)$, $f\left(\frac{\alpha}{2n+1}\right)$, $F\left(\frac{\alpha}{2n+1}\right)$ en $\Phi\alpha$, $f\alpha$, $F\alpha$, il faut résoudre les équations
\[ \Phi\alpha=\frac{P_{2n+1}}{Q_{2n+1}},\ f\alpha=\frac{P'_{2n+1}}{Q_{2n+1}},\ F\alpha=\frac{P''_{2n+1}}{Q_{2n+1}}, \]
qui toutes sont du degré $(2n+1)^{2}$. Nous allons voir, qu'il est toujours possible, d'effectuer algébriquement cette résolution.

\origpage{128}

Soient
\[ \varphi_{1}\beta=\overset{+n}{\underset{-n}{\Sigma}}_{m}\varphi\left(\beta+\frac{2m\omega}{2n+1}\right) \tag{68.} \]
et
\[ \psi\beta=\overset{+n}{\underset{-n}{\Sigma}}_{\mu}\theta^{\mu}.\varphi_{1}\left(\beta+\frac{2\mu\varpi i}{2n+1}\right),\ \psi_{1}\beta=\overset{+n}{\underset{-n}{\Sigma}}_{\mu}\theta^{\mu}\varphi_{1}\left(\beta-\frac{2\mu\varpi i}{2n+1}\right), \tag{69.} \]
où $\theta$ est une racine imaginaire quelconque de l'équation $\theta^{2n+1}-1=0$.
Cela posé, je dis que les deux quantités
\[ \psi\beta.\psi_{1}\beta\ \text{et}\ (\psi\beta)^{2n+1}+(\psi_{1}\beta)^{2n+1}, \]
pourront être exprimées rationellement en $\varphi(2n+1)\beta$.

D'abord en écrivant $\varphi_{1}\beta$ comme suit:
\[ \begin{aligned}
\varphi_{1}\beta&=\varphi\beta.\overset{n}{\underset{1}{\Sigma}}.\left\{\varphi\left(\beta+\frac{2m\omega}{2n+1}\right)+\varphi\left(\beta-\frac{2m\omega}{2n+1}\right)\right\}\\
&=\varphi\beta.\overset{n}{\underset{1}{\Sigma}}(-1)^{m}.\frac{2\varphi\beta.f\left(\frac{2m\omega}{2n+1}\right).F\left(\frac{2m\omega}{2n+1}\right)}{1+e^{2}c^{2}\varphi^{2}\left(\frac{2m\omega}{2n+1}\right).\varphi^{2}\beta},
\end{aligned} \]
\ednote{Printed with a point after $\varphi\beta$ before the sum in both lines, where the sense requires a plus sign; the first line shows a blot over the point and a further point after the sum sign.}
on voit, que $\varphi_{1}\beta$ peut s'exprimer rationellement en $\varphi\beta$. Soit donc $\varphi_{1}\beta=\chi(\varphi\beta)$, on a de même:
\[ \begin{aligned}
\varphi_{1}\left(\beta\pm\frac{2\mu\varpi i}{2n+1}\right)&=\chi\left(\varphi\left(\beta\pm\frac{2\mu\varpi i}{2n+1}\right)\right)\\
&=\chi\left(\frac{\varphi\beta.f\left(\frac{2\mu\varpi i}{2n+1}\right).F\left(\frac{2\mu\varpi i}{2n+1}\right)\pm\varphi\left(\frac{2\mu\varpi i}{2n+1}\right).f\beta.F\beta}{1+e^{2}c^{2}.\varphi^{2}\left(\frac{2\mu\varpi i}{2n+1}\right).\varphi^{2}\beta}\right),
\end{aligned} \]
ou bien, en faisant $\Phi\beta=x$:
\[ f\left(\frac{2\mu\varpi i}{2n+1}\right).F\left(\frac{2\mu\varpi i}{2n+1}\right)=a,\ \varphi\left(\frac{2\mu\varpi i}{2n+1}\right)=b, \]
et en substituant pour $f\beta$ et $F\beta$ leurs valeurs $\sqrt{(1-c^{2}x^{2})}$ et $\sqrt{(1+e^{2}x^{2})}$:
\[ \varphi_{1}\left(\beta\pm\frac{2\mu\varpi i}{2n+1}\right)=\chi\left(\frac{ax\pm b.\sqrt{[(1-c^{2}x^{2})(1+e^{2}x^{2})]}}{1+e^{2}c^{2}b^{2}.x^{2}}\right); \]
or, $\chi$ désignant une fonction rationelle, le second membre de cette équation peut se mettre sous la forme
\[ R_{\mu}\pm R'_{\mu}.\sqrt{[(1-c^{2}x^{2})(1+e^{2}x^{2})]}, \]
où $R_{\mu}$ et $R'_{\mu}$ sont des fonctions rationelles de $x$.

Donc on a
\[ \varphi_{1}\left(\beta\pm\frac{2\mu\varpi i}{2n+1}\right)=R_{\mu}\pm R'_{\mu}.\sqrt{[(1-c^{2}x^{2})(1+e^{2}x^{2})]}. \]
En substituant dans les expressions de $\psi\beta$ et $\psi_{1}\beta$, il viendra:

\origpage{129}
\[ \left\{\begin{aligned}
\psi\beta&=\overset{+n}{\underset{-n}{\Sigma}}_{\mu}(+\theta)^{\mu}.R_{\mu}+\sqrt{[(1-c^{2}x^{2})(1+e^{2}x^{2})]}.\overset{+n}{\underset{-n}{\Sigma}}_{\mu}(+\theta)^{\mu}.R'_{\mu},\\
\psi_{1}\beta&=\overset{+n}{\underset{-n}{\Sigma}}_{\mu}(+\theta)^{\mu}.R_{\mu}-\sqrt{[(1-c^{2}x^{2})(1+e^{2}x^{2})]}.\overset{+n}{\underset{-n}{\Sigma}}_{\mu}(+\theta)^{\mu}.R'_{\mu}.
\end{aligned}\right. \tag{70.} \]
\ednote{The plus sign inside the parentheses before $\theta$ is printed overstruck with a short horizontal bar throughout (70.); kept here as a plain plus.}

Maintenant $R_{\mu}$ et $R'_{\mu}$ étant des fonctions rationelles de $x$, les quantités $\overset{+n}{\underset{-n}{\Sigma}}_{\mu}(+\theta)^{\mu}.R_{\mu}$ et $\overset{+n}{\underset{-n}{\Sigma}}_{\mu}(+\theta)^{\mu}.R'_{\mu}$ le sont également. En élevant donc $\psi\beta$ et $\psi_{1}\beta$ à la $(2n+1)^{\text{ième}}$ puissance: les deux quantités $(\psi\beta)^{2n+1}$ et $(\psi_{1}\beta)^{2n+1}$ pourront se mettre sous la forme:
\[ \begin{aligned}
(\psi\beta)^{2n+1}&=t+t'.\sqrt{[(1-c^{2}x^{2})(1+e^{2}x^{2})]},\\
(\psi_{1}\beta)^{2n+1}&=t-t'.\sqrt{[(1-c^{2}x^{2})(1+e^{2}x^{2})]},
\end{aligned} \]
$t$ et $t'$ étant des fonctions rationelles de $x$. En prenant la somme des valeurs de $(\psi\beta)^{2n+1}$ et $(\psi_{1}\beta)^{2n+1}$, on aura
\[ (\psi\beta)^{2n+1}+(\psi_{1}\beta)^{2n+1}=2t. \]

Donc la quantité $(\psi\beta)^{2n+1}+(\psi_{1}\beta)^{2n+1}$ peut être exprimée rationellement en $x$. Il en est de même du produit $\psi\beta.\psi_{1}\beta$, comme on voit par les équations (70.).

Donc on peut faire
\[ \left\{\begin{aligned}
\psi\beta.\psi_{1}\beta&=\lambda(x),\\
(\psi\beta)^{2n+1}+(\psi_{1}\beta)^{2n+1}&=\lambda_{1}(x),
\end{aligned}\right. \tag{71.} \]
$\lambda(x)$ et $\lambda_{1}(x)$ désignant des fonctions rationelles de $x$. Maintenant ces fonctions ont la propriété, de ne pas changer de valeur, lorsqu'on met à la place de $x$ une autre racine quelconque de l'équation
\[ \Phi(2n+1)\beta=\frac{P_{2n+1}}{Q_{2n+1}}. \]

Considérons d'abord la fonction $\lambda(x)$. En remettant la valeur de $x=\Phi\beta$, on aura
\[ \psi\beta.\psi_{1}\beta=\lambda(\Phi\beta), \]
d'où l'on tire, en mettant $\beta+\frac{2k\omega}{2n+1}+\frac{2k'\varpi i}{2n+1}$ au lieu de $\beta$:
\[ \lambda\left(\varphi\left(\beta+\frac{2k\omega}{2n+1}+\frac{2k'\varpi i}{2n+1}\right)\right)=\psi\left(\beta+\frac{2k'\varpi i}{2n+1}+\frac{2k\omega}{2n+1}\right).\psi_{1}\left(\beta+\frac{2k'\varpi i}{2n+1}+\frac{2k\omega}{2n+1}\right). \]
Cela posé, en remarquant, que:
\[ \overset{+n}{\underset{-n}{\Sigma}}\psi(m+k)=\overset{+n}{\underset{-n}{\Sigma}}_{m}\psi(m)+\overset{k}{\underset{1}{\Sigma}}_{m}.(\psi(m+n)-\psi(m-n-1)), \tag{72.} \]
on aura, en faisant dans l'expression de $\psi\beta$, $\beta=\beta+\frac{2k\omega}{2n+1}$:

\origpage{130}
\[ \begin{aligned}
\varphi_{1}\left(\beta+\frac{2k\omega}{2n+1}\right)&=\overset{+n}{\underset{-n}{\Sigma}}_{m}\varphi\left(\beta+\frac{2(k+m)}{2n+1}\omega\right)\\
&=\varphi_{1}\beta+\overset{k}{\underset{1}{\Sigma}}_{m}\left\{\varphi\left(\beta+\frac{2(m+n)\omega}{2n+1}\right)-\varphi\left(\beta+\frac{2(m-n-1)\omega}{2n+1}\right)\right\},
\end{aligned} \]
or: $\varphi\left(\beta+\frac{2(m-n-1)}{2n+1}\omega\right)=\varphi\left(\beta+\frac{2(m+n)}{2n+1}\omega-2\omega\right)=\varphi\left(\beta+\frac{2(m+n)}{2n+1}\omega\right)$,
donc:
\[ \varphi_{1}\left(\beta+\frac{2k\omega}{2n+1}\right)=\varphi_{1}\beta. \tag{73.} \]

En mettant dans l'expression de $\psi\beta$, $\beta+\frac{2k'\varpi i}{2n+1}+\frac{2k\omega}{2n+1}$ au lieu de $\beta$, on trouvera
\[ \psi\left(\beta+\frac{2k'\varpi i}{2n+1}+\frac{2k\omega}{2n+1}\right)=\overset{+n}{\underset{-n}{\Sigma}}_{\mu}\theta^{\mu}.\varphi_{1}\left(\beta+\frac{2(k'+\mu)\varpi i}{2n+1}+\frac{2k\omega}{2n+1}\right), \]
or en vertu de (73.) on a:
\[ \varphi_{1}\left(\beta+\frac{2(k'+\mu)\varpi i}{2n+1}+\frac{2k\omega}{2n+1}\right)=\varphi_{1}\left(\beta+\frac{2(k'+\mu)\varpi i}{2n+1}\right), \]
donc
\[ \psi\left(\beta+\frac{2k'\varpi i}{2n+1}+\frac{2k\omega}{2n+1}\right)=\overset{+n}{\underset{-n}{\Sigma}}_{\mu}\theta^{\mu}.\varphi_{1}\left(\beta+\frac{2(k'+\mu)\varpi i}{2n+1}\right). \]
En vertu de (72.) on a
\[ \begin{aligned}
&\overset{+n}{\underset{-n}{\Sigma}}_{\mu}\theta^{\mu}.\varphi_{1}\left(\beta+\frac{2(k'+\mu)\varpi i}{2n+1}\right)\\
&=\theta^{-k'}.\overset{+n}{\underset{-n}{\Sigma}}_{\mu}\theta^{\mu}.\varphi_{1}\left(\beta+\frac{2\mu\varpi i}{2n+1}\right)+\overset{k'}{\underset{1}{\Sigma}}_{\mu}.\theta^{n+\mu-k'}.\varphi_{1}\left(\beta+\frac{2(\mu+n)\varpi i}{2n+1}\right)\\
&\quad-\overset{k'}{\underset{1}{\Sigma}}_{\mu}\theta^{\mu-n-1-k'}.\varphi_{1}\left(\beta+\frac{2(\mu-n-1)\varpi i}{2n+1}\right),
\end{aligned} \]
donc, en remarquant que $\theta^{n+\mu-k'}=\theta^{\mu-n-1-k'}$ et
\[ \varphi_{1}\left(\beta+\frac{2(\mu-n-1)\varpi i}{2n+1}\right)=\varphi_{1}\left(\beta+\frac{2(\mu+n)\varpi i}{2n+1}-2\varpi i\right)=\varphi_{1}\left(\beta+\frac{2(\mu+n)\varpi i}{2n+1}\right), \]
il viendra
\[ \psi\left(\beta+\frac{2k'\varpi i}{2n+1}+\frac{2k\omega}{2n+1}\right)=\theta^{-k'}.\psi\beta. \tag{74.} \]
De la même manière on trouvera aussi
\[ \psi_{1}\left(\beta+\frac{2k'\varpi i}{2n+1}+\frac{2k\omega}{2n+1}\right)=\theta^{k'}.\psi_{1}\beta. \]
Ces deux équations donneront
\[ \begin{aligned}
&\psi\left(\beta+\frac{2k\omega+2k'\varpi i}{2n+1}\right).\psi_{1}\left(\beta+\frac{2k\omega+2k'\varpi i}{2n+1}\right)=\psi\beta.\psi_{1}\beta,\\
&\left\{\psi\left(\beta+\frac{2k\omega+2k'\varpi i}{2n+1}\right)\right\}^{2n+1}+\left\{\psi_{1}\left(\beta+\frac{2k\omega+2k'\varpi i}{2n+1}\right)\right\}^{2n+1}=(\psi\beta)^{2n+1}+(\psi_{1}\beta)^{2n+1}.
\end{aligned} \]

\origpage{131}
En vertu de ces équations on obtiendra, en mettant dans les valeurs de $\lambda(\Phi\beta)$ et $\lambda_{1}(\Phi\beta)$, $\beta+\frac{2k\omega+2k'\varpi i}{2n+1}$ au lieu de $\beta$:
\[ \begin{aligned}
\lambda(\Phi\beta)&=\lambda\left[\Phi\left(\beta+\frac{2k\omega+2k'\varpi i}{2n+1}\right)\right],\\
\lambda_{1}(\Phi\beta)&=\lambda_{1}\left[\Phi\left(\beta+\frac{2k\omega+2k'\varpi i}{2n+1}\right)\right].
\end{aligned} \]
Or, $\Phi\left(\beta+\frac{2k\omega+2k'\varpi i}{2n+1}\right)$ exprime une racine quelconque de l'équation
\[ \Phi(2n+1)\beta=\frac{P_{2n+1}}{Q_{2n+1}}. \]

Donc comme nous avons dit, les fonctions $\lambda(x)$ et $\lambda_{1}(x)$ auront les mêmes valeurs, quelle que soit la racine, qu'on met à la place de $x$.

Soient donc $x_{0}$, $x_{1}$, $x_{2}....x_{2n+1}$ ces racines, on aura
\[ \begin{aligned}
\lambda(x)&=\frac{1}{2n+1}\{\lambda(x_{0})+\lambda(x_{1})+....+\lambda(x_{2n})\},\\
\lambda_{1}(x)&=\frac{1}{2n+1}\{\lambda_{1}(x_{0})+\lambda_{1}(x_{1})+....+\lambda_{1}(x_{2n})\}.
\end{aligned} \]
\ednote{The list of roots is printed as ending with $x_{2n+1}$, while the sums run only to $x_{2n}$.}
Or le second membre de ces équations est une fonction \emph{rationelle et symétrique} des racines de l'équation $\Phi(2n+1)\beta=\frac{P_{2n+1}}{Q_{2n+1}}$, donc $\lambda(x)$ et $\lambda_{1}(x)$ pourront s'exprimer rationellement en $\Phi(2n+1)\beta$. En faisant
\[ \lambda(x)=B,\ \lambda_{1}(x)=2A, \]
les équations (71.) donneront
\[ (\psi\beta)^{2n+1}.(\psi_{1}\beta)^{2n+1}=B^{2n+1};\ (\psi\beta)^{2n+1}+(\psi_{1}\beta)^{2n+1}=2A, \]
d'où l'on tire
\[ \psi\beta=\sqrt[2n+1]{[A+\sqrt{(A^{2}-B^{2n+1})}]}=\overset{+n}{\underset{-n}{\Sigma}}_{\mu}\theta^{\mu}.\Phi_{1}\left(\beta+\frac{2\mu\varpi i}{2n+1}\right). \tag{75.} \]

\subsection*{15.}

Ayant trouvé la valeur de $\psi\beta$, on en déduira facilement celles de $\Phi_{1}\beta$.

En effet, en prenant pour $\theta$ successivement toutes les racines imaginaires de l'équation $\theta^{2n+1}-1=0$, et désignant les valeurs correspondantes de $A$ et $B$ par $A_{1}$, $B_{1}$, $A_{2}$, $B_{2}$ etc., on obtiendra:
\[ \begin{aligned}
\sqrt[2n+1]{[A_{1}+\sqrt{(A_{1}^{2}-B_{1}^{2n+1})}]}&=\overset{+n}{\underset{-n}{\Sigma}}_{\mu}\theta_{1}^{\mu}.\Phi_{1}\left(\beta+\frac{2\mu\varpi i}{2n+1}\right),\\
\sqrt[2n+1]{[A_{2}+\sqrt{(A_{2}^{2}-B_{2}^{2n+1})}]}&=\overset{+n}{\underset{-n}{\Sigma}}_{\mu}\theta_{2}^{\mu}.\Phi_{1}\left(\beta+\frac{2\mu\varpi i}{2n+1}\right),\\
\cdots\cdots\cdots&\cdots\cdots\cdots\\
\sqrt[2n+1]{[A_{2n}+\sqrt{(A_{2n}^{2}-B_{2n}^{2n+1})}]}&=\overset{+n}{\underset{-n}{\Sigma}}_{\mu}\theta_{2n}^{\mu}.\Phi_{1}\left(\beta+\frac{2\mu\varpi i}{2n+1}\right).
\end{aligned} \]

\origpage{132}
De même on connoît la somme des racines:
\[ \overset{+n}{\underset{-n}{\Sigma}}_{m}\overset{+n}{\underset{-n}{\Sigma}}_{\mu}\varphi\left(\beta+\frac{2m\omega}{2n+1}+\frac{2\mu\varpi i}{2n+1}\right)=\overset{+n}{\underset{-n}{\Sigma}}_{\mu}\varphi_{1}\left(\beta+\frac{2\mu\varpi i}{2n+1}\right), \]
qui est égale à $(2n+1)\varphi(2n+1)\beta$, comme nous le verrons dans la suite. En ajoutant ces équations membre à membre, après avoir multiplié la première par $\theta_{0}^{-k}$, la seconde par $\theta_{1}^{-k}$, la troisième par $\theta_{2}^{-k}....$ et la $(2n)^{\text{ième}}$ par $\theta_{2n}^{-k}$, il viendra:
\[ \begin{aligned}
&\overset{+n}{\underset{-n}{\Sigma}}_{\mu}.(\theta_{0}^{\mu-k}+\theta_{1}^{\mu-k}+\theta_{2}^{\mu-k}+....+\theta_{2n}^{\mu-k}+1).\varphi_{1}\left(\beta+\frac{2\mu\varpi i}{2n+1}\right)\\
&=(2n+1).\varphi(2n+1)\beta+\overset{2n}{\underset{1}{\Sigma}}_{\mu}\theta_{\mu}^{-k}.\sqrt[2n+1]{[A_{\mu}+\sqrt{(A_{\mu}^{2}-B_{\mu}^{2n+1})}]};
\end{aligned} \]
\ednote{In the first sum the term $\theta_{0}^{\mu-k}$ and the final $1$ are both printed, so 1 is counted twice; the next sum is printed without $\theta_{0}^{\mu-k}$.}
or la somme
\[ 1+\theta_{1}^{\mu-k}+\theta_{2}^{\mu-k}+....+\theta_{2n}^{\mu-k} \]
se réduit à zéro pour toutes les valeurs de $k$, excepté pour $k=\mu$. Dans ce cas elle devient égale à $2n+1$. Donc le premier membre de l'équation précédente devient
\[ (2n+1).\varphi_{1}\left(\beta+\frac{2k\varpi i}{2n+1}\right), \]
donc, en substituant et divisant par $(2n+1)$, on a:
\[ \begin{aligned}
\varphi_{1}\left(\beta+\frac{2k\varpi i}{2n+1}\right)=&\\
\varphi(2n+1)\beta+\frac{1}{2n+1}.&\left\{\theta_{1}^{-k}.\sqrt[2n+1]{[A_{1}+\sqrt{(A_{1}^{2}-B_{1}^{2n+1})}]}+\theta_{2}^{-k}\sqrt[2n+1]{[A_{2}+\sqrt{(A_{2}^{2}-B_{2}^{2n+1})}]}+....\right.\\
&\left.....+\theta_{2n}^{-k}\sqrt[2n+1]{[A_{2n}+\sqrt{(A_{2n}^{2}-B_{2n}^{2n+1})}]}\right\}.
\end{aligned} \tag{76.} \]
Pour $k=0$, on a:
\[ \begin{aligned}
\varphi_{1}\beta=\varphi(2n+1)\beta+\frac{1}{2n+1}&\left\{\sqrt[2n+1]{[A_{1}+\sqrt{(A_{1}^{2}-B_{1}^{2n+1})}]}+\sqrt[2n+1]{[A_{2}+\sqrt{(A_{2}^{2}-B_{2}^{2n+1})}]}+\ldots\right.\\
&\left.....+\sqrt[2n+1]{[A_{2n}+\sqrt{(A_{2n}^{2}-B_{2n}^{2n+1})}]}\right\}.
\end{aligned} \tag{77.} \]

\subsection*{16.}

Ayant ainsi trouvé la valeur de $\varphi_{1}\beta$, il s'agit, d'en tirer celle de $\varphi\beta$. Or cela peut se faire aisément comme suit:

Soit
\[ \psi_{2}\beta=\overset{+n}{\underset{-n}{\Sigma}}_{m}\theta^{m}.\varphi\left(\beta+\frac{2m\omega}{2n+1}\right);\ \psi_{3}\beta=\overset{+n}{\underset{-n}{\Sigma}}_{m}\theta^{m}.\varphi\left(\beta-\frac{2m\omega}{2n+1}\right), \tag{78.} \]
on a
\[ \varphi\left(\beta\pm\frac{2m\omega}{2n+1}\right)=\frac{\varphi\beta.f\left(\frac{2m\omega}{2n+1}\right).F\left(\frac{2m\omega}{2n+1}\right)\pm f\beta.F\beta.\varphi\left(\frac{2m\omega}{2n+1}\right)}{1+e^{2}c^{2}.\varphi^{2}\left(\frac{2m\omega}{2n+1}\right).\varphi^{2}\beta}. \]

\origpage{133}
De là suit, qu'on peut faire
\[ \psi_{2}\beta=r+f\beta.F\beta.s;\ \psi_{3}\beta=r-f\beta.F\beta.s, \]
où $r$ et $s$ sont des fonctions rationelles de $\Phi\beta$.

De là on tire
\[ \left\{\begin{aligned}
\psi_{2}\beta.\psi_{3}\beta&=\chi(\Phi\beta),\\
(\psi_{2}\beta)^{2n+1}+(\psi_{3}\beta)^{2n+1}&=\chi_{1}(\Phi\beta),
\end{aligned}\right. \tag{79.} \]
$\chi(\Phi\beta)$ et $\chi_{1}(\Phi\beta)$ étant deux fonctions rationelles de $\Phi\beta$.

Cela posé je dis, que $\chi(\Phi\beta)$ et $\chi_{1}(\Phi\beta)$ pourront s'exprimer rationellement en $\Phi_{1}\beta$.

On a vu, que
\[ \varphi_{1}\beta=\varphi\beta+\overset{n}{\underset{1}{\Sigma}}_{m}.\left\{\frac{2\varphi\beta.f\left(\frac{2m\omega}{2n+1}\right).F\left(\frac{2m\omega}{2n+1}\right)}{1+e^{2}c^{2}.\varphi^{2}\left(\frac{2m\omega}{2n+1}\right)\varphi^{2}\beta}\right\}. \tag{80.} \]
En faisant $\Phi\beta=x$, on aura une équation en $x$ du degré $(2n+1)$. Une racine de cette équation est $x=\Phi\beta$; or, en mettant $\beta+\frac{2k\omega}{2n+1}$ au lieu de $\beta$, $\Phi_{1}\beta$ ne change pas de valeur, donc $x=\Phi\left(\beta+\frac{2k\omega}{2n+1}\right)$ sera une racine, quel que soit le nombre entier $k$. Or, en donnant à $k$ toutes les valeurs entières de $-n$ jusqu'á $+n$, $\Phi\left(\beta+\frac{2k\omega}{2n+1}\right)$ prendra $2n+1$ valeurs différentes, donc ces $2n+1$ quantités seront précisément les $2n+1$ racines de l'équation en $x$.

Cela posé, en mettant $\beta+\frac{2k\omega}{2n+1}$ au lieu de $\beta$ dans l'expression de $\psi_{2}\beta$, il viendra en vertu de (72.):
\[ \begin{aligned}
\psi_{2}\left(\beta+\frac{2k\omega}{2n+1}\right)&=\overset{+n}{\underset{-n}{\Sigma}}_{m}\theta^{m}\Phi\left(\beta+\frac{2(k+m)\omega}{2n+1}\right)\\
&=\theta^{-k}.\psi_{2}3+\overset{k}{\underset{1}{\Sigma}}_{m}\theta^{m+n-k}.\Phi\left(\beta+\frac{2(m+n)\omega}{2n+1}\right)\\
&\quad-\overset{k}{\underset{1}{\Sigma}}_{m}\theta^{m-n-1-k}.\Phi\left(\beta+\frac{2(m-n-1)\omega}{2n+1}\right),
\end{aligned} \]
\ednote{Printed $\psi_{2}3$ with a digit three instead of $\psi_{2}\beta$; an apparent misprint.}
donc en ayant attention, que $\theta^{m+n-k}=\theta^{m-n-1-k}$ et $\Phi\left(\beta+\frac{2(m-n-1)}{2n+1}\omega\right)=\Phi\left(\beta+\frac{2(m+n)\omega}{2n+1}\right)$, il en résultera
\[ \psi_{2}\left(\beta+\frac{2k\omega}{2n+1}\right)=\theta^{-k}.\psi_{2}\beta. \tag{81.} \]

\origpage{134}
De même on aura
\[ \psi_{3}\left(\beta+\frac{2k\omega}{2n+1}\right)=\theta^{+k}.\psi_{3}\beta. \]
On voit en vertu de ces relations, que les équations, qui donnent les valeurs des fonctions $\chi(\Phi\beta)$ et $\chi_{1}(\Phi\beta)$\ednote{The second printed argument looks like a damaged or misprinted $\beta$.}, conduisent à ces deux égalités:
\[ \begin{aligned}
\chi_{1}\left[\Phi\left(\beta+\frac{2k\omega}{2n+1}\right)\right]&=\chi_{1}(\Phi\beta),\\
\chi\left[\Phi\left(\beta+\frac{2k\omega}{2n+1}\right)\right]&=\chi(\Phi\beta).
\end{aligned} \]
De là on tire
\[ \begin{aligned}
\chi(\Phi\beta)&=\frac{1}{2n+1}.\overset{+n}{\underset{-n}{\Sigma}}_{k}.\chi\left[\Phi\left(\beta+\frac{2k\omega}{2n+1}\right)\right],\\
\chi_{1}(\Phi\beta)&=\frac{1}{2n+1}.\overset{+n}{\underset{-n}{\Sigma}}_{k}.\chi_{1}\left[\Phi\left(\beta+\frac{2k\omega}{2n+1}\right)\right].
\end{aligned} \]

Or, ces valeurs de $\chi(\Phi\beta)$ et $\chi_{1}(\Phi\beta)$ sont des fonctions rationelles et symétriques de toutes les racines de l'équation (80.). Donc elles peuvent être exprimées rationellement par les coefficiens de la même équation, c'est-à-dire rationellement en $\Phi_{1}\beta$.

Soit
\[ \chi_{1}(\Phi\beta)=2C,\ \chi(\Phi\beta)=D, \]
\ednote{The subscripts of both printed $\chi$ are damaged and hard to read; they are read here as $\chi_1$ for $2C$ and $\chi$ for $D$, as (79.) requires.}
les équations (79.) donneront
\[ \psi_{2}\beta=\sqrt[2n+1]{[C+\sqrt{(C^{2}-D^{2n+1})}]}, \]
d'où, en remettant la valeur de $\psi_{2}\beta$:
\[ \sqrt[2n+1]{[C+\sqrt{(C^{2}-D^{2n+1})}]}=\overset{+n}{\underset{-n}{\Sigma}}_{m}\theta^{m}.\Phi\left(\beta+\frac{2m\omega}{2n+1}\right). \tag{82.} \]
De là on tire, en mettant $\theta_{\mu}$ au lieu de $\theta$, et désignant les valeurs correspondantes de $C$ et $D$ par $C_{\mu}$ et $D_{\mu}$:
\[ \theta_{\mu}^{-k}\sqrt[2n+1]{[C_{\mu}+\sqrt{(C_{\mu}^{2}-D_{\mu}^{2n+1})}]}=\overset{+n}{\underset{-n}{\Sigma}}_{m}\theta_{\mu}^{m-k}.\Phi\left(\beta+\frac{2m\omega}{2n+1}\right). \]
En y joignant l'équation
\[ \Phi_{1}\beta=\overset{+n}{\underset{-n}{\Sigma}}_{m}\Phi\left(\beta+\frac{2m\omega}{2n+1}\right), \]
on en tirera facilement:
\[ (2n+1).\Phi\left(\beta+\frac{2k\omega}{2n+1}\right)=\Phi_{1}\beta+\overset{2n}{\underset{1}{\Sigma}}_{\mu}\theta_{\mu}^{-k}.\sqrt[2n+1]{[C_{\mu}+\sqrt{(C_{\mu}^{2}-D_{\mu}^{2n+1})}]}. \tag{83.} \]
En supposant $k=0$, il viendra:
\[ \Phi\beta=\frac{1}{2n+1}\left\{\Phi_{1}\beta+\sqrt[2n+1]{[C_{1}+\sqrt{(C_{1}^{2}-D_{1}^{2n+1})}]}+\ldots+\sqrt[2n+1]{[C_{2n}+\sqrt{(C_{2n}^{2}-D_{2n}^{2n+1})}]}\right\}. \tag{84.} \]

\origpage{135}

Cette équation donne $\Phi\beta$ en fonction algébrique de $\Phi_{1}\beta$, or précédemment nous avons trouvé $\Phi_{1}\beta$ en fonction algébrique de $\Phi(2n+1)\beta$. Donc en mettant $\frac{\alpha}{2n+1}$ au lieu de $\beta$, on aura $\Phi\left(\frac{\alpha}{2n+1}\right)$ en fonction algébrique de $\Phi\alpha$.

Par une analyse toute semblable on trouvera $f\left(\frac{\alpha}{2n+1}\right)$ en $f\alpha$ et $F\left(\frac{\alpha}{2n+1}\right)$ en $F\alpha$.

\subsection*{17.}

Les valeurs, que nous venons de trouver des quantités $\Phi_{1}\beta$ et $\Phi\beta$, la première en $\Phi(2n+1)\beta$ et la seconde en $\Phi_{1}\beta$, contiennent chacune la somme de $2n$ radicaux différentes du $(2n+1)$ degré. Il en résultera pour $\Phi\beta$, $\Phi_{1}\beta....$ un nombre $(2n+1)^{2n}$ de valeurs, tandisque chacune de ces quantités est la racine d'une équation du $(2n+1)^{\text{ème}}$ degré. Or on peut donner aux expressions de $\Phi\beta$ et $\Phi_{1}\beta$ une telle forme, que le nombre des valeurs de ces quantités soit précisément égal à $2n+1$.

Pour cela soit
\[ \theta=\cos\frac{2\pi}{2n+1}+i\sin\frac{2\pi}{2n+1}, \]
on peut faire
\[ \theta_{1}=\theta,\ \theta_{2}=\theta^{2},\ \theta_{3}=\theta^{3}....\theta_{2n}=\theta^{2n}. \]

Soient de même
\[ \left\{\begin{aligned}
\psi^{k}(\beta)&=\overset{+n}{\underset{-n}{\Sigma}}_{\mu}\theta^{k\mu}.\Phi_{1}\left(\beta+\frac{2\mu\varpi i}{2n+1}\right),\\
\psi_{1}^{k}(\beta)&=\overset{+n}{\underset{-n}{\Sigma}}_{\mu}\theta^{k\mu}.\Phi_{1}\left(\beta-\frac{2\mu\varpi i}{2n+1}\right),
\end{aligned}\right. \tag{85.} \]
on aura en vertu de (74.)
\[ \begin{aligned}
\psi^{k}\left(\beta+\frac{2\nu\varpi i}{2n+1}\right)&=\theta^{-k\nu}.\psi^{k}(\beta),\\
\psi_{1}^{k}\left(\beta+\frac{2\nu\varpi i}{2n+1}\right)&=\theta^{+k\nu}.\psi_{1}^{k}(\beta),\\
\psi^{1}\left(\beta+\frac{2\nu\varpi i}{2n+1}\right)&=\theta^{-\nu}.\psi^{1}(\beta),\\
\psi_{1}^{1}\left(\beta+\frac{2\nu\varpi i}{2n+1}\right)&=\theta^{+\nu}.\psi_{1}^{1}(\beta).
\end{aligned} \]

Soit maintenant
\[ \left\{\begin{aligned}
\frac{\psi^{k}(\beta)}{(\psi^{1}\beta)^{k}}+\frac{\psi_{1}^{k}(\beta)}{(\psi_{1}^{1}\beta)^{k}}&=P(\Phi\beta),\\
\frac{\psi^{k}(\beta)}{(\psi^{1}\beta)^{k-2n-1}}+\frac{\psi_{1}^{k}(\beta)}{\psi_{1}^{1}(\beta)^{k-2n-1}}&=Q(\Phi\beta),
\end{aligned}\right. \tag{86.} \]

\origpage{136}
$P(\Phi\beta)$ et $Q(\Phi\beta)$ seront des fonctions rationelles de $\Phi\beta$; or, en mettant $\beta+\frac{2m\omega+2\mu\varpi i}{2n+1}$ au lieu de $\beta$, il est clair en vertu des formules, que $P$ et $Q$ ne changent pas de valeurs; donc on aura
\[ P(\Phi\beta)=\frac{1}{(2n+1)^{2}}.\overset{+n}{\underset{-n}{\Sigma}}_{m}\overset{+n}{\underset{-n}{\Sigma}}_{\mu}.P\left[\Phi\left(\beta+\frac{2m\omega+2\mu\varpi i}{2n+1}\right)\right]; \]
or, le second membre étant une fonction symétrique et rationelle des racines de l'équation $\Phi(2n+1)\beta=\frac{P_{2n+1}}{Q_{2n+1}}$, $P(\Phi\beta)$ pourra s'exprimer rationellement en $\Phi(2n+1)\beta$. Il en est de même de $Q(\Phi\beta)$. Connoissant ces deux quantités, les équations (86.) donneront
\[ \frac{\psi^{k}(\beta)}{(\psi^{1}\beta)^{k}}\left\{1-\left(\frac{\psi_{1}^{1}\beta}{\psi_{1}^{1}\beta}\right)^{2n+1}\right\}=P(\Phi\beta)-\frac{Q(\varphi\beta)}{(\psi_{1}^{1}\beta)^{2n+1}}, \]
or
\[ \begin{aligned}
\psi_{1}^{1}\beta&=A_{1}+\sqrt{(A_{1}^{2}-B_{1}^{2n+1})},\\
\psi_{1}^{1}\beta&=A_{1}-\sqrt{(A_{1}^{2}-B_{1}^{2n+1})},
\end{aligned} \]
\ednote{In the first of these two equations the subscript of the printed $\psi$ is damaged; it should read $\psi^{1}\beta$ without subscript, and the quotient in the preceding display should likewise have $\psi^{1}\beta$ in the numerator.}
donc
\[ \frac{\psi^{k}(\beta)}{(\psi^{1}\beta)^{k}}.2\sqrt{(A_{1}^{2}-B_{1}^{2n+1})}=Q(\Phi\beta)-[A_{1}-\sqrt{(A_{1}^{2}-B_{1}^{2n+1})}].P(\Phi\beta). \]
Donc on aura
\[ \psi^{k}(\beta)=(\psi^{1}\beta)^{k}.\{F_{k}+H_{k}.\sqrt{(A_{1}^{2}-B_{1}^{2n+1})}\}, \]
où $F_{k}$ et $H_{k}$ sont des fonctions rationelles de $\Phi(2n+1)\beta$. En remplaçant $A_{1}$ et $B_{1}$ et substituant les valeurs de $\psi^{k}(\beta)$ et $(\psi^{1}\beta)^{k}$, il viendra:
\[ \sqrt[2n+1]{[A_{k}+\sqrt{(A_{k}^{2}-B_{k}^{2n+1})}]}=[A+\sqrt{(A^{2}-B^{2n+1})}]^{\frac{k}{2n+1}}[F_{k}+H_{k}.\sqrt{(A^{2}-B^{2n+1})}], \]
donc la valeur de $\Phi_{1}\beta$ deviendra:
\[ \begin{aligned}
\Phi_{1}\beta=\Phi(2n+1)\beta+\frac{1}{2n+1}.&\left\{[A+\sqrt{(A^{2}-B^{2n+1})}]^{\frac{1}{2n+1}}\right.\\
&+[F_{2}+H_{2}\sqrt{(A^{2}-B_{2n}^{2n+1})}]\,[A+\sqrt{(A^{2}-B^{2n+1})}]^{\frac{2}{2n+1}}\\
&\left.+....[F_{2n}+H_{2n}\sqrt{(A^{2}-B^{2n+1})}]\,[A+\sqrt{(A^{2}-B^{2n+1})}]^{\frac{2n}{2n+1}}\right\}.
\end{aligned} \tag{87.} \]
\ednote{In the second line the printed $B$ under the radical carries the subscript $2n$, which the first and third lines lack; an apparent misprint.}
Par un procédé tout semblable on trouvera
\[ \begin{aligned}
\Phi\beta=\frac{1}{2n+1}&\left\{\Phi_{1}\beta+[C+\sqrt{(C^{2}-D^{2n+1})}]^{\frac{1}{2n+1}}\right.\\
&+[K_{2}+L_{2}\sqrt{(C^{2}-D^{2n+1})}]\,[C+\sqrt{(C^{2}-D^{2n+1})}]^{\frac{2}{2n+1}}\\
&\left.+....[K_{2n}+L_{2n}\sqrt{(C^{2}-D^{2n+1})}]\,[C+\sqrt{(C^{2}-D^{2n+1})}]^{\frac{2n}{2n+1}}\right\},
\end{aligned} \tag{88.} \]
où $K_{2}$, $L_{2}$, $K_{3}$, $L_{3}....K_{2n}$, $L_{2n}$ sont des fonctions rationelles de $\Phi_{1}\beta$.

\origpage{137}

Ces expressions de $\Phi_{1}\beta$ et $\Phi\beta$ n'ont que $2n+1$ valeurs différentes, qu'on obtiendra en attribuant aux radicaux leurs $2n+1$ valeurs. --- Il suit de notre analyse, qu'on peut prendre $\sqrt{(A^{2}-B^{2n+1})}$ et $\sqrt{(C^{2}-D^{2n+1})}$ avec tel signe qu'on voudra.

\subsection*{18.}

La valeur, que nous avons trouvée pour $\Phi(\beta)$ ou $\Phi\left(\frac{\alpha}{2n+1}\right)$ contient, encore outre la fonction $\Phi\alpha$, les suivantes:
\[ \begin{gathered}
e,\ c,\ \theta,\\
\Phi\left(\frac{m\omega}{2n+1}\right),\ \Phi\left(\frac{m\varpi i}{2n+1}\right),\ f\left(\frac{m\omega}{2n+1}\right),\\
f\left(\frac{m\varpi i}{2n+1}\right),\ F\left(\frac{m\omega}{2n+1}\right),\ F\left(\frac{m\varpi i}{2n+1}\right),
\end{gathered} \]
pour des valeurs quelconques de $m$ depuis $1$ jusqu'à $2n$. Maintenant quelle que soit la valeur de $m$, on peut toujours exprimer algébriquement $\Phi\left(\frac{m\omega}{2n+1}\right)$, $f\left(\frac{m\omega}{2n+1}\right)$, $F\left(\frac{m\omega}{2n+1}\right)$ en $\Phi\left(\frac{\omega}{2n+1}\right)$, et $\Phi\left(\frac{\mu\varpi i}{2n+1}\right)$, $f\left(\frac{\mu\varpi i}{2n+1}\right)$, $F\left(\frac{\mu\varpi i}{2n+1}\right)$ en $\Phi\left(\frac{\varpi i}{2n+1}\right)$. Tout est donc connu dans l'expression de $\Phi\left(\frac{\alpha}{2n+1}\right)$, excepté les deux quantités indépendantes de $\alpha$; $\Phi\left(\frac{\omega}{2n+1}\right)$, $\Phi\left(\frac{\varpi i}{2n+1}\right)$. --- Ces quantités dépendent seulement de $c$ et $e$, et elles peuvent être trouvées par la résolution d'une équation du degré $(2n+1)^{2}-1$, savoir de l'équation $\frac{P_{2n+1}}{x}=0$. --- Nous allons voir dans le paragraphe suivant, comment on peut en ramener la résolution à celle d'équations moins élevées.

\section*{§. V.}

C. \emph{Sur l'équation} $P_{2n+1}=0$.

\subsection*{19.}

L'expression, que nous venons de trouver pour $\Phi\left(\frac{\alpha}{2n+1}\right)$ contiendra, comme nous avons vu, les deux quantités constantes $\Phi\left(\frac{\omega}{2n+1}\right)$ et $\Phi\left(\frac{\varpi i}{2n+1}\right)$. On trouvera ces quantités en résolvant l'équation
\[ P_{2n+1}=0, \]
dont les racines seront représentées par
\[ x=\Phi\left(\frac{m\omega+\mu\varpi i}{2n+1}\right), \]

\origpage{138}
où $m$ et $\mu$ pourront être tous les nombres entiers depuis $-n$ jusqu'à $+n$. Une de ces racines, qui répond à $m=0$, $\mu=0$, est égale à zéro. Donc $P_{2n+1}$ est divisible par $x$. En écartant ce facteur, on aura une équation
\[ R=0,\ \text{du degré }(2n+1)^{2}-1. \]
En faisant $x^{2}=r$, l'équation (90.) $R=0$, en $r$, sera du degré $\frac{(2n+1)^{2}-1}{2}=2.n(n+1)$, et les racines de cette équation seront
\[ r=\Phi^{2}\left(\frac{m\omega\pm\mu\varpi i}{2n+1}\right), \tag{91.} \]
$\mu$ et $m$ ayant toutes les valeurs positives, au dessous de $n$, en faisant abstraction de la racine zéro.

Nous allons voir maintenant, comment on peut ramener la résolution de l'équation $R=0$ à celle de deux équations, l'une du degré $n$ et l'autre du degré $2n+2$.

D'abord, je dis, qu'on peut représenter toutes les valeurs de $r$ par
\[ \Phi^{2}\left(\frac{m\omega}{2n+1}\right)\ \text{et}\ \Phi^{2}\left(m\frac{\mu\omega+\varpi i}{2n+1}\right)\ldots\ldots, \tag{92.} \]
en donnant à $\mu$ toutes les valeurs entières depuis zéro jusqu'à $2n$, et à $m$ toutes celles depuis $1$ jusqu'à $n$.

En effet $\Phi^{2}\left(\frac{m\omega}{2n+1}\right)$ représente d'abord un nombre $n$ de valeurs de $r$; or les autres peuvent être représentées par $\Phi^{2}\left(m.\frac{\mu\omega+\varpi i}{2n+1}\right)$. Soit, pour le démontrer, $m\mu=(2n+1)k+m'$, où $m'$ est un nombre entier compris entre les limites $-n$ et $+n$. En substituant, on aura
\[ \begin{aligned}
\Phi^{2}\left(m.\frac{\mu\omega+\varpi i}{2n+1}\right)&=\Phi^{2}\left(k.\omega+\frac{m'\omega+\mu\varpi i}{2n+1}\right)\\
&=\Phi^{2}\left(\frac{m'\omega+\mu\varpi i}{2n+1}\right)=\Phi^{2}\left(\frac{-m'\omega-\mu\varpi i}{2n+1}\right).
\end{aligned} \]
$\Phi^{2}\left(m.\frac{\mu\omega+\varpi i}{2n+1}\right)$ est donc une valeur de $r$; maintenant à chaque valeur de $m$, répond une valeur différente de $m'$. Car si l'on avoit
\[ m_{1}\mu=(2n+1)k_{1}+m', \]
il s'en suivrait
\[ (m_{1}-m)\mu=(2n+1).(k-k_{1}), \]
ce qui est impossible, en remarquant que $2n+1$ est un nombre premier. Donc $\Phi^{2}\left(m.\frac{\mu\omega+\varpi i}{2n+1}\right)$ combiné avec $\Phi^{2}\left(\frac{m\omega}{2n+1}\right)$ représente toutes les valeurs de $r$.

\origpage{139}

Cela posé, soit
\[ \begin{aligned}
&\left[r-\Phi^{2}\left(\frac{\omega}{2n+1}\right)\right]\left[r-\Phi^{2}\left(\frac{2\omega}{2n+1}\right)\right]\ldots\left[r-\Phi^{2}\left(\frac{n\omega}{2n+1}\right)\right]\\
&=r^{n}+p_{n-1}.r^{n-1}+p_{n-2}.r^{n-2}+\ldots\ldots+p_{1}.r+p_{0}.
\end{aligned} \tag{93.} \]
Les quantités $p_{0}$, $p_{1}$, $\ldots p_{n-1}$, seront des fonctions rationelles et symétriques de $\Phi^{2}\left(\frac{\omega}{2n+1}\right)$, $\Phi^{2}\left(\frac{2\omega}{2n+1}\right)\ldots\Phi^{2}\left(\frac{n\omega}{2n+1}\right)$; or ces fonctions peuvent être trouvées au moyen d'une équation de degré $2n+2$.

Soit $p$ une fonction rationelle et symétrique quelconque de $\Phi^{2}\left(\frac{\omega'}{2n+1}\right)$, $\Phi^{2}\left(\frac{2\omega'}{2n+1}\right)\ldots\Phi^{2}\left(\frac{n\omega'}{2n+1}\right)$, où $\omega'$ désigne la quantité $m\omega+\mu\varpi i$.

Par les formules que nous avons données plus haut pour exprimer $\Phi(n\beta)$ en $\Phi\beta$, il est clair, qu'on peut exprimer $\Phi^{2}\left(m'.\frac{\omega'}{2n+1}\right)$ en fonction rationelle de $\Phi^{2}\left(\frac{\omega'}{2n+1}\right)$. Donc on peut faire
\[ p=\psi.\left[\Phi^{2}\left(\frac{\omega'}{2n+1}\right)\right]=\theta\left[\Phi^{2}\left(\frac{\omega'}{2n+1}\right),\ \Phi^{2}\left(\frac{2\omega'}{2n+1}\right),\ \ldots\ \Phi^{2}\left(\frac{n\omega'}{2n+1}\right)\right], \tag{94.} \]
$\theta$ désignant une fonction symétrique et rationelle. En mettant $\nu\omega'$ au lieu de $\omega'$, il viendra
\[ \psi\left[\Phi^{2}\left(\frac{\nu\omega'}{2n+1}\right)\right]=\theta\left[\Phi^{2}\left(\frac{\nu\omega'}{2n+1}\right),\ \Phi^{2}\left(\frac{2\nu\omega'}{2n+1}\right),\ \ldots\ \Phi^{2}\left(\frac{n\nu\omega'}{2n+1}\right)\right], \tag{95.} \]
or, en faisant:
\[ a.\nu=(2n+1).k_{a}'+k_{a}, \]
où $k_{a}$ est entier et compris entre $-n$ et $+n$, la serie
\[ k_{1},\ k_{2}\ \ldots\ldots\ k_{n} \]
aura au signe près les mêmes termes que celle-ci:
\[ 1,\ 2,\ 3\ \ldots\ldots\ n; \]
donc il est clair, que le second membre de l'équation (95.) aura la même valeur que $p$. Donc:
\[ \psi\left[\Phi^{2}\left(\frac{\nu\omega'}{2n+1}\right)\right]=\psi\left[\Phi^{2}\left(\frac{\omega'}{2n+1}\right)\right], \tag{96.} \]
\ednote{The exponent of the second Phi in (96.) is damaged in the print and appears as a raised dot; read as the exponent 2 of the first member.}
équation, qui en faisant $\omega'=\omega$ et $\omega'=m\omega+\varpi i$, donnera ces deux-ci:
\[ \left\{\begin{aligned}
&\psi\left[\Phi^{2}\left(\frac{\nu\omega}{2n+1}\right)\right]=\psi.\left[\Phi^{2}\left(\frac{\omega}{2n+1}\right)\right],\\
&\psi\left[\Phi^{2}\left(\nu.\frac{m\omega+\varpi i}{2n+1}\right)\right]=\psi\left[\Phi^{2}\left(\frac{m\omega+\varpi i}{2n+1}\right)\right],
\end{aligned}\right. \tag{97.} \]
ou bien, en faisant, pour abréger,
\[ \Phi^{2}\left(\frac{\nu\omega}{2n+1}\right)=r_{\nu},\ \Phi^{2}\left(\nu.\frac{m\omega+\varpi i}{2n+1}\right)=r_{\nu,m}, \tag{98.} \]

\origpage{140}
il viendra:
\[ \psi.r_{\nu}=\psi r_{1};\ \psi r_{\nu,m}=\psi r_{1,m}. \tag{99.} \]

Cela posé, soit
\[ \left\{\begin{aligned}
&(p-\psi r_{1})(p-\psi r_{1,0})(p-\psi r_{1,1})(p-\psi r_{1,2})\ldots\ (p-\psi r_{1,2n})\\
&=q_{0}+q_{1}.p+q_{2}.p^{2}+\ldots q_{2n+1}.p^{2n+1}+p^{2n}.
\end{aligned}\right. \tag{100.} \]
\ednote{The exponent of the last term is partly obscured by a smudge and reads $2n$; the degree $2n+2$ is expected from (104.).}

Je dis, qu'on peut exprimer les coefficiens $q_{0}$, $q_{1}$ etc. rationellement en $e$ et $c$.

D'abord en vertu des formules connues on peut exprimer rationellement ces coefficiens en $t_{1}$, $t_{2}\ldots t_{2n}$, si l'on fait, pour abréger,
\[ t_{k}=(\psi r_{1})^{k}+(\psi r_{1,0})^{k}+(\psi r_{1,1})^{k}+\ldots+(\psi r_{1,2n})^{k}. \tag{101.} \]

Il s'agit donc de trouver les quantités $t_{1}$, $t_{2}\ldots$, or cela se pourra aisément au moyen des relations (99.). En effet, en y faisant successivement $\nu=1$, $2\ldots n$, après avoir élevé les deux membres à la $k^{\text{ème}}$ puissance, on en tirera sur le champ:
\[ \left\{\begin{aligned}
(\psi r_{1})^{k}&=\frac{1}{n}\left[(\psi r_{1})^{k}+(\psi r_{2})^{k}+\ldots+(\psi r_{n})^{k}\right],\\
(\psi r_{1,m})^{k}&=\frac{1}{n}\left[(\psi r_{1,m})^{k}+(\psi r_{2,m})^{k}+\ldots+(\psi r_{n,m})^{k}\right].
\end{aligned}\right. \tag{102.} \]
Donc en mettant pour $m$ tous les nombres entiers $0$, $1$, $\ldots$ $2n$, et ensuite substituant dans l'expression de $t_{k}$, il viendra:
\[ \left\{\begin{aligned}
n.t_{k}&=(\psi r_{1})^{k}+(\psi r_{2})^{k}+\ldots\ldots+(\psi r_{n})^{k}\\
&+(\psi r_{1,0})^{k}+(\psi r_{2,0})^{k}+\ldots\ldots+(\psi r_{n,0})^{k}\\
&+(\psi r_{1,1})^{k}+(\psi r_{2,1})^{k}+\ldots\ldots+(\psi r_{n,1})^{k}\\
&+\cdots\cdots\cdots\cdots\cdots\cdots\cdots\\
&+(\psi r_{1,2n})^{k}+(\psi r_{2,2n})^{k}+\ldots\ldots+(\psi r_{n,2n})^{k}.
\end{aligned}\right. \tag{103.} \]

Cette valeur de $t_{k}$ est, comme on voit, une fonction rationelle et symétrique des $2n+2$ quantités $r_{1}$, $r_{2}$, $\ldots r_{n}$; $r_{1,0}$, $r_{2,0}$, $\ldots r_{n,0}\ldots$ $r_{1,2n}$, $r_{2,2n}\ldots r_{n,2n}$, qui sont les $2n+2$ racines de l'équation $R=0$. Donc, comme on sait, $t_{k}$ pourra s'exprimer rationellement par les coefficiens de cette équation, et par suite en fonction rationelle de $e$ et $c$. Ayant ainsi trouvé les quantités $t_{k}$, on en tire les valeurs de $q_{0}$, $q_{1}$, $\ldots$ $\ldots q_{2n+1}$; qui seront également des fonctions rationelles de $e$ et $c$.

\subsection*{20.}

Cela posé, en supposant
\[ 0=q_{0}+q_{1}.p+q_{2}.p^{2}+\ldots\ldots+q_{2n+1}.p^{2n+1}+p^{2n+2}, \tag{104.} \]
on aura une équation du $(2n+2)^{\text{ème}}$ degré, dont les racines seront
\[ \psi r_{1};\ \psi r_{1,0};\ \psi r_{1,1};\ \psi r_{1,2}\ \ldots\ldots\ \psi r_{1,2n}. \]

\origpage{141}

La fonction $\psi r_{1}$, c'est-à-dire, une fonction quelconque rationelle et symétrique des racines $r_{1}$, $r_{2}$, $r_{3}$, $\ldots\ldots r_{n}$ pourra donc être trouvée au moyen d'une équation du degré $2n+2$.

Donc on aura de cette manière les coefficiens $p_{0}$, $p_{1}$, $\ldots\ldots p_{n-1}$, en résolvant un nombre $n$ d'équations, chacune du $(2n+2)^{\text{ème}}$ degré.

Ayant déterminé $p_{0}$, $p_{1}$, $\ldots\ldots$, on aura, en résolvant l'équation
\[ 0=p_{0}+p_{1}r+\ldots\ldots+p_{n-1}.r^{n-1}+r^{n}, \tag{105.} \]
la valeur des quantités
\[ r_{1},\ r_{2},\ \ldots\ldots r_{n};\ r_{1,0},\ r_{2,0},\ \ldots\ldots r_{n,0};\ r_{1,1},\ r_{2,1},\ \ldots\ldots r_{n,1}\ \text{etc. etc.,} \]
dont la première est égale à $\Phi^{2}\left(\frac{\omega}{2n+1}\right)$. Donc la détermination de cette quantité, ou bien la résolution de l'équation $R=0$, qui est du degré $(2n+2).n$, est réduite à celle d'équations du degré $(2n+2)$ et $n$.

Mais on peut encore simplifier le procédé précédent. En effet, comme nous le verrons, pour avoir les quantités $p_{0}$, $p_{1}$, $\ldots\ldots$, il suffit de connoître l'une quelconque d'entre elles, et alors on peut exprimer les autres rationellement par celle-là.

Soient généralement $p$, $q$ deux fonctions rationelles et symétriques des quantités $r_{1}$, $r_{2}$, $\ldots\ldots r_{n}$, on peut faire, comme nous l'avons vu:
\[ p=\psi r_{1};\ q=\theta r_{1}, \]
$\psi r_{1}$ et $\theta r_{1}$ désignent deux fonctions rationelles de $r_{1}$, qui ont cette propriété: de rester les mêmes, si l'on change $r_{1}$ en une autre quelconque des quantités $r_{1}$, $r_{2}$, $\ldots\ldots r_{n}$.

Supposons maintenant:
\[ s_{k}=(\psi r_{1})^{k}.\theta r_{1}+(\psi r_{1,0})^{k}.\theta r_{1,0}+(\psi r_{1,1})^{k}.\theta r_{1,1}+\ldots\ldots+(\psi r_{1,2n})^{k}.\theta r_{1,2n}, \]
je dis que $s_{k}$ pourra être exprimé rationellement en $e$ et $c$.

En effet, on a
\[ (\psi r_{1})^{k}.\theta r_{1}=(\psi r_{\nu})^{k}.\theta r_{\nu}=\frac{1}{n}\left\{(\psi r_{1})^{k}.\theta r_{1}+(\psi r_{2})^{k}.\theta r_{2}+\ldots\ldots+(\psi r_{n})^{k}.\theta r_{n}\right\}, \]
\[ (\psi r_{1,m})^{k}.\theta r_{1,m}=(\psi r_{\nu,m})^{k}.\theta r_{\nu,m}=\frac{1}{n}\left\{(\psi r_{1,m})^{k}.\theta r_{1,m}+(\psi r_{2,m})^{k}.\theta r_{2,m}+\ldots+(\psi r_{n,m})^{k}.\theta r_{n,m}\right\}. \]
En faisant $m=0$, $1$, $2$, $\ldots\ldots 2n$, et substituant dans l'expression de $s_{k}$, on verra, que $s_{k}$ sera une fonction rationelle et symétrique des racines $r_{0}$, $r_{1}$, $\ldots\ldots r_{1,0}$ etc., etc. de l'équation $R=0$; donc $s_{k}$ pourra s'exprimer rationellement en $e$ et $c$.

Connoissant $s_{k}$, on obtiendra, en faisant $k=0$, $1$, $2$, $\ldots\ldots 2n$, $2n+1$ équations, desquelles on tirera aisément la valeur de $\theta r_{1}$ en fonction

\origpage{142}
rationelle de $\psi_{1}r$.\ednote{Printed $\psi_{1}r$ with the subscript on the psi; $\psi r_{1}$ is expected, as on the previous page.} Donc, une fonction de la forme $p$ étant donnée, on peut exprimer une autre fonction quelconque de la même forme en fonction rationelle de $p$. Donc, comme nous l'avons dit, on peut exprimer les coefficiens $p_{0}$, $p_{1}$, $\ldots\ldots p_{n-1}$ rationellement par l'un quelconque d'entre eux. Donc enfin, pour en avoir les valeurs, il suffit de résoudre une seule équation du degré $2n+2$, et par conséquent, pour avoir les racines de l'équation $R=0$, il suffit de résoudre une équation du degré $2n+2$, et $2n+2$ équations du degré $n$.

\subsection*{21.}

Maintenant, parmi les équations, dont dépend la détermination des quantités $\varphi\left(\frac{\omega}{2m+1}\right)$; $\varphi\left(\frac{\varpi i}{2n+1}\right)$,\ednote{Printed $2m+1$ in the first denominator; $2n+1$ is apparently expected.} celles du degré $n$ peuvent être résolues algébriquement. Le procédé, par lequel nous allons effectuer cette résolution est entièrement semblable à celui, qui est dû à M. \emph{Gauss} pour la résolution de l'équation
\[ \theta^{2n+1}-1=0. \]
Soit proposée l'équation
\[ 0=p_{1}+p_{1}.r+p_{2}.r^{2}+\ldots\ldots+p_{n-1}.r^{n-1}+r^{n}, \tag{106.} \]
\ednote{The first coefficient is printed $p_{1}$; $p_{0}$ is expected.}
dont les racines sont:
\[ \Phi^{2}\left(\frac{\omega'}{2n+1}\right);\ \Phi^{2}\left(\frac{2\omega'}{2n+1}\right);\ \ldots\ldots\ \Phi^{2}\left(\frac{n\omega'}{2n+1}\right), \]
où $\omega'$ a une des valeurs de $\omega$, $m\omega+\varpi i$. Désignons par $\alpha$ une des racines primitives du nombre $2n+1$, c'est-à-dire, un nombre entier tel, que $\mu=2n$\uncertain{$\,+1$}\ednote{The characters after $2n$ are overstruck by a heavy bar and hard to read; $\mu=2n$ is expected.} est le nombre le plus petit, qui rende $\alpha^{\mu}-1$ divisible par $2n+1$, je dis, que les racines de l'équation (106.) peuvent aussi être représentées par
\[ \Phi^{2}(\varepsilon),\ \Phi^{2}(\alpha\varepsilon),\ \Phi^{2}(\alpha^{2}\varepsilon),\ \Phi^{2}(\alpha^{3}\varepsilon)\ldots\ \Phi^{2}(\alpha^{n-1}.\varepsilon), \tag{107.} \]
où $\varepsilon=\frac{\omega'}{2n+1}$.

Soit
\[ \alpha^{m}=(2n+1)k_{m}\pm a_{m}, \]
où $k$ est entier et $a_{m}$ entier, positif et moindre que $n+1$, je dis, que les termes de la série
\[ 1,\ a_{1},\ a_{2},\ \ldots\ldots\ a_{n-1} \]
seront tous différens entre eux.

En effet, si l'on a
\[ a_{m}=a_{\mu}, \]
il en résulte,

\origpage{143}
\[ \begin{aligned}
\text{ou }\ \alpha^{m}-\alpha^{\mu}&=(2n+1)(k_{m}-k_{\mu}),\\
\text{ou }\ \alpha^{m}+\alpha^{\mu}&=(2n+1)(k_{m}+k_{\mu}).
\end{aligned} \]
Il faut donc que l'une des quantités $\alpha^{m}-\alpha^{\mu}$, $\alpha^{m}+\alpha^{\mu}$ soit divisible par $2n+1$; or soit posé $m>\mu$, ce qui est permis, il faut que $\alpha^{m-\mu}-1$ ou $\alpha^{m-\mu}+1$ soit divisible par $n$,\ednote{Printed $n$ here; $2n+1$ is apparently expected.} or cela est impossible, car $m-\mu$ est moindre que $n$.

Donc les quantités $1$, $a_{1}$, $a_{2}\ldots a_{n-1}$ sont différentes entre elles et par conséquent elles coincident, mais dans un ordre différent, avec les nombres $1$, $2$, $3$, $4\ldots n$.

Donc, en remarquant que
\[ \varphi^{2}\{[(2n+1)k_{m}\pm a_{m}]\varepsilon\}=\varphi^{2}(a_{m}\varepsilon), \]
on voit que les quantités (107.) sont les mêmes que celles-ci:
\[ \Phi^{2}(\varepsilon),\ \Phi^{2}(2\varepsilon)\ \ldots\ \Phi^{2}(n\varepsilon), \]
c'est-à-dire les racines de l'équation (106.) c.\ q.\ f.\ d.

Il y a encore à remarquer, qu'ayant
\[ \alpha^{n}=(2n+1)k_{n}-1, \]
on aura
\[ \alpha^{n+m}=(2n+1)k_{n}\alpha^{m}-\alpha^{m}, \]
donc
\[ a_{n+m}=\pm a_{m} \]
et
\[ \Phi^{2}(\alpha^{n+m}\varepsilon)=\Phi^{2}(\alpha^{m}\varepsilon). \]
Cela posé, soit $\theta$ une racine imaginaire quelconque de l'équation
\[ \theta^{n}-1=0 \]
et
\[ \psi(\varepsilon)=\Phi^{2}(\varepsilon)+\Phi^{2}(\alpha\varepsilon)\theta+\Phi^{2}(\alpha^{2}\varepsilon)\theta^{2}+\ldots+\Phi^{2}(\alpha^{n-1}\varepsilon)\theta^{n-1}. \tag{108.} \]
En vertu de ce que nous avons vu précédemment, le second membre de cette équation peut être transformé en une fonction \emph{rationelle} de $\Phi^{2}(\varepsilon)$. Faisons
\[ \psi(\varepsilon)=\chi[\Phi^{2}(\varepsilon)]. \tag{109.} \]
En mettant dans la première expression de $\psi(\varepsilon)$, $\alpha^{m}\varepsilon$ au lieu de $\varepsilon$, il viendra:
\[ \begin{aligned}
\psi(\alpha^{m}\varepsilon)&=\varphi^{2}(\alpha^{m}\varepsilon)+\varphi^{2}(\alpha^{m+1}\varepsilon).\theta+\varphi^{2}(\alpha^{m+2}\varepsilon).\theta^{2}+\ldots\\
&\ldots+\Phi^{2}(\alpha^{n-1}\varepsilon).\theta^{n-m-1}+\Phi^{2}(\alpha^{n}\varepsilon).\theta^{n-m}+\ldots+\Phi^{2}(\alpha^{n+m-1}\varepsilon)\theta^{n-1},
\end{aligned} \]
mais nous avons vu, que $\Phi^{2}(\alpha^{n+m}\varepsilon)=\Phi^{2}(\alpha^{m}\varepsilon)$; donc:
\[ \begin{aligned}
\psi(\alpha^{m}\varepsilon)&=\theta^{n-m}.\Phi^{2}(\varepsilon)+\theta^{n-m+1}.\Phi^{2}(\alpha\varepsilon)+\theta^{n-m+2}.\Phi^{2}(\alpha^{2}\varepsilon)+\ldots\\
&\ldots+\theta^{n-1}.\Phi^{2}(\alpha^{m-1}\varepsilon)+\Phi^{2}(\alpha^{m}\varepsilon)+\theta.\Phi^{2}(\alpha^{m+1}\varepsilon)+\ldots+\theta^{n-m-1}.\Phi^{2}(\alpha^{n-1}\varepsilon).
\end{aligned} \]

\origpage{144}
En multipliant par $\theta^{m}$, le second membre deviendra égal a $\psi(\varepsilon)$, donc:
\[ \psi(\alpha^{m}\varepsilon)=\theta^{-m}.\psi(\varepsilon), \tag{110.} \]
ou bien:
\[ \psi(\varepsilon)=\theta^{m}.\chi[\Phi^{2}(\alpha^{m}\varepsilon)], \]
d'où, en élevant les deux à la $n^{\text{ème}}$ puissance, on tire en vertu de la relation $\theta^{+mn}=1$:
\[ [\psi(\varepsilon)]^{n}=\{\chi[\Phi^{2}(\alpha^{m}\varepsilon)]\}^{n}. \tag{111.} \]

Cette formule donne, en faisant successivement $m=0$, $1$, $2$, $3\ldots\ldots n-1$, $n$ équations, qui, ajoutées membres à membres donneront la suivante:
\[ n[\psi(\varepsilon)]^{n}=\{\chi[\Phi^{2}(\varepsilon)]\}^{n}+\{\chi[\Phi^{2}(\alpha\varepsilon)]\}^{n}+\{\chi[\Phi^{2}(\alpha^{2}\varepsilon)]\}^{n}+\ldots+\{\chi[\Phi^{2}(\alpha^{n-1}\varepsilon)]\}^{n}; \tag{112.} \]
or, le second membre de cette équation est une fonction rationelle et symétrique des quantités $\Phi^{2}(\varepsilon)$, $\Phi^{2}(\alpha\varepsilon)\ldots\Phi^{2}(\alpha^{n-1}\varepsilon)$, c'est-à-dire des racines de l'équation (106.); donc $(\psi\varepsilon)^{n}$ peut être exprimé en fonction rationelle de $p_{0}$, $p_{1}\ldots p_{n-1}$, par conséquent en fonction rationelle de l'une quelconque de ces quantités. Soit $v$ la valeur de $\psi(\varepsilon)$, on aura
\[ \sqrt[n]{v}=\Phi^{2}(\varepsilon)+\theta.\Phi^{2}(\alpha\varepsilon)+\theta^{2}.\Phi^{2}(\alpha^{2}\varepsilon)+\ldots+\theta^{n-1}.\Phi^{2}(\alpha^{n-1}\varepsilon). \tag{113.} \]
Cela posé, soit $\theta=\cos\frac{2\pi}{n}+i\sin\frac{2\pi}{n}$. Les racines imaginaires de l'équation $\theta^{n}-1$ peuvent être représentées par
\[ \theta^{1},\ \theta^{2}\ \ldots\ldots\ \theta^{n-1}. \]
Donc en faisant successivement $\theta$ égal à chacune de ces racines et désignant les valeurs correspondantes de $v$ par $v_{1}$, $v_{2}\ldots v_{n-1}$, il viendra
\[ \begin{aligned}
\sqrt[n]{v_{1}}&=\Phi^{2}(\varepsilon)+\theta.\Phi^{2}(\alpha\varepsilon)+\ldots\ldots+\theta^{n-1}.\Phi^{2}(\alpha^{n-1}\varepsilon),\\
\sqrt[n]{v_{2}}&=\Phi^{2}(\varepsilon)+\theta^{2}.\Phi^{2}(\alpha\varepsilon)+\ldots\ldots+\theta^{2n-2}.\Phi^{2}(\alpha^{n-1}\varepsilon)\\
&\cdots\cdots\cdots\cdots\cdots\cdots\cdots\cdots\cdots\\
\sqrt[n]{v_{n-1}}&=\Phi^{2}(\varepsilon)+\theta^{n-1}.\Phi^{2}(\alpha\varepsilon)+\ldots\ldots+\theta^{(n-1)^{2}}.\Phi^{2}(\alpha^{n-1}\varepsilon),
\end{aligned} \]
En combinant ces équations avec la suivante:
\[ -p_{n-1}=\Phi^{2}(\varepsilon)+\Phi^{2}(\alpha\varepsilon)+\ldots\ldots+\Phi^{2}(\alpha^{n-1}\varepsilon), \]
on en tire aisément:
\[ \begin{aligned}
\Phi^{2}(\alpha^{m}\varepsilon)=+\frac{1}{n}&\left\{-p_{n-1}+\theta^{-m}.\sqrt[n]{v_{1}}+\theta^{-2m}.\sqrt[n]{v_{2}}+\theta^{-3m}.\sqrt[n]{v_{3}}+\ldots\right.\\
&\left.\ldots+\theta^{-(n-1)m}\sqrt[n]{v_{n-1}}\right\},
\end{aligned} \tag{114.} \]
et pour $m=0$:
\[ \Phi^{2}(\varepsilon)=\frac{1}{n}\left\{-p_{n-1}+\sqrt[n]{v_{1}}+\sqrt[n]{v_{2}}+\ldots+\sqrt[n]{v_{n-1}}\right\}. \tag{115.} \]

\origpage{145}

\subsection*{22.}

Toutes les racines de l'équation (106.) sont contenues dans la formule (115.), mais puisque leur nombre n'est que $n$, il reste encore à donner à $\Phi^{2}(\varepsilon)$ une forme, qui ne contienne pas des racines étrangères à la question. Or cela se fait aisément, comme suit:

Soit
\[ s_{k}=\frac{\sqrt[n]{v_{k}}}{(\sqrt[n]{v_{1}})^{k}}. \]
En posant ici $\varepsilon$ au lieu de $\alpha^{m}\varepsilon$, $\sqrt[n]{v_{k}}$ se changera en $\theta^{-km}.\sqrt[n]{v_{k}}$, et $v_{1}$ en $\theta^{-m}.v_{1}$, donc $s_{k}$ se changera en:
\[ \frac{\theta^{-km}.\sqrt[n]{v_{k}}}{(\theta^{-m}\sqrt[n]{v_{1}})^{k}}=\frac{\sqrt[n]{v_{k}}}{(\sqrt[n]{v_{1}})^{k}}. \]

La fonction $s_{k}$, comme on voit, ne change pas de valeur, en mettant $\alpha^{m}\varepsilon$ au lieu de $\varepsilon$. Or $s_{k}$ est une fonction rationelle de $\Phi^{2}(\varepsilon)$. Donc, en désignant $s_{k}$ par $\lambda[\Phi^{2}(\varepsilon)]$, on aura
\[ s_{k}=\lambda[\Phi^{2}(\alpha^{m}\varepsilon)], \]
quel que soit le nombre entier $m$. De là on tirera de la même manière, et comme nous avons trouvé $(\psi\varepsilon)^{n}$, la valeur de $s_{k}$, en fonction rationelle de l'une des quantités $p_{0}$, $p_{1}$, $....p_{n-1}$. Connoissant $s_{k}$, on a:
\[ \sqrt[n]{v_{k}}=s_{k}.(\sqrt[n]{v_{1}})^{k}. \]
Donc, en mettant $v$ au lieu de $v_{1}$, l'expression de $\Phi^{2}(\alpha^{m}\varepsilon)$ deviendra:
\[ \Phi^{2}(\alpha^{m}\varepsilon)=\frac{1}{n}\left\{-p_{n-1}+\theta^{-m}.v^{\frac{1}{n}}+s_{2}\theta^{-2m}.v^{\frac{2}{n}}+....+s_{n-1}\theta^{-(n-1)m}.v^{\frac{n-1}{n}}\right\}, \tag{116.} \]
pour $m=0$:
\[ \Phi^{2}(\varepsilon)=\frac{1}{n}\left\{-p_{n-1}+v^{\frac{1}{n}}+s_{2}.v^{\frac{2}{n}}+s_{3}.v^{\frac{3}{n}}+....+s_{n-1}.v^{\frac{n-1}{n}}\right\}. \tag{117.} \]
Cette valeur n'a que $n$ valeurs différentes, qui répondent aux $n$ valeurs de $v^{\frac{1}{n}}$. Donc en dernier lieu la résolution de l'équation $P_{2n+1}=0$ est réduite à celle d'une seule équation du degré $2n+2$; mais cette équation ne paraît pas \emph{en général} être résoluble algébriquement. Néanmoins on peut la résoudre complètement dans plusieurs cas particuliers, p.\ ex., lorsque $e=c$, $e=c\sqrt{3}$, $e=c(2\pm\sqrt{3})$ etc. Dans le cours de ce mémoire je m'occuperai de ces cas, dont le premier surtout est remarquable, tant par la simplicité de la solution, que par sa belle application dans la géométrie.

\origpage{146}

En effet entre autres je suis parvenu à ce théorème:

„On peut diviser la circonférence entière de la lemniscate par la \emph{règle et le compas seuls}, en $m$ parties égales, si $m$ est de la forme $2^{n}$ ou \uncertain{$2^{n}+1$}, le dernier nombre étant en même temps premier; ou bien si $m$ est un produit de plusieurs nombres de ces deux formes.”

Ce théorème est, comme on voit, précisément le même que celui de M. \emph{Gauss}, relativement au cercle.

\section*{§. VI.}

\emph{Expressions diverses des fonctions} $\Phi(n\beta)$, $f(n\beta)$, $F(n\beta)$.

\subsection*{23.}

En faisant usage des formules connues, qui donnent les valeurs des coefficiens d'une équation algébrique en fonction des racines, on peut tirer plusieurs expressions des fonctions $\varphi(n\beta)$, $f(n\beta)$, $F(n\beta)$ des formules du paragraphe précédent.

Je vais considérer les plus remarquables.

Pour abréger les formules, je me servirai des notations suivantes: Je désignerai:

1) Par $\overset{k'}{\underset{k}{\Sigma}}_{m}\psi(m)$ la somme, et par $\overset{k'}{\underset{k}{\Pi}}_{m}\psi(m)$ le produit de toutes les quantités de la forme $\psi(m)$, qu'on obtiendra, en donnant à $m$ toutes les valeurs entières, depuis $k$ jusqu'à $k'$, les limites $k$ et $k'$ y compris.

2) Par $\overset{k'}{\underset{k}{\Sigma}}_{m}\overset{\nu'}{\underset{\nu}{\Sigma}}_{\mu}\psi(m,\mu)$ la somme, et par $\overset{k'}{\underset{k}{\Pi}}_{m}\overset{\nu'}{\underset{\nu}{\Pi}}_{\mu}\psi(m,\mu)$ le produit de toutes les quantités de la forme $\psi(m,\mu)$, qu'on obtiendra, en donnant à $m$ toutes les valeurs entières de $k$ à $k'$, et à $\mu$ les valeurs entières de $\nu$ à $\nu'$, en y comprenant toujours les limites.

D'après cela il est clair, qu'on aura:
\[ \overset{k'}{\underset{k}{\Sigma}}_{m}\psi(m)=\psi(k)+\psi(k+1)+.....\psi(k'), \tag{119.} \]
\ednote{No equation number 118. is printed; the numbering passes from 117. to 119.}
\[ \overset{k'}{\underset{k}{\Pi}}_{m}\psi(m)=\psi(k).\psi(k+1)+.....\psi(k'), \tag{120.} \]
\ednote{A plus sign is printed before the row of dots in the product; a multiplication dot is expected.}
\[ \overset{k'}{\underset{k}{\Sigma}}_{m}\overset{\nu'}{\underset{\nu}{\Sigma}}_{\mu}\psi(m,\mu)=\overset{\nu'}{\underset{\nu}{\Sigma}}_{\mu}\psi(k,\mu)+\overset{\nu'}{\underset{\nu}{\Sigma}}_{\mu}\psi(k+1,\mu)+.....+\overset{\nu'}{\underset{\nu}{\Sigma}}_{\mu}\psi(k',\mu), \tag{121.} \]
\[ \overset{k'}{\underset{k}{\Pi}}_{m}\overset{\nu'}{\underset{\nu}{\Pi}}_{\mu}\psi(m,\mu)=\overset{\nu'}{\underset{\nu}{\Pi}}_{\mu}\psi(k,\mu).\overset{\nu'}{\underset{\nu}{\Pi}}_{\mu}\psi(k+1,\mu).........\overset{\nu'}{\underset{\nu}{\Pi}}_{\mu}\psi(k',\mu). \tag{122.} \]

\origpage{147}

Cela posé, considérons les équations
\[ \left\{\begin{aligned}
\Phi(2n+1)\beta&=\frac{P_{2n+1}}{Q_{2n+1}},\\
f(2n+1)\beta&=\frac{P'_{2n+1}}{Q_{2n+1}},\\
F(2n+1)\beta&=\frac{P''_{2n+1}}{Q_{2n+1}}.
\end{aligned}\right. \tag{123.} \]

Nous avons vu, que $P_{2n+1}$ est une fonction rationelle de $x$ du degré $(2n+1)^{2}$ et de la forme $x.\psi(x^{2})$. De même $P'_{2n+1}$ et $F''_{2n+1}$ sont des fonctions de cette même forme, la première de $y$ et la seconde de $z$.\ednote{Printed $F''_{2n+1}$ where $P''_{2n+1}$ is expected.} Enfin $Q_{2n+1}$ est une fonction qui, exprimée indistinctement en $x$, $y$ ou $z$, sera du degré $(2n+1)^{2}-1$, et contiendra seulement des puissances paires.
Donc on aura
\[ \begin{aligned}
P_{2n+1}&=A.x^{(2n+1)^{2}}+......+B.x;&\quad P'_{2n+1}&=A'.y^{(2n+1)^{2}}+......+B'.y;\\
P''_{2n+1}&=A''.z^{(2n+1)^{2}}+......+B''.z;&\quad Q_{2n+1}&=C.x^{(2n+1)^{2}-1}+....+D;\\
Q_{2n+1}&=C'.y^{(2n+1)^{2}-1}+....+D';&\quad Q_{2n+1}&=C''.z^{(2n+1)^{2}-1}+....+D''.
\end{aligned} \]

En substituant ces valeurs dans (123.), il viendra
\[ \begin{aligned}
&\{A.x^{(2n+1)^{2}}+.....+B.x\}-\varphi(2n+1)\beta.\{C.x^{(2n+1)^{2}-1}+....+D\},\\
&\{A'.y^{(2n+1)^{2}}+.....+B'.y\}-f(2n+1)\beta.\{C'.y^{(2n+1)^{2}-1}+....+D'\},\\
&\{A''.z^{(2n+1)^{2}}+.....+B''.z\}-F(2n+1)\beta.\{C''.z^{(2n+1)^{2}-1}+....+D''\}.
\end{aligned} \]

Dans la première de ces équations $A$ est le coefficient du premier terme, $-\varphi(2n+1)\beta.C$ celui du second et $-\varphi(2n+1)\beta.D$ le dernier terme. Donc $\frac{C}{A}\varphi(2n+1)\beta$ est égal à la somme et $\pm\frac{D}{A}\varphi(2n+1)\beta$ égal au produit des racines de l'équation dont il s'agit, équation, qui est la même, que celle-ci:
\[ \Phi(2n+1)\beta=\frac{P_{2n+1}}{Q_{2n+1}}. \tag{124.} \]

Donc en remarquant que $A$, $C$ et $D$ (et en général tous les coefficiens) sont indépendans de $\beta$, on voit que $\Phi(2n+1)\beta$ est (d'un coefficient constant près) égal à la somme et au produit de toutes les racines de l'équation (124.).

De la même manière on voit que $f(2n+1)\beta$ et $F(2n+1)\beta$ sont respectivement égaux au produit ou à la somme des racines des équations
\[ f(2n+1)\beta=\frac{P'_{2n+1}}{Q_{2n+1}},\ F(2n+1)\beta=\frac{P''_{2n+1}}{Q_{2n+1}}, \]

\origpage{148}
en ayant attention de multiplier le résultat par un coefficient constant, choisi convenablement.

Maintenant les racines des équations (123.) d'après le No.\ 11. sont respectivement:
\[ \begin{aligned}
x&=(-1)^{m+\mu}.\Phi\left(\beta+\frac{m}{2n+1}\omega+\frac{\mu}{2n+1}\varpi i\right),\\
y&=(-1)^{m}.f\left(\beta+\frac{m}{2n+1}\omega+\frac{\mu}{2n+1}\varpi i\right),\\
z&=(-1)^{\mu}.F\left(\beta+\frac{m}{2n+1}\omega+\frac{\mu}{2n+1}\varpi i\right),
\end{aligned} \]
où les limites de $m$ et $\mu$ sont $-n$ et $+n$.

Donc en vertu de ce qu'on vient de voir, et en faisant usage des notations adoptées, on aura les formules suivantes:
\[ \left\{\begin{aligned}
\Phi(2n+1)\beta&=A.\overset{+n}{\underset{-n}{\Sigma}}_{m}\overset{+n}{\underset{-n}{\Sigma}}_{\mu}(-1)^{m+\mu}.\Phi\left(\beta+\frac{m.\omega+\mu\varpi i}{2n+1}\right),\\
f(2n+1)\beta&=A'.\overset{+n}{\underset{-n}{\Sigma}}_{m}\overset{+n}{\underset{-n}{\Sigma}}_{\mu}(-1)^{m}.f\left(\beta+\frac{m\omega+\mu\varpi i}{2n+1}\right),\\
F(2n+1)\beta&=A''.\overset{+n}{\underset{-n}{\Sigma}}_{m}\overset{+n}{\underset{-n}{\Sigma}}_{\mu}(-1)^{\mu}.F\left(\beta+\frac{m\omega+\mu\varpi i}{2n+1}\right);\\
\Phi(2n+1)\beta&=B.\overset{+n}{\underset{-n}{\Pi}}_{m}\overset{+n}{\underset{-n}{\Pi}}_{\mu}.\Phi\left(\beta+\frac{m\omega+\mu\varpi i}{2n+1}\right),\\
f(2n+1)\beta&=B'.\overset{+n}{\underset{-n}{\Pi}}_{m}\overset{+n}{\underset{-n}{\Pi}}_{\mu}.f\left(\beta+\frac{m\omega+\mu\varpi i}{2n+1}\right),\\
F(2n+1)\beta&=B''.\overset{+n}{\underset{-n}{\Pi}}_{m}\overset{+n}{\underset{-n}{\Pi}}_{\mu}.F\left(\beta+\frac{m\omega+\mu\varpi i}{2n+1}\right).
\end{aligned}\right. \tag{125.} \]

Pour déterminer les quantités constantes $A$, $A'$, $A''$, $B$, $B'$, $B''$, il faudra donner à $\beta$ une valeur particulière. Ainsi en faisant dans les trois premières formules $\beta=\frac{\omega}{2}+\frac{\varpi}{2}i$, après avoir divisé les deux membres par $\Phi\beta$, il viendra, en remarquant que $\Phi\left(\frac{\omega}{2}+\frac{\varpi}{2}i\right)=\frac{1}{0}$:
\[ \left.\begin{aligned}
A&=\frac{\varphi(2n+1)\beta}{\varphi\beta}\\
A'&=\frac{f(2n+1)\beta}{f\beta}\\
A''&=\frac{F(2n+1)\beta}{f\beta}
\end{aligned}\right\}\ \text{pour }\beta=\frac{\omega}{2}+\frac{\varpi}{2}i. \]
\ednote{The denominator of the third fraction is printed $f\beta$; $F\beta$ is apparently expected.}

\origpage{149}

Soit $\beta=\frac{\omega}{2}+\frac{\varpi}{2}i+\alpha$, on a:
\[ \left.\begin{aligned}
A&=\frac{\varphi\left((2n+1)\alpha+n\omega+n\varpi i+\frac{\omega}{2}+\frac{\varpi}{2}i\right)}{\varphi\left(\alpha+\frac{\omega}{2}+\frac{\varpi}{2}i\right)}\\
&=\frac{\varphi\left((2n+1)\alpha+\frac{\omega}{2}+\frac{\varpi}{2}i\right)}{\varphi\left(\alpha+\frac{\omega}{2}+\frac{\varpi}{2}i\right)}=\frac{\varphi\alpha}{\varphi(2n+1)\alpha},\\
A'&=\frac{f\left((2n+1)\alpha+n\omega+n\varpi i+\frac{\omega}{2}+\frac{\varpi}{2}i\right)}{f\left(\alpha+\frac{\omega}{2}+\frac{\varpi}{2}i\right)}\\
&=(-1)^{n}.\frac{f\left((2n+1)\alpha+\frac{\omega}{2}+\frac{\varpi}{2}i\right)}{f\left(\alpha+\frac{\omega}{2}+\frac{\varpi}{2}i\right)}=(-1)^{n}.\frac{f\left(\alpha+\frac{\omega}{2}\right)}{f\left((2n+1)\alpha+\frac{\omega}{2}\right)},\\
A''&=\frac{F\left((2n+1)\alpha+n\omega+n\varpi i+\frac{\omega}{2}+\frac{\varpi}{2}i\right)}{F\left(\alpha+\frac{\omega}{2}+\frac{\varpi}{2}i\right)}\\
&=(-1)^{n}.\frac{F\left((2n+1)\alpha+\frac{\omega}{2}+\frac{\varpi}{2}i\right)}{F\left(\alpha+\frac{\omega}{2}+\frac{\varpi}{2}i\right)}=(-1)^{n}.\frac{F\left(\alpha+\frac{\varpi}{2}i\right)}{F\left((2n+1)\alpha+\frac{\varpi}{2}i\right)}.
\end{aligned}\right\}\ \text{pour }\alpha=0. \]

Ces expressions de $A$, $A'$, $A''$ deviendront de la forme $\frac{0}{0}$, en faisant $\alpha=0$, donc on trouvera d'après les règles connues:
\[ A=\frac{1}{2n+1},\quad A'=A''=\frac{(-1)^{n}}{2n+1}. \]
D'après cela les trois premières formules deviendront:
\[ \left\{\begin{aligned}
\Phi(2n+1)\beta&=\frac{1}{2n+1}.\overset{+n}{\underset{-n}{\Sigma}}_{m}\overset{+n}{\underset{-n}{\Sigma}}_{\mu}(-1)^{m+\mu}.\varphi\left(\beta+\frac{m\omega+\mu\varpi i}{2n+1}\right),\\
f(2n+1)\beta&=\frac{(-1)^{n}}{2n+1}.\overset{+n}{\underset{-n}{\Sigma}}_{m}\overset{+n}{\underset{-n}{\Sigma}}_{\mu}(-1)^{m}.f\left(\beta+\frac{m\omega+\mu\varpi i}{2n+1}\right),\\
F(2n+1)\beta&=\frac{(-1)^{n}}{2n+1}.\overset{+n}{\underset{-n}{\Sigma}}_{m}\overset{+n}{\underset{-n}{\Sigma}}_{\mu}(-1)^{\mu}.F\left(\beta+\frac{m\omega+\mu\varpi i}{2n+1}\right).
\end{aligned}\right. \tag{126.} \]

\origpage{150}

Pour avoir la valeur des constantes $B$, $B'$, $B''$, je remarque, qu'on aura:
\[ \begin{aligned}
&\overset{+n}{\underset{-n}{\Pi}}_{m}\overset{+n}{\underset{-n}{\Pi}}_{\mu}\psi(m,\mu)=\psi(0,0).\overset{n}{\underset{1}{\Pi}}_{m}\psi(m,0).\overset{n}{\underset{1}{\Pi}}_{\mu}\psi(0,\mu)\\
&\quad\times\overset{n}{\underset{1}{\Pi}}_{m}\overset{n}{\underset{1}{\Pi}}_{\mu}\psi(m,\mu).\psi(-m,-\mu).\overset{n}{\underset{1}{\Pi}}_{m}\overset{n}{\underset{1}{\Pi}}_{\mu}\psi(m,-\mu).\psi(-m,\mu).
\end{aligned} \tag{127.} \]

En appliquant cette transformation aux formules (126.), divisant la première par $\Phi\beta$, la seconde par $f\beta$ et la troisième par $F\beta$, faisant ensuite dans la première $\beta=0$, dans la seconde $\beta=\frac{\omega}{2}$ et dans la troisième $\beta=\frac{\varpi}{2}i$, et remarquant que $\frac{\varphi(2n+1)\beta}{\varphi\beta}=2n+1$, pour $\beta=0$, que $\frac{f(2n+1)\beta}{f\beta}=2n+1$, pour $\beta=\frac{\omega}{2}$, et que $\frac{F(2n+1)\beta}{F\beta}=2n+1$, pour $\beta=\frac{\varpi}{2}i$: on trouvera:
\[ \left\{\begin{aligned}
(2n+1)&=B.\overset{n}{\underset{1}{\Pi}}_{m}\varphi^{2}\left(\frac{m\omega}{2n+1}\right).\overset{n}{\underset{1}{\Pi}}_{\mu}\varphi^{2}\left(\frac{\mu\varpi i}{2n+1}\right)\\
&\quad\times\overset{n}{\underset{1}{\Pi}}_{m}\overset{n}{\underset{1}{\Pi}}_{\mu}\varphi^{2}\left(\frac{m\omega+\mu\varpi i}{2n+1}\right).\varphi^{2}\left(\frac{m\omega-\mu\varpi i}{2n+1}\right),\\
(2n+1)&=B'.\overset{n}{\underset{1}{\Pi}}_{m}f^{2}\left(\frac{\omega}{2}+\frac{m\omega}{2n+1}\right).\overset{n}{\underset{1}{\Pi}}_{\mu}f^{2}\left(\frac{\omega}{2}+\frac{\mu\varpi i}{2n+1}\right)\\
&\quad\times\overset{n}{\underset{1}{\Pi}}_{m}\overset{n}{\underset{1}{\Pi}}_{\mu}f^{2}\left(\frac{\omega}{2}+\frac{m\omega+\mu\varpi i}{2n+1}\right).f^{2}\left(\frac{\omega}{2}+\frac{m\omega-\mu\varpi i}{2n+1}\right),\\
(2n+1)&=B''.\overset{n}{\underset{1}{\Pi}}_{m}F^{2}\left(\frac{\varpi}{2}i+\frac{m\omega}{2n+1}\right).\overset{n}{\underset{1}{\Pi}}_{\mu}F^{2}\left(\frac{\varpi}{2}i+\frac{\mu\varpi i}{2n+1}\right)\\
&\quad\times\overset{n}{\underset{1}{\Pi}}_{m}\overset{n}{\underset{1}{\Pi}}_{\mu}F^{2}\left(\frac{\varpi}{2}i+\frac{m\omega+\mu\varpi i}{2n+1}\right).F^{2}\left(\frac{\varpi}{2}i+\frac{m\omega-\mu\varpi i}{2n+1}\right).
\end{aligned}\right. \tag{128.} \]

En tirant de ces équations les valeurs de $B$, $B'$, $B''$, et les substituant ensuite dans les formules transformées, il viendra:

\origpage{151}
\[ \left\{\begin{aligned}
&\Phi(2n+1)\beta=\\
&(2n+1)\Phi\beta.\overset{n}{\underset{1}{\Pi}}_{m}\frac{\varphi\left(\beta+\frac{m\omega}{2n+1}\right).\varphi\left(\beta-\frac{m\omega}{2n+1}\right)}{\varphi^{2}\left(\frac{m\omega}{2n+1}\right)}.\overset{n}{\underset{1}{\Pi}}_{\mu}\frac{\varphi\left(\beta+\frac{\mu\varpi i}{2n+1}\right).\varphi\left(\beta-\frac{\mu\varpi i}{2n+1}\right)}{\varphi^{2}\left(\frac{\mu\varpi i}{2n+1}\right)}\\
&\quad\times\overset{n}{\underset{1}{\Pi}}_{m}\overset{n}{\underset{1}{\Pi}}_{\mu}\frac{\varphi\left(\beta+\frac{m\omega+\mu\varpi i}{2n+1}\right).\varphi\left(\beta-\frac{m\omega+\mu\varpi i}{2n+1}\right)}{\varphi^{2}\left(\frac{m\omega+\mu\varpi i}{2n+1}\right)}.\frac{\varphi\left(\beta+\frac{m\omega-\mu\varpi i}{2n+1}\right).\varphi\left(\beta-\frac{m\omega-\mu\varpi i}{2n+1}\right)}{\varphi^{2}\left(\frac{m\omega-\mu\varpi i}{2n+1}\right)},\\
&f(2n+1)\beta=\\
&(2n+1)f\beta.\overset{n}{\underset{1}{\Pi}}_{m}\frac{f\left(\beta+\frac{m\omega}{2n+1}\right).f\left(\beta-\frac{m\omega}{2n+1}\right)}{f^{2}\left(\frac{\omega}{2}+\frac{m\omega}{2n+1}\right)}.\overset{n}{\underset{1}{\Pi}}_{\mu}\frac{f\left(\beta+\frac{\mu\varpi i}{2n+1}\right).f\left(\beta-\frac{\mu\varpi i}{2n+1}\right)}{f^{2}\left(\frac{\omega}{2}+\frac{\mu\varpi i}{2n+1}\right)}\\
&\quad\times\overset{n}{\underset{1}{\Pi}}_{m}\overset{n}{\underset{1}{\Pi}}_{\mu}\frac{f\left(\beta+\frac{m\omega+\mu\varpi i}{2n+1}\right).f\left(\beta-\frac{m\omega+\mu\varpi i}{2n+1}\right)}{f^{2}\left(\frac{\omega}{2}+\frac{m\omega+\mu\varpi i}{2n+1}\right)}.\frac{f\left(\beta+\frac{m\omega-\mu\varpi i}{2n+1}\right).f\left(\beta-\frac{m\omega-\mu\varpi i}{2n+1}\right)}{f^{2}\left(\frac{\omega}{2}+\frac{m\omega-\mu\varpi i}{2n+1}\right)},\\
&F(2n+1)\beta=\\
&(2n+1)F\beta.\overset{n}{\underset{1}{\Pi}}_{m}\frac{F\left(\beta+\frac{m\omega}{2n+1}\right).F\left(\beta-\frac{m\omega}{2n+1}\right)}{F^{2}\left(\frac{\varpi}{2}i+\frac{m\omega}{2n+1}\right)}.\overset{n}{\underset{1}{\Pi}}_{\mu}\frac{F\left(\beta+\frac{\mu\varpi i}{2n+1}\right).F\left(\beta-\frac{\mu\varpi i}{2n+1}\right)}{F^{2}\left(\frac{\varpi}{2}i+\frac{\mu\varpi i}{2n+1}\right)}\\
&\quad\times\overset{n}{\underset{1}{\Pi}}_{m}\overset{n}{\underset{1}{\Pi}}_{\mu}\frac{F\left(\beta+\frac{m\omega+\mu\varpi i}{2n+1}\right).F\left(\beta-\frac{m\omega+\mu\varpi i}{2n+1}\right)}{F^{2}\left(\frac{\varpi}{2}i+\frac{m\omega+\mu\varpi i}{2n+1}\right)}.\frac{F\left(\beta+\frac{m\omega-\mu\varpi i}{2n+1}\right).F\left(\beta-\frac{m\omega-\mu\varpi i}{2n+1}\right)}{F^{2}\left(\frac{\varpi}{2}i+\frac{m\omega-\mu\varpi i}{2n+1}\right)}.
\end{aligned}\right. \tag{129.} \]

On peut donner à ces formules des formes plus simples, en faisant usage des formules suivantes:
\[ \frac{\varphi(\beta+\alpha).\varphi(\beta-\alpha)}{\varphi^{2}\alpha}=-\frac{1-\frac{\varphi^{2}\beta}{\varphi^{2}\alpha}}{1-\frac{\varphi^{2}\beta}{\varphi^{2}\left(\alpha+\frac{\omega}{2}+\frac{\varpi}{2}i\right)}}, \]
\[ \frac{f(\beta+\alpha).f(\beta-\alpha)}{f^{2}\left(\frac{\omega}{2}+\alpha\right)}=-\frac{1-\frac{f^{2}\beta}{f^{2}\left(\frac{\omega}{2}+\alpha\right)}}{1-\frac{f^{2}\beta}{f^{2}\left(\alpha+\frac{\omega}{2}+\frac{\varpi}{2}i\right)}}, \]
\[ \frac{F^{2}(\beta+\alpha).F(\beta-\alpha)}{F^{2}\left(\frac{\varpi}{2}i+\alpha\right)}=-\frac{1-\frac{F^{2}\beta}{F^{2}\left(\frac{\varpi}{2}i+\alpha\right)}}{1-\frac{F^{2}\beta}{F^{2}\left(\frac{\omega}{2}+\frac{\varpi}{2}i+\alpha\right)}}, \]
\ednote{The numerator of the third left-hand fraction is printed $F^{2}(\beta+\alpha).F(\beta-\alpha)$; $F(\beta+\alpha).F(\beta-\alpha)$ is apparently expected.}

\origpage{152}
qu'on verifiera aisément au moyen des formules (13.), (16.), (18.).

En vertu de ces formules il est clair, qu'on peut mettre les équations (129.) sous la forme:
\[ \left\{\begin{aligned}
&\Phi(2n+1)\beta=\\
&(2n+1)\Phi\beta.\overset{n}{\underset{1}{\Pi}}_{m}\frac{1-\frac{\varphi^{2}\beta}{\varphi^{2}\left(\frac{m\omega}{2n+1}\right)}}{1-\frac{\varphi^{2}\beta}{\varphi^{2}\left(\frac{\omega}{2}+\frac{\varpi}{2}i+\frac{m\omega}{2n+1}\right)}}.\overset{n}{\underset{1}{\Pi}}_{\mu}\frac{1-\frac{\varphi^{2}\beta}{\varphi^{2}\left(\frac{\mu\varpi i}{2n+1}\right)}}{1-\frac{\varphi^{2}\beta}{\varphi^{2}\left(\frac{\omega}{2}+\frac{\varpi}{2}i+\frac{\mu\varpi i}{2n+1}\right)}}\\
&\quad\times\overset{n}{\underset{1}{\Pi}}_{m}\overset{n}{\underset{1}{\Pi}}_{\mu}\frac{1-\frac{\varphi^{2}\beta}{\varphi^{2}\left(\frac{m\omega+\mu\varpi i}{2n+1}\right)}}{1-\frac{\varphi^{2}\beta}{\varphi^{2}\left(\frac{\omega}{2}+\frac{\varpi}{2}i+\frac{m\omega+\mu\varpi i}{2n+1}\right)}}.\frac{1-\frac{\varphi^{2}\beta}{\varphi^{2}\left(\frac{m\omega-\mu\varpi i}{2n+1}\right)}}{1-\frac{\varphi^{2}\beta}{\varphi^{2}\left(\frac{\omega}{2}+\frac{\varpi}{2}i+\frac{m\omega-\mu\varpi i}{2n+1}\right)}},\\
&f(2n+1)\beta=\\
&(2n+1)f\beta.\overset{n}{\underset{1}{\Pi}}_{m}\frac{1-\frac{f^{2}\beta}{f^{2}\left(\frac{m\omega}{2n+1}+\frac{\omega}{2}\right)}}{1-\frac{f^{2}\beta}{f^{2}\left(\frac{\omega}{2}+\frac{\varpi}{2}i+\frac{m\omega}{2n+1}\right)}}.\overset{n}{\underset{1}{\Pi}}_{\mu}\frac{1-\frac{f^{2}\beta}{f^{2}\left(\frac{\omega}{2}+\frac{\mu\varpi i}{2n+1}\right)}}{1-\frac{f^{2}\beta}{f^{2}\left(\frac{\omega}{2}+\frac{\varpi}{2}i+\frac{\mu\varpi i}{2n+1}\right)}}\\
&\quad\times\overset{n}{\underset{1}{\Pi}}_{m}\overset{n}{\underset{1}{\Pi}}_{\mu}\frac{1-\frac{f^{2}\beta}{f^{2}\left(\frac{\omega}{2}+\frac{m\omega+\mu\varpi i}{2n+1}\right)}}{1-\frac{f^{2}\beta}{f^{2}\left(\frac{\omega}{2}+\frac{\varpi}{2}i+\frac{m\omega+\mu\varpi i}{2n+1}\right)}}.\frac{1-\frac{f^{2}\beta}{f^{2}\left(\frac{\omega}{2}+\frac{m\omega-\mu\varpi i}{2n+1}\right)}}{1-\frac{f^{2}\beta}{f^{2}\left(\frac{\omega}{2}+\frac{\varpi}{2}i+\frac{m\omega-\mu\varpi i}{2n+1}\right)}},\\
&F(2n+1)\beta=\\
&(2n+1)F\beta.\overset{n}{\underset{1}{\Pi}}_{m}\frac{1-\frac{F^{2}\beta}{F^{2}\left(\frac{\varpi}{2}i+\frac{m\omega}{2n+1}\right)}}{1-\frac{F^{2}\beta}{F^{2}\left(\frac{\omega}{2}+\frac{\varpi}{2}i+\frac{m\omega}{2n+1}\right)}}.\overset{n}{\underset{1}{\Pi}}_{\mu}\frac{1-\frac{F^{2}\beta}{F^{2}\left(\frac{\varpi}{2}i+\frac{\mu\varpi i}{2n+1}\right)}}{1-\frac{F^{2}\beta}{F^{2}\left(\frac{\omega}{2}+\frac{\varpi}{2}i+\frac{\mu\varpi i}{2n+1}\right)}}\\
&\quad\times\overset{n}{\underset{1}{\Pi}}_{m}\overset{n}{\underset{1}{\Pi}}_{\mu}\frac{1-\frac{F^{2}\beta}{F^{2}\left(\frac{\varpi}{2}i+\frac{m\omega+\mu\varpi i}{2n+1}\right)}}{1-\frac{F^{2}\beta}{F^{2}\left(\frac{\omega}{2}+\frac{\varpi}{2}i+\frac{m\omega+\mu\varpi i}{2n+1}\right)}}.\frac{1-\frac{F^{2}\beta}{F^{2}\left(\frac{\varpi}{2}i+\frac{m\omega-\mu\varpi i}{2n+1}\right)}}{1-\frac{F^{2}\beta}{F^{2}\left(\frac{\omega}{2}+\frac{\varpi}{2}i+\frac{m\omega-\mu\varpi i}{2n+1}\right)}}.
\end{aligned}\right. \tag{130.} \]

\origpage{153}

Ces formules donnent, comme on voit, les valeurs de $\Phi(2n+1)\beta$, $f(2n+1)\beta$ et $F(2n+1)\beta$, exprimées respectivement en fonctions rationelles de $\Phi\beta$, $f\beta$ et $F\beta$, sous la forme de produits.

Nous donnerons encore les valeurs de $f(2n+1)\beta$, $F(2n+1)\beta$ sous une autre forme, qui sera utile dans la suite.

On a $f^{2}\beta=1-c^{2}\varphi^{2}\beta$, donc
\[ 1-\frac{f^{2}\beta}{f^{2}\alpha}=\frac{c^{2}(\varphi^{2}\beta-\varphi^{2}\alpha)}{f^{2}\alpha}=\frac{c^{2}}{f^{2}\alpha}\Phi^{2}\beta-\frac{c^{2}\varphi^{2}\alpha}{f^{2}\alpha} \]
et
\[ 1-\frac{f^{2}\beta}{f^{2}\left(\frac{\varpi}{2}i+\alpha\right)}=\frac{c^{2}\left[\varphi^{2}\beta-\varphi^{2}\left(\frac{\varpi}{2}i+\alpha\right)\right]}{f^{2}\left(\frac{\varpi}{2}i+\alpha\right)}, \]
or en vertu de (18.) on a:
\[ f^{2}\left(\frac{\varpi}{2}i+\alpha\right)=\frac{e^{2}+c^{2}}{e^{2}}.\frac{1}{f^{2}\alpha}; \]
donc:
\[ \frac{1-\frac{f^{2}\beta}{f^{2}\alpha}}{1-\frac{f^{2}\beta}{f^{2}\left(\frac{\varpi}{2}i+\alpha\right)}}=\frac{1}{f^{4}\alpha}.\frac{e^{2}+c^{2}}{e^{2}}.\frac{1-\frac{\varphi^{2}\beta}{\varphi^{2}\alpha}}{1-\frac{\varphi^{2}\beta}{\varphi^{2}\left(\frac{\varpi}{2}i+\alpha\right)}}. \]
On trouvera de même:
\[ \frac{1-\frac{F^{2}\beta}{F^{2}\alpha}}{1-\frac{F^{2}\beta}{F^{2}\left(\frac{\omega}{2}+\alpha\right)}}=\frac{1}{F^{4}\alpha}.\frac{e^{2}+c^{2}}{e^{2}}.\frac{1-\frac{\varphi^{2}\beta}{\varphi^{2}\alpha}}{1-\frac{\varphi^{2}\beta}{\varphi^{2}\left(\frac{\omega}{2}+\alpha\right)}}. \]

En vertu de ces formules, et en faisant $\beta=0$, pour déterminer le facteur constant, il est clair, qu'on peut écrire les expressions de $f(2n+1)\beta$, $F(2n+1)\beta$, comme suit:

\origpage{154}
\[ \left\{\begin{aligned}
&f(2n+1)\beta=f\beta.\overset{n}{\underset{1}{\Pi}}_{m}\frac{1-\frac{\varphi^{2}\beta}{\varphi^{2}\left(\frac{\omega}{2}+\frac{m\omega}{2n+1}\right)}}{1-\frac{\varphi^{2}\beta}{\varphi^{2}\left(\frac{\omega}{2}+\frac{\varpi}{2}i+\frac{m\omega}{2n+1}\right)}}.\overset{n}{\underset{1}{\Pi}}_{\mu}\frac{1-\frac{\varphi^{2}\beta}{\varphi^{2}\left(\frac{\omega}{2}i+\frac{\mu\varpi i}{2n+1}\right)}}{1-\frac{\varphi^{2}\beta}{\varphi^{2}\left(\frac{\omega}{2}+\frac{\varpi}{2}i+\frac{\mu\varpi i}{2n+1}\right)}}\\
&\quad\times\overset{n}{\underset{1}{\Pi}}_{m}\overset{n}{\underset{1}{\Pi}}_{\mu}\frac{1-\frac{\varphi^{2}\beta}{\varphi^{2}\left(\frac{\omega}{2}+\frac{m\omega+\mu\varpi i}{2n+1}\right)}}{1-\frac{\varphi^{2}\beta}{\varphi^{2}\left(\frac{\omega}{2}+\frac{\varpi}{2}i+\frac{m\omega+\mu\varpi i}{2n+1}\right)}}.\frac{1-\frac{\varphi^{2}\beta}{\varphi^{2}\left(\frac{\omega}{2}+\frac{m\omega-\mu\varpi i}{2n+1}\right)}}{1-\frac{\varphi^{2}\beta}{\varphi^{2}\left(\frac{\omega}{2}+\frac{\varpi}{2}i+\frac{m\omega-\mu\varpi i}{2n+1}\right)}},\\
&F(2n+1)\beta=F\beta.\overset{n}{\underset{1}{\Pi}}_{m}\frac{1-\frac{\varphi^{2}\beta}{\varphi^{2}\left(\frac{\varpi}{2}i+\frac{m\omega}{2n+1}\right)}}{1-\frac{\varphi^{2}\beta}{\varphi^{2}\left(\frac{\omega}{2}+\frac{\varpi}{2}i+\frac{m\omega}{2n+1}\right)}}.\overset{n}{\underset{1}{\Pi}}_{\mu}\frac{1-\frac{\varphi^{2}\beta}{\varphi^{2}\left(\frac{\varpi}{2}i+\frac{\mu\varpi i}{2n+1}\right)}}{1-\frac{\varphi^{2}\beta}{\varphi^{2}\left(\frac{\omega}{2}+\frac{\varpi}{2}i+\frac{\mu\varpi i}{2n+1}\right)}}\\
&\quad\times\overset{n}{\underset{1}{\Pi}}_{m}\overset{n}{\underset{1}{\Pi}}_{\mu}\frac{1-\frac{\varphi^{2}\beta}{\varphi^{2}\left(\frac{\varpi}{2}i+\frac{m\omega+\mu\varpi i}{2n+1}\right)}}{1-\frac{\varphi^{2}\beta}{\varphi^{2}\left(\frac{\omega}{2}+\frac{\varpi}{2}i+\frac{m\omega+\mu\varpi i}{2n+1}\right)}}.\frac{1-\frac{\varphi^{2}\beta}{\varphi^{2}\left(\frac{\varpi}{2}i+\frac{m\omega-\mu\varpi i}{2n+1}\right)}}{1-\frac{\varphi^{2}\beta}{\varphi^{2}\left(\frac{\omega}{2}+\frac{\varpi}{2}i+\frac{m\omega-\mu\varpi i}{2n+1}\right)}}.
\end{aligned}\right. \tag{130$'$.} \]
\ednote{In the numerator of the second product of the first line, the print sets an $i$ after $\omega/2$; there is apparently no $i$ there, by analogy with the first product. Reproduced as printed.}

Dans ce paragraphe nous n'avons considéré les fonctions $\Phi(n\beta)$, $f(n\beta)$, $F(n\beta)$, que dans le cas des valeurs impaires de $n$. On pourroit trouver des expressions analogues de ces fonctions pour des valeurs paires de $n$; mais comme il n'y a dans cela aucune difficulté et que d'ailleurs les formules, aux quelles nous sommes parvenus, sont celles qui nous seront le plus utiles dans la suite, je ne m'en occuperai pas.

\section*{§. VII.}

\emph{Développement des fonctions} $\Phi\alpha$, $f\alpha$, $F\alpha$ \emph{en séries et en produits infinis.}

\subsection*{24.}

En faisant dans les formules du paragraphe précédent $\beta=\frac{\alpha}{2n+1}$, on obtiendra des expressions des fonctions $\Phi\alpha$, $f\alpha$, $F\alpha$, qui, à cause du nombre indéterminé $n$, peuvent être variées d'une infinité de manières.

\origpage{155}
Parmi toutes les formules, qu'on obtiendra ainsi, celles qui résultent de la supposition de $n$ infini, sont les plus remarquables. Alors les fonctions $\Phi$, $f$, $F$ disparaîtront des valeurs de $\Phi\alpha$, $f\alpha$, $F\alpha$, et on obtiendra pour ces fonctions des expressions algébriques, mais composées d'une infinité de termes. Pour avoir ces expressions, il faut faire dans les formules (126.), (130.) $\beta=\frac{\alpha}{2n+1}$, et ensuite chercher la limite du second membre de ces équations pour des valeurs toujours croissantes de $n$. Pour abréger, soit $v$ une quantité, dont la limite est zéro pour des valeurs toujours croissantes de $n$. Cela posé considérons successivement les trois formules (126.).

En faisant dans la première des formules (126.) $\beta=\frac{\alpha}{2n+1}$, et remarquant que
\[ \begin{aligned}
&\overset{+n}{\underset{-n}{\Sigma}}_{m}\overset{+n}{\underset{-n}{\Sigma}}_{\mu}\theta(m,\mu)=\theta(0,0)+\overset{n}{\underset{1}{\Sigma}}_{m}\theta(m,0)+\overset{n}{\underset{1}{\Sigma}}_{\mu}\theta(0,\mu)\\
&\quad+\overset{n}{\underset{1}{\Sigma}}_{m}\overset{n}{\underset{1}{\Sigma}}_{\mu}\{\theta(m,\mu)+\theta(-m,-\mu)+\theta(m,-\mu)+\theta(-m,\mu)\},
\end{aligned} \tag{131.} \]
il est clair, qu'on peut mettre la formule, dont il s'agit, sous la forme:
\[ \begin{aligned}
\Phi\alpha={}&\frac{1}{2n+1}.\varphi\left(\frac{\alpha}{2n+1}\right)+\frac{1}{2n+1}\overset{n}{\underset{1}{\Sigma}}_{m}(-1)^{m}\left\{\Phi\left(\frac{\alpha+m\omega}{2n+1}\right)+\Phi\left(\frac{\alpha-m\omega}{2n+1}\right)\right\}\\
&+\frac{1}{2n+1}\overset{n}{\underset{1}{\Sigma}}_{\mu}(-1)^{\mu}\left\{\Phi\left(\frac{\alpha+\mu\varpi i}{2n+1}\right)+\Phi\left(\frac{\alpha-\mu\varpi i}{2n+1}\right)\right\}-\frac{i}{ec}.\overset{n}{\underset{1}{\Sigma}}_{m}\overset{n}{\underset{1}{\Sigma}}_{\mu}(-1)^{m+\mu}.\psi(n-m,n-\mu)\\
&+\frac{i}{ec}\overset{n}{\underset{1}{\Sigma}}_{m}\overset{n}{\underset{1}{\Sigma}}_{\mu}(-1)^{m+\mu}.\psi_{1}(n-m,n-\mu),
\end{aligned} \tag{132.} \]
où l'on a fait, pour abréger,
\[ \left\{\begin{aligned}
\psi(m,\mu)&=\frac{1}{2n+1}.\left\{\frac{1}{\varphi\left(\frac{\alpha+(m+\tfrac{1}{2})\omega+(\mu+\tfrac{1}{2})\varpi i}{2n+1}\right)}+\frac{1}{\varphi\left(\frac{\alpha-(m+\tfrac{1}{2})\omega-(\mu+\tfrac{1}{2})\varpi i}{2n+1}\right)}\right\},\\
\psi_{1}(m,\mu)&=\frac{1}{2n+1}.\left\{\frac{1}{\varphi\left(\frac{\alpha+(m+\tfrac{1}{2})\omega-(\mu+\tfrac{1}{2})\varpi i}{2n+1}\right)}+\frac{1}{\varphi\left(\frac{\alpha-(m+\tfrac{1}{2})\omega+(\mu+\tfrac{1}{2})\varpi i}{2n+1}\right)}\right\}.
\end{aligned}\right. \tag{133.} \]

Maintenant, en remarquant, que
\[ \Phi\left(\frac{\alpha+m\omega}{2n+1}\right)+\Phi\left(\frac{\alpha-m\omega}{2n+1}\right)=\frac{2\varphi\left(\frac{\alpha}{2n+1}\right).f\left(\frac{m\omega}{2n+1}\right).F\left(\frac{m\omega}{2n+1}\right)}{1+e^{2}c^{2}.\varphi^{2}\left(\frac{m\omega}{2n+1}\right).\varphi^{2}\left(\frac{\alpha}{2n+1}\right)}=\frac{A_{m}}{2n+1}, \]
\[ \Phi\left(\frac{\alpha+\mu\varpi i}{2n+1}\right)+\Phi\left(\frac{\alpha-\mu\varpi i}{2n+1}\right)=\frac{2\varphi\left(\frac{\alpha}{2n+1}\right).f\left(\frac{\mu\varpi i}{2n+1}\right).F\left(\frac{\mu\varpi i}{2n+1}\right)}{1+e^{2}c^{2}.\varphi^{2}\left(\frac{\mu\varpi i}{2n+1}\right).\varphi^{2}\left(\frac{\alpha}{2n+1}\right)}=\frac{B_{\mu}}{2n+1}, \]

\origpage{156}
où $A_{m}$ et $B_{\mu}$ sont des quantités finies, la partie de l'équation (132.) jusqu'au membre, qui a le signe $-$, prendra la forme:
\[ \frac{1}{2n+1}.\varphi\left(\frac{\alpha}{2n+1}\right)+\frac{1}{(2n+1)^{2}}.\overset{n}{\underset{1}{\Sigma}}_{m}(-1)^{m}(A_{m}+B_{\mu}); \]
or la limite de cette quantité est évidemment zéro; donc, en prenant la limite de la formule (132.), on aura:
\[ \begin{aligned}
\Phi\alpha&=-\frac{i}{ec}\ \text{lim.}\ \overset{n}{\underset{1}{\Sigma}}_{m}\overset{n}{\underset{1}{\Sigma}}_{\mu}(-1)^{m+\mu}.\psi(n-m,n-\mu)\\
&\quad+\frac{i}{ec}\ \text{lim.}\ \overset{n}{\underset{1}{\Sigma}}_{m}\overset{n}{\underset{1}{\Sigma}}_{\mu}(-1)^{m+\mu}.\psi_{1}(n-m,n-\mu),
\end{aligned} \]
ou bien:
\[ \begin{aligned}
\Phi\alpha&=-\frac{i}{ec}\ \text{lim.}\ \overset{n-1}{\underset{0}{\Sigma}}_{m}\overset{n-1}{\underset{0}{\Sigma}}_{\mu}(-1)^{m+\mu}.\psi(m,\mu)\\
&\quad+\frac{i}{ec}\ \text{lim.}\ \overset{n-1}{\underset{0}{\Sigma}}_{m}\overset{n-1}{\underset{0}{\Sigma}}_{\mu}(-1)^{m+\mu}.\psi_{1}(m,\mu).
\end{aligned} \tag{134.} \]

Il suffit de connoître l'une de ces limites, car on aura l'autre, en changeant seulement le signe de $i$.

Cherchons la limite de
\[ \overset{n-1}{\underset{0}{\Sigma}}_{m}\overset{n-1}{\underset{0}{\Sigma}}_{\mu}(-1)^{m+\mu}.\psi(m,\mu). \]
Pour cela, il faut essayer de mettre la quantité précédente sous la forme
\[ P+v, \]
où $P$ est indépendant de $n$, et $v$ une quantité, qui a zéro pour limite; car alors la quantité $P$ sera précisement la limite, dont il s'agit.

\subsection*{25.}

Considérons d'abord l'expression:
\[ \overset{n-1}{\underset{0}{\Sigma}}_{\mu}(-1)^{\mu}.\psi(m,\mu). \]

Soit
\[ \theta(m,\mu)=\frac{2\alpha}{\alpha^{2}-(m\omega+\mu\varpi i)^{2}}, \tag{135.} \]
et faisons
\[ \psi(m,\mu)-\theta(m,\mu)=\frac{2\alpha}{(2n+1)^{2}}.R_{\mu}, \tag{136.} \]
on aura:
\[ \overset{n-1}{\underset{0}{\Sigma}}_{\mu}(-1)^{\mu}.\psi(m,\mu)-\overset{n-1}{\underset{0}{\Sigma}}_{\mu}(-1)^{\mu}.\theta(m,\mu)=2\alpha.\overset{n-1}{\underset{0}{\Sigma}}_{\mu}(-1)^{n}.\frac{R_{\mu}}{(2n+1)^{2}}. \tag{137.} \]
\ednote{The exponent of the factor $-1$ in the last sum is apparently printed $n$; $\mu$ is expected, as in the two sums on the left.}

Cela posé, je dis que le second membre de cette équation est une quantité de la forme $\frac{v}{2n+1}$.

\origpage{157}

D'après (12.), (13.) on aura
\[ \frac{1}{\varphi(\beta+\varepsilon)}+\frac{1}{\varphi(\beta-\varepsilon)}=\frac{\varphi(\beta+\varepsilon)+\varphi(\beta-\varepsilon)}{\varphi(\beta+\varepsilon).\varphi(\beta-\varepsilon)}=\frac{2\varphi\beta.f\varepsilon.F\varepsilon}{\varphi^{2}\beta-\varphi^{2}\varepsilon}, \]
donc en faisant $\beta=\frac{\alpha}{2n+1}$ et $\varepsilon=\frac{m\omega+\mu\varpi i}{2n+1}=\frac{\varepsilon_{\mu}}{2n+1}$ et $f\varepsilon.F\varepsilon=\theta\varepsilon$, on a:
\[ \psi(m,\mu)=\frac{1}{2n+1}.\frac{2\varphi\left(\frac{\alpha}{2n+1}\right).\theta\left(\frac{\varepsilon_{\mu}}{2n+1}\right)}{\varphi^{2}\left(\frac{\alpha}{2n+1}\right)-\varphi^{2}\left(\frac{\varepsilon_{\mu}}{2n+1}\right)}. \]
Maintenant on a:
\[ \varphi\left(\frac{\alpha}{2n+1}\right)=\frac{\alpha}{2n+1}+\frac{A\alpha^{3}}{(2n+1)^{3}}, \]
donc:
\[ \psi(m,\mu)=\frac{\theta\left(\frac{\varepsilon_{\mu}}{2n+1}\right)}{\varphi^{2}\left(\frac{\alpha}{2n+1}\right)-\varphi^{2}\left(\frac{\varepsilon_{\mu}}{2n+1}\right)}.\left\{\frac{2A\alpha^{3}}{(2n+1)^{4}}+\frac{2\alpha}{(2n+1)^{2}}\right\}, \]
et par conséquent:
\[ \begin{aligned}
\psi(m,\mu)-\theta(m,\mu)&=\frac{2\alpha}{(2n+1)^{2}}\left\{\frac{\theta\left(\frac{\varepsilon_{\mu}}{2n+1}\right)}{\varphi^{2}\left(\frac{\alpha}{2n+1}\right)-\varphi^{2}\left(\frac{\varepsilon_{\mu}}{2n+1}\right)}-\frac{1}{\left(\frac{\alpha}{2n+1}\right)^{2}-\left(\frac{\varepsilon_{\mu}}{2n+1}\right)^{2}}\right\}\\
&\quad+\frac{2A\alpha^{3}}{(2n+1)^{4}}.\frac{\theta\left(\frac{\varepsilon_{\mu}}{2n+1}\right)}{\varphi^{2}\left(\frac{\alpha}{2n+1}\right)-\varphi^{2}\left(\frac{\varepsilon_{\mu}}{2n+1}\right)}.
\end{aligned} \]
Donc la valeur de $R_{\mu}$ deviendra
\[ R_{\mu}=\frac{\theta\left(\frac{\varepsilon_{\mu}}{2n+1}\right)}{\varphi^{2}\left(\frac{\alpha}{2n+1}\right)-\varphi^{2}\left(\frac{\varepsilon_{\mu}}{2n+1}\right)}.\left\{1+\frac{A\alpha^{2}}{(2n+1)^{2}}\right\}-\frac{1}{\left(\frac{\alpha}{2n+1}\right)^{2}-\left(\frac{\varepsilon_{\mu}}{2n+1}\right)^{2}}. \tag{138.} \]

Cela posé, il y a deux cas à considérer, savoir si $\frac{\varepsilon_{\mu}}{2n+1}$ a zéro pour limite ou non.

$a)$ Si $\frac{\varepsilon_{\mu}}{2n+1}$ a zéro pour limite, on aura:
\[ \begin{aligned}
\varphi^{2}\left(\frac{\varepsilon_{\mu}}{2n+1}\right)&=\frac{\varepsilon_{\mu}^{2}}{(2n+1)^{2}}+\frac{B_{\mu}.\varepsilon_{\mu}^{4}}{(2n+1)^{4}},\\
\theta\left(\frac{\varepsilon_{\mu}}{2n+1}\right)&=\sqrt{\left[1-c^{2}\varphi^{2}\left(\frac{\varepsilon_{\mu}}{2n+1}\right)\right]}.\sqrt{\left[1+e^{2}\varphi^{2}\left(\frac{\varepsilon_{\mu}}{2n+1}\right)\right]}=1+\frac{C_{\mu}.\varepsilon_{\mu}^{2}}{(2n+1)^{2}},\\
\varphi^{2}\left(\frac{\alpha}{2n+1}\right)&=\frac{\alpha^{2}}{(2n+1)^{2}}+\frac{D.\alpha^{4}}{(2n+1)^{4}},
\end{aligned} \]
ou $B_{\mu}$, $C_{\mu}$, $D$ ont des limites finies; donc en substituant:

\origpage{158}
\[ \begin{aligned}
R_{\mu}&=A\alpha^{2}.\frac{\frac{1}{\varepsilon_{\mu}^{2}}+\frac{C_{\mu}}{(2n+1)^{2}}}{\frac{\alpha^{2}}{\varepsilon_{\mu}^{2}}-1+\frac{D\alpha^{4}}{(2n+1)^{2}.\varepsilon_{\mu}^{2}}-B_{\mu}\frac{\varepsilon_{\mu}^{2}}{(2n+1)^{2}}}\\
&\quad+\frac{C_{\mu}.\frac{\alpha^{2}}{\varepsilon_{\mu}^{2}}-\frac{D\alpha^{4}}{\varepsilon_{\mu}^{4}}-C_{\mu}+B_{\mu}}{\left(1-\frac{\alpha^{2}}{\varepsilon_{\mu}^{2}}\right)^{2}-\left(1-\frac{\alpha^{2}}{\varepsilon_{\mu}^{2}}\right)\left\{\frac{D\alpha^{4}}{(2n+1)^{2}.\varepsilon_{\mu}^{2}}-B_{\mu}.\frac{\varepsilon_{\mu}^{2}}{(2n+1)^{2}}\right\}},
\end{aligned} \tag{139.} \]
or soit que $\varepsilon_{\mu}$ soit fini ou infini, il est clair, que cette quantité convergera toujours vers une quantité finie pour des valeurs toujours croissantes de $n$. Donc on aura
\[ R_{\mu}=r_{\mu}+v_{\mu}, \tag{140.} \]
où $r_{\mu}$ est une quantité finie indépendante de $n$.

$b)$ Si $\frac{\varepsilon_{\mu}}{2n+1}$ a pour limite une quantité finie, il est clair, qu'en nommant cette limite $\delta_{\mu}$, on aura:
\[ R_{\mu}=-\frac{\theta(\delta_{\mu})}{\varphi^{2}(\delta_{\mu})}+\frac{1}{\delta_{\mu}^{2}}+v_{\mu}^{1}. \tag{141.} \]
Cela posé, considérons l'expression $\overset{n}{\underset{1}{\Sigma}}_{\mu}(-1)^{\mu}.\frac{R_{\mu}}{(2n+1)^{2}}$. On a
\[ \begin{aligned}
\overset{n}{\underset{1}{\Sigma}}_{\mu}(-1)^{\mu}.\frac{R_{\mu}}{(2n+1)^{2}}&=\frac{1}{(2n+1)^{2}}.\{R_{0}-R_{1}+R_{2}-R_{3}+....\\
&\quad....+(-1)^{\nu-1}.R_{\nu-1}+(-1)^{\nu+1}[R_{\nu+1}-R_{\nu+1}+R_{\nu+2}-R_{\nu+3}...+(-1)^{n-\nu-1}.R_{n-1}]\}.
\end{aligned} \tag{142.} \]
\ednote{The first term inside the square bracket is printed $R_{\nu+1}$; $R_{\nu}$ is apparently expected.}
Supposons d'abord, que $\frac{\varepsilon_{\mu}}{2n+1}$ ait pour limite une quantité finie, quelle que soit la valeur de $\mu$. Alors en remarquant, que
\[ \delta_{\mu+1}=\delta_{\mu}, \]
on aura
\[ R_{\mu}-R_{\mu+1}=v_{\mu}^{1}-v_{\mu+1}^{1}, \]
donc:
\[ \overset{n-1}{\underset{0}{\Sigma}}_{\mu}(-1)^{\mu}.\frac{R_{\mu}}{(2n+1)^{2}}=\frac{1}{(2n+1)^{2}}(v_{0}^{1}-v_{1}^{1}+v_{2}^{1}-v_{3}^{1}+....+v_{k-2}^{1}-v_{k-1}^{1})+\frac{B}{(2n+1)^{2}}, \]
où $k=n$ ou $n-1$, selon que $n$ est pair ou impair. La quantité $B$ a toujours pour limite une quantité finie, savoir $B=0$, si $n$ est pair, et $B=R_{n-1}$, si $n$ est impair.

Maintenant on sait, qu'une somme telle que
\[ v_{0}^{1}-v_{1}^{1}+v_{2}^{1}-....+v_{k-2}^{1}-v_{k-1}^{1}, \]
peut être mise sous la forme $k.v$, $v$ ayant zéro pour limite. Donc en substituant:

\origpage{159}
\[ \overset{n-1}{\underset{0}{\Sigma}}_{\mu}(-1)^{\mu}.\frac{R_{\mu}}{(2n+1)^{2}}=\frac{k.v+B}{(2n+1)^{2}}; \]
or, $k$ étant égal à $n$ ou à $n-1$, et $B$ fini, la limite de $\frac{k.v+B}{2n+1}$ sera zéro, donc:
\[ \overset{n-1}{\underset{0}{\Sigma}}_{\mu}(-1)^{\mu}.\frac{R_{\mu}}{(2n+1)^{2}}=\frac{v}{2n+1}. \tag{143.} \]

Supposons maintenant, que $\frac{m}{2n+1}$ a zéro pour limite. Alors $\frac{\varepsilon_{\mu}}{2n+1}$ a également zéro pour limite, à moins qu'en même tems $\frac{\mu}{2n+1}$ n'ait pour limite une quantité finie. Soit dans ce cas $\nu$ le nombre entier immédiatement inferieur à $\sqrt{n}$, et considérons la somme
\[ R_{0}-R_{1}+R_{2}-R_{3}+....+(-1)^{\nu-1}R_{\nu-1}. \]
En supposant, que $\mu$ est un des nombres $0$, $1$, $....\nu$, il est clair, que $\frac{\varepsilon_{\mu}}{(2n+1)}=\frac{m\omega+\mu\varpi i}{2n+1}$ a zéro pour limite; donc, selon ce qu'on a vu, $R_{\mu}$ sera une quantité finie, et par conséquent
\[ R_{0}-R_{1}+R_{2}-....+(-1)^{\nu-1}R_{\nu-1}=\nu.R, \]
où $R$ est également une quantité finie.

Considérons maintenant la somme
\[ (-1)^{\nu}\{R_{\nu}-R_{\nu+1}+R_{\nu+2}-....+(-1)^{n-\nu-1}.R_{n-1}\}. \]

Si $\frac{\varepsilon_{\mu}}{2n+1}$ a pour limite une quantité différente de zéro, on a, comme on verra:
\[ R_{\mu}-R_{\mu+1}=v_{\mu}^{1}-v_{\mu+1}^{1}; \]
si au contraire $\frac{\varepsilon_{\mu}}{2n+1}$ a pour limite zéro, on a:
\[ R_{\mu}=v_{\mu}+v_{\mu}^{1}; \]
or, si en même tems $\mu>\sqrt{n}$, il est clair, qu'en vertu de la valeur de $R_{\mu}$,
\[ v_{\mu}=B_{\mu}-C_{\mu}; \]
or il est clair, que $B_{\mu}$ et $C_{\mu}$, tous deux, ont pour limites des quantités indépendantes de $\mu$, donc en nommant ces limites $B$ et $C$, on aura:
\[ R_{\mu}=B-C+v_{\mu}, \]
et par suite, aussi dans ce cas,
\[ R_{\mu}-R_{\mu+1}=v_{\mu}-v_{\mu+1}. \]
\ednote{The last subscript is printed indistinctly, apparently $v_{\mu+1}$ with the $\mu$ damaged in print.}

Donc comme dans le cas où $\frac{\varepsilon_{\mu}}{2n+1}$ auroit une limite différente de zéro pour toutes les valeurs de $\mu$, on démontrera que
\[ \frac{(-1)^{\nu}}{(2n+1)^{2}}\{R_{\nu}-R_{\nu+1}+....+(-1)^{n-\nu-1}R_{n-1}\}=\frac{v}{(2n+1)}. \]

\origpage{160}
Maintenant en combinant les équations ci-dessus, on en tirera
\[ \overset{n-1}{\underset{0}{\Sigma}}_{\mu}(-1)^{\mu}.\frac{R_{\mu}}{(2n+1)^{2}}=\frac{1}{(2n+1)^{2}}.(\nu R)+\frac{v}{2n+1}; \]
or $\frac{v}{2n+1}$ a zéro pour limite, donc
\[ \overset{n-1}{\underset{0}{\Sigma}}_{\mu}(-1)^{\mu}.\frac{R_{\mu}}{(2n+1)^{2}}=\frac{v}{2n+1}. \]
Donc cette formule a toujours lieu, et par conséquent la formule (137.) deviendra:
\[ \overset{n-1}{\underset{0}{\Sigma}}_{\mu}(-1)^{\mu}\psi(m,\mu)-\overset{n-1}{\underset{0}{\Sigma}}_{\mu}(-1)^{\mu}.\theta(m,\mu)=\frac{v}{2n+1}. \tag{144.} \]
Cela posé, il s'agit de mettre $\overset{n-1}{\underset{0}{\Sigma}}_{\mu}(-1)^{\mu}\theta(m,\mu)$ sous la forme $P+\frac{v}{2n+1}$.
Or c'est ce qu'on peut faire comme il suit. On a:
\[ \left\{\begin{aligned}
&\overset{n-1}{\underset{0}{\Sigma}}_{\mu}(-1)^{\mu}\theta(m,\mu)=\overset{\frac{1}{0}}{\underset{0}{\Sigma}}_{\mu}(-1)^{\mu}\theta(m,\mu)-\overset{\frac{1}{0}}{\underset{n}{\Sigma}}_{\mu}(-1)^{\mu}\theta(m,\mu)\ \text{et}\\
&\overset{\frac{1}{0}}{\underset{n}{\Sigma}}_{\mu}(-1)^{\mu}.\theta(m,\mu)=(-1)^{n}\{\theta(m,n)-\theta(m,n+1)+\theta(m,n+2)-....\text{etc.}\}.
\end{aligned}\right. \tag{145.} \]
Or d'après une formule connue on a:
\[ \begin{aligned}
&\theta(m,n)-\theta(m,n+1)+\theta(m,n+2)-......\\
&=\tfrac{1}{2}\theta(m,n)+A\frac{\partial\theta(m,n)}{\partial n}+B.\frac{\partial^{3}\theta(m,n)}{\partial^{3}n}+......,
\end{aligned} \]
où $A$, $B$$....$ sont des nombres; or
\[ \theta(m,n)=\frac{2\alpha}{\alpha^{2}-(m\omega+n\varpi i)^{2}}, \]
donc en substituant:
\[ \theta(m,n)-\theta(m,n+1)+....=\frac{\alpha}{\alpha^{2}-(m\omega+n\varpi i)^{2}}+\frac{4A\alpha\varpi i.n}{\{\alpha^{2}-(m\omega+n\varpi i)^{2}\}^{2}}+.... \]
\ednote{The exponent outside the braces in the denominator of the second term is printed indistinctly; $2$ is read here.}
De là il suit, que
\[ \theta(m,n)-\theta(m,n+1)+....=\frac{\alpha}{\varpi^{2}.n^{2}}+\frac{v}{n^{2}}=\frac{v}{2n+1}. \]
Donc en vertu des équations (145.)
\[ \overset{n-1}{\underset{0}{\Sigma}}_{\mu}(-1)^{\mu}\theta(m,\mu)=\overset{\frac{1}{0}}{\underset{0}{\Sigma}}_{\mu}(-1)^{\mu}.\theta(m,\mu)+\frac{v}{2n+1}, \]
et par conséquent
\[ \overset{n-1}{\underset{0}{\Sigma}}_{\mu}(-1)^{\mu}\psi(m,\mu)=\overset{\infty}{\underset{0}{\Sigma}}_{\mu}(-1)^{\mu}.\theta(m,\mu)+\frac{v}{2n+1}. \tag{146.} \]

\subsection*{26.}

Ayant transformé de cette sorte la quantité $\overset{n-1}{\underset{0}{\Sigma}}_{\mu}(-1)^{\mu}.\psi(m,\mu)$, on tire de l'équation (146.):

\origpage{161}
\[ \overset{n-1}{\underset{0}{\Sigma}}_{m}\overset{n-1}{\underset{0}{\Sigma}}_{\mu}(-1)^{m+\mu}\psi(m,\mu)=\overset{n-1}{\underset{0}{\Sigma}}(-1)^{m}.\varrho_{m}+\overset{n-1}{\underset{0}{\Sigma}}.\frac{v_{m}}{2n+1}, \tag{147.} \]
en faisant
\[ \varrho_{m}=\overset{\infty}{\underset{0}{\Sigma}}_{\mu}(-1)^{\mu}.\theta(m,\mu); \tag{148.} \]
or
\[ \overset{n-1}{\underset{0}{\Sigma}}.\frac{v_{m}}{2n+1}=\frac{v_{0}+v_{1}+v_{2}\ldots v_{n-1}}{2n+1}=\frac{nv}{2n+1}=\frac{v}{2}, \]
$v$ ayant zéro pour limite. Donc l'équation (147.) donnera, en faisant $n$ infini:
\[ \text{lim.}\overset{n-1}{\underset{0}{\Sigma}}_{m}\overset{n-1}{\underset{0}{\Sigma}}_{\mu}(-1)^{m+\mu}.\psi(m,\mu)=\overset{\infty}{\underset{0}{\Sigma}}(-1)^{m}.\varrho_{m}. \tag{149.} \]
De même, si l'on fait, pour abréger:
\[ \left\{\begin{aligned}\theta_{1}(m,\mu)&=\frac{2\alpha}{\alpha^{2}-(m\omega-\mu\varpi i)^{2}},\\ \varrho'_{m}&=\overset{\infty}{\underset{0}{\Sigma}}_{\mu}(-1)^{\mu}.\theta_{1}(m,\mu),\end{aligned}\right. \tag{150.} \]
on aura:
\[ \text{lim.}\overset{n-1}{\underset{0}{\Sigma}}_{m}\overset{n-1}{\underset{0}{\Sigma}}_{\mu}(-1)^{m+\mu}.\psi_{1}(m,\mu)=\overset{\infty}{\underset{0}{\Sigma}}(-1)^{m}.\varrho'_{m}. \tag{151.} \]
Ayant trouvé ces deux quantités, dont l'expression de $\Phi\alpha$ est composée, on aura, en substituant:
\[ \Phi\alpha=-\frac{i}{ec}.\overset{\infty}{\underset{0}{\Sigma}}(-1)^{m}.\varrho_{m}+\frac{i}{ec}.\overset{\infty}{\underset{0}{\Sigma}}(-1)^{m}.\varrho'_{m}=\frac{i}{ec}.\overset{\infty}{\underset{0}{\Sigma}}(-1)^{m}\{\varrho'_{m}-\varrho_{m}\}, \]
ou bien, en remettant les valeurs de $\varrho'_{m}$ et $\varrho_{m}$,
\[ \Phi\alpha=\frac{i}{ec}.\overset{\infty}{\underset{0}{\Sigma}}_{m}(-1)^{m}\left\{\overset{\infty}{\underset{0}{\Sigma}}_{\mu}\left\{\frac{2\alpha}{\alpha^{2}-[(m+\tfrac{1}{2})\omega-(\mu+\tfrac{1}{2})\varpi i]^{2}}-\frac{2\alpha}{\alpha^{2}-[(m+\tfrac{1}{2})\omega+(\mu+\tfrac{1}{2})\varpi i]^{2}}\right\}\right\}. \tag{152.} \]
Maintenant
\[ \frac{2\alpha}{\alpha^{2}-[(m+\tfrac{1}{2})\omega\pm(\mu+\tfrac{1}{2})\varpi i]^{2}}=\frac{1}{\alpha-(m+\tfrac{1}{2})\omega\mp(\mu+\tfrac{1}{2})\varpi i}+\frac{1}{\alpha+(m+\tfrac{1}{2})\omega\pm(\mu+\tfrac{1}{2})\varpi i}, \]
donc:
\[ \begin{aligned}&\frac{2\alpha}{\alpha^{2}-[(m+\tfrac{1}{2})\omega-(\mu+\tfrac{1}{2})\varpi i]^{2}}-\frac{2\alpha}{\alpha^{2}-[(m+\tfrac{1}{2})\omega+(\mu+\tfrac{1}{2})\varpi i]^{2}}\\ &=+\frac{(2\mu+1)\varpi i}{[\alpha+(m+\tfrac{1}{2})\omega]^{2}+(\mu+\tfrac{1}{2})^{2}\varpi^{2}}-\frac{(2\mu+1)\varpi i}{[\alpha-(m+\tfrac{1}{2})\omega]^{2}+(\mu+\tfrac{1}{2})^{2}\varpi^{2}},\end{aligned} \]
donc l'expression de $\Phi\alpha$ deviendra sous une forme réelle:
\[ \Phi\alpha=\frac{1}{ec}.\overset{\infty}{\underset{0}{\Sigma}}_{m}(-1)^{m}.\overset{\infty}{\underset{0}{\Sigma}}_{\mu}(-1)^{\mu}.\left\{\frac{(2\mu+1)\varpi}{[\alpha-(m+\tfrac{1}{2})\omega]^{2}+(\mu+\tfrac{1}{2})^{2}\varpi^{2}}-\frac{(2\mu+1)\varpi}{[\alpha+(m+\tfrac{1}{2})\omega]^{2}+(\mu+\tfrac{1}{2})^{2}\varpi^{2}}\right\}, \tag{153.} \]
c'est-à-dire, on aura:

\origpage{162}
\[ \begin{aligned}\Phi\alpha&=\frac{\varpi}{ec}(\delta_{0}-\delta_{1}+\delta_{2}-\delta_{3}+....+(-1)^{m}\delta_{m}-....)\\ &-\frac{\varpi}{ec}(\delta'_{0}-\delta'_{1}+\delta'_{2}-\delta'_{3}+....+(-1)^{m}\delta'_{m}-....),\end{aligned} \tag{154.} \]
ou bien:
\[ \left\{\begin{aligned}\delta_{m}&=\frac{1}{[\alpha-(m+\tfrac{1}{2})\omega]^{2}+\frac{\varpi^{2}}{4}}-\frac{3}{[\alpha-(m+\tfrac{1}{2})\omega]^{2}+\frac{9\varpi^{2}}{4}}+\frac{5}{[\alpha-(m+\tfrac{1}{2})\omega]^{2}+\frac{25\varpi^{2}}{4}}-\ldots\\ \delta'_{\mu}&=\frac{1}{[\alpha+(m+\tfrac{1}{2})\omega]^{2}+\frac{\varpi^{2}}{4}}-\frac{3}{[\alpha+(m+\tfrac{1}{2})\omega]^{2}+\frac{9\varpi^{2}}{4}}+\frac{5}{[\alpha+(m+\tfrac{1}{2})\omega]^{2}+\frac{25\varpi^{2}}{4}}-\ldots\end{aligned}\right. \tag{155.} \]

Si l'on commence la recherche de la limite de la fonction $\overset{n-1}{\underset{0}{\Sigma}}_{m}\overset{n-1}{\underset{0}{\Sigma}}_{\mu}(-1)^{m+\mu}.\psi(m,\mu)$ par celle de $\overset{n-1}{\underset{0}{\Sigma}}(-1)^{m}\psi(m,\mu)$ au lieu de celle de $\overset{n-1}{\underset{0}{\Sigma}}(-1)^{\mu}\psi(m,\mu)$, comme nous l'avons fait, on trouvera au lieu de la formule (53.)\ednote{Printed (53.); the formula meant is presumably (153.), an apparent misprint.} la suivante:
\[ \Phi\alpha=\frac{1}{ec}.\overset{\infty}{\underset{0}{\Sigma}}_{\mu}(-1)^{\mu}.\overset{\infty}{\underset{0}{\Sigma}}_{m}(-1)^{m}.\left\{\frac{(2\mu+1)\varpi}{[\alpha-(m+\tfrac{1}{2})\omega]^{2}+(\mu+\tfrac{1}{2})^{2}\varpi^{2}}-\frac{(2\mu+1)\varpi}{[\alpha+(m+\tfrac{1}{2})\omega]^{2}+(\mu+\tfrac{1}{2})^{2}\varpi^{2}}\right\}, \tag{156.} \]
c'est-à-dire:
\[ \begin{aligned}\Phi\alpha&=\frac{\varpi}{ec}(\varepsilon_{0}-3\varepsilon_{1}+5\varepsilon_{2}-7\varepsilon_{3}+....+(-1)^{\mu}.(2\mu+1)\varepsilon_{\mu}-....)\\ &-\frac{\varpi}{ec}(\varepsilon'_{0}-3\varepsilon'_{1}+5\varepsilon'_{2}-7\varepsilon'_{3}+....+(-1)^{\mu}.(2\mu+1)\varepsilon_{\mu}-....),\end{aligned} \tag{157.} \]
\ednote{In the second row the last term is printed with $\varepsilon_{\mu}$ without the prime; $\varepsilon'_{\mu}$ is expected, an apparent misprint.}
où
\[ \left\{\begin{aligned}\varepsilon_{\mu}&=\frac{1}{\left(\alpha-\frac{\omega}{2}\right)^{2}+(\mu+\tfrac{1}{2})^{2}\varpi^{2}}-\frac{1}{\left(\alpha-\frac{3\omega}{2}\right)^{2}+(\mu+\tfrac{1}{2})^{2}\varpi^{2}}+\frac{1}{\left(\alpha-\frac{5\omega}{2}\right)^{2}+(\mu+\tfrac{1}{2})^{2}\varpi^{2}}-....\\ \varepsilon'_{\mu}&=\frac{1}{\left(\alpha+\frac{\omega}{2}\right)^{2}+(\mu+\tfrac{1}{2})^{2}\varpi^{2}}-\frac{1}{\left(\alpha+\frac{3\omega}{2}\right)^{2}+(\mu+\tfrac{1}{2})^{2}\varpi^{2}}+\frac{1}{\left(\alpha+\frac{5\omega}{2}\right)^{2}+(\mu+\tfrac{1}{2})^{2}\varpi^{2}}-....\end{aligned}\right. \tag{158.} \]

\subsection*{27.}

Cherchons maintenant l'expression de $f\alpha$, au moyen de la deuxième des formules (126.). En y faisant $\beta=\frac{\alpha}{2n+1}$, et remarquant, qu'alors la limite des quantités contenues dans les deux premières lignes, devient égale à zéro, on aura:
\[ \begin{aligned}f\alpha&=\text{lim.}(-1)^{n}\overset{n}{\underset{1}{\Sigma}}_{m}\overset{n}{\underset{1}{\Sigma}}_{\mu}(-1)^{m}.\psi(n-m,n-\mu)\\ &+\text{lim.}(-1)^{n}\overset{n}{\underset{1}{\Sigma}}_{m}\overset{n}{\underset{1}{\Sigma}}_{\mu}(-1)^{m}.\psi_{1}(n-m,n-\mu),\end{aligned} \tag{159.} \]

\origpage{163}
ou l'on a fait pour abréger:
\[ \begin{aligned}\psi(n-m,n-\mu)&=\frac{1}{2n+1}\left\{f\left(\frac{\alpha+m\omega+\mu\varpi i}{2n+1}\right)+f\left(\frac{\alpha-m\omega-\mu\varpi i}{2n+1}\right)\right\},\\ \psi_{1}(n-m,n-\mu)&=\frac{1}{2n+1}\left\{f\left(\frac{\alpha+m\omega-\mu\varpi i}{2n+1}\right)+f\left(\frac{\alpha-m\omega+\mu\varpi i}{2n+1}\right)\right\}.\end{aligned} \]
Maintenant on a:
\[ f(\beta-\varepsilon)+f(\beta-\varepsilon)=\frac{2f\beta.f\varepsilon}{1+e^{2}c^{2}\varphi^{2}\varepsilon.\varphi^{2}\beta}=\frac{f\varepsilon}{e^{2}c^{2}\varphi^{2}\varepsilon}.\frac{f\beta}{\varphi^{2}\beta+\frac{1}{e^{2}c^{2}\varphi^{2}\varepsilon}}. \]
\ednote{Printed $f(\beta-\varepsilon)$ in both terms; the first term is presumably $f(\beta+\varepsilon)$, an apparent misprint.}

Soit
\[ \varepsilon=\frac{m\omega+\mu\varpi i}{2n+1}, \]
on aura:
\[ \begin{aligned}\frac{1}{\varphi\varepsilon}&=-iec.\Phi\left(\frac{\omega}{2}+\frac{\varpi}{2}i-\varepsilon\right)=-iec.\Phi\left(\frac{(n-m+\tfrac{1}{2})\omega+(n-\mu+\tfrac{1}{2})\varpi i}{2n+1}\right),\\ \frac{f\varepsilon}{\varphi\varepsilon}&=-\frac{i\sqrt{(e^{2}+c^{2})}}{F\left(\varepsilon-\frac{\varpi}{2}i\right)}=-i\sqrt{(e^{2}+c^{2})}.\frac{c}{\sqrt{(e^{2}+c^{2})}}.F\left(\varepsilon-\frac{\omega}{2}-\frac{\varpi}{2}i\right)\\ &=-ci.F.\left(\frac{(n-m+\tfrac{1}{2})\omega+(n-\mu+\tfrac{1}{2})\varpi i}{2n+1}\right).\end{aligned} \]
Donc on aura, en substituant et mettant $m$ et $\mu$ respectivement au lieu de $n-m$ et $n-\mu$:
\[ \psi(m,\mu)=\frac{1}{e}.\frac{2\varphi\left(\frac{(m+\tfrac{1}{2})\omega+(\mu+\tfrac{1}{2})\varpi i}{2n+1}\right).F\left(\frac{(m+\tfrac{1}{2})\omega+(\mu+\tfrac{1}{2})\varpi i}{2n+1}\right)}{(2n+1)\left[\varphi^{2}\left(\frac{\alpha}{2n+1}\right)-\varphi^{2}\left(\frac{(m+\tfrac{1}{2})\omega+(\mu+\tfrac{1}{2})\varpi i}{2n+1}\right)\right]}. \]
On aura la valeur de $\psi_{1}(m,\mu)$, en changeant seulement le signe de $i$. En faisant maintenant
\[ \theta(m,\mu)=\frac{(2m+1)\omega+(2\mu+1)\varpi i}{\alpha^{2}-[(m+\tfrac{1}{2})\omega+(\mu+\tfrac{1}{2})\varpi i]^{2}} \]
et
\[ \theta_{1}(m,\mu)=\frac{(2m+1)\omega-(2\mu+1)\varpi i}{\alpha^{2}-[(m+\tfrac{1}{2})\omega-(\mu+\tfrac{1}{2})\varpi i]^{2}}, \]
et cherchant ensuite la limite de la fonction
\[ \overset{n-1}{\underset{0}{\Sigma}}_{m}\overset{n-1}{\underset{0}{\Sigma}}_{\mu}(-1)^{m}.\psi(m,\mu), \]
de la même manière que précédemment, on trouvera:
\[ \text{lim.}\overset{n-1}{\underset{0}{\Sigma}}_{m}\overset{n-1}{\underset{0}{\Sigma}}_{\mu}(-1)^{m}.\psi(m,\mu)=\frac{1}{e}.\overset{\infty}{\underset{0}{\Sigma}}_{\mu}\left\{\overset{\infty}{\underset{0}{\Sigma}}_{m}(-1)^{m}.\theta(m,\mu)\right\} \]
et
\[ \text{lim.}\overset{n-1}{\underset{0}{\Sigma}}_{m}\overset{n-1}{\underset{0}{\Sigma}}_{\mu}(-1)^{m}.\psi_{1}(m,\mu)=\frac{1}{e}.\overset{\infty}{\underset{0}{\Sigma}}_{\mu}\left\{\overset{\infty}{\underset{0}{\Sigma}}_{m}(-1)^{m}.\theta_{1}(m,\mu)\right\}; \]

\origpage{164}
donc en substituant dans (159.), et remettant les valeurs de $\theta(m,\mu)$ et $\theta_{1}(m,\mu)$, on a:
\[ f\alpha=\frac{1}{e}.\overset{\infty}{\underset{0}{\Sigma}}_{\mu}\overset{\infty}{\underset{0}{\Sigma}}_{m}(-1)^{m}\left\{\frac{(2m+1)\omega+(2\mu+1)\varpi i}{\alpha^{2}-[(m+\tfrac{1}{2})\omega+(\mu+\tfrac{1}{2})\varpi i]^{2}}+\frac{(2m+1)\omega-(2\mu+1)\varpi i}{\alpha^{2}-[(m+\tfrac{1}{2})\omega-(\mu+\tfrac{1}{2})\varpi i]^{2}}\right\}. \tag{160.} \]

La quantité renfermée entre les crochets peut aussi se mettre sous la forme:
\[ \frac{2[\alpha-(m+\tfrac{1}{2})\omega]}{[\alpha-(m+\tfrac{1}{2})\omega]^{2}+(\mu+\tfrac{1}{2})^{2}\varpi^{2}}-\frac{2[\alpha+(m+\tfrac{1}{2})\omega]}{[\alpha+(m+\tfrac{1}{2})\omega]^{2}+(\mu+\tfrac{1}{2})^{2}\varpi^{2}}, \]
donc aussi:
\[ f\alpha=\frac{1}{e}.\overset{\infty}{\underset{0}{\Sigma}}_{\mu}\left\{\overset{\infty}{\underset{0}{\Sigma}}_{m}(-1)^{m}.\frac{2[\alpha-(m+\tfrac{1}{2})\omega]}{[\alpha-(m+\tfrac{1}{2})\omega]^{2}+(\mu+\tfrac{1}{2})^{2}\varpi^{2}}-\overset{\infty}{\underset{0}{\Sigma}}_{m}(-1)^{m}\frac{2[\alpha+(m+\tfrac{1}{2})\omega]}{[\alpha+(m+\tfrac{1}{2})\omega]^{2}+(\mu+\tfrac{1}{2})^{2}\varpi^{2}}\right\}. \tag{161.} \]

On aura de la même manière:
\[ F(\alpha)=\frac{1}{c}.\overset{\infty}{\underset{0}{\Sigma}}_{m}\left\{\overset{\infty}{\underset{0}{\Sigma}}_{\mu}(-1)^{\mu}.\frac{(2\mu+1)\varpi}{[\alpha-(m+\tfrac{1}{2})\omega]^{2}+(\mu+\tfrac{1}{2})\varpi^{2}}+\overset{\infty}{\underset{0}{\Sigma}}(-1)^{\mu}.\frac{(2\mu+1)\varpi}{[\alpha+(m+\tfrac{1}{2})\omega]^{2}+(\mu+\tfrac{1}{2})\varpi^{2}}\right\}. \tag{162.} \]
\ednote{Both denominators are printed with $(\mu+\tfrac{1}{2})\varpi^{2}$; the square on the bracket, $(\mu+\tfrac{1}{2})^{2}\varpi^{2}$, is expected as in (161.), an apparent misprint.}

\subsection*{28.}

Venons maintenant aux formules (130.). Pour trouver la valeur du second membre, après avoir fait $\beta=\frac{\alpha}{2n+1}$, et supposé $n$ infini, nous allons d'abord chercher la limite de l'expression suivante:
\[ t=\overset{n}{\underset{1}{\Pi}}_{m}\overset{n}{\underset{1}{\Pi}}_{\mu}\left\{\frac{1-\frac{\varphi^{2}\left(\frac{\alpha}{2n+1}\right)}{\varphi^{2}\left(\frac{m\omega+\mu\varpi i+k}{2n+1}\right)}}{1-\frac{\varphi^{2}\left(\frac{\alpha}{2n+1}\right)}{\varphi^{2}\left(\frac{m\omega+\mu\varpi i+l}{2n+1}\right)}}\right\}, \tag{163.} \]
où $k$ et $l$ sont deux quantités indépendantes de $n$, $m$, $\mu$.

En prenant le logarithme, et faisant, pour abréger:
\[ \psi(m,\mu)=\log\left\{\frac{1-\frac{\varphi^{2}\left(\frac{\alpha}{2n+1}\right)}{\varphi^{2}\left(\frac{m\omega+\mu\varpi i+k}{2n+1}\right)}}{1-\frac{\varphi^{2}\left(\frac{\alpha}{2n+1}\right)}{\varphi^{2}\left(\frac{m\omega+\mu\varpi i+l}{2n+1}\right)}}\right\}, \tag{164.} \]

\origpage{165}
on aura:
\[ \log t=\overset{n}{\underset{1}{\Sigma}}_{m}\overset{n}{\underset{1}{\Sigma}}_{\mu}.\psi(m,\mu). \tag{165.} \]

Considérons d'abord l'expression: $\overset{n}{\underset{1}{\Sigma}}_{\mu}\psi(m,\mu)$. Soit
\[ \theta(m,\mu)=\log\left\{\frac{1-\frac{\alpha^{2}}{(m\omega+\mu\varpi i+k)^{2}}}{1-\frac{\alpha^{2}}{(m\omega+\mu\varpi i+l)^{2}}}\right\}, \tag{166.} \]
on aura:
\[ \psi(m,\mu)-\theta(m,\mu)=\log\left\{\frac{1-\frac{\varphi^{2}\left(\frac{\alpha}{2n+1}\right)}{\varphi^{2}\left(\frac{m\omega+\mu\varpi i+k}{2n+1}\right)}}{1-\frac{\alpha^{2}}{(m\omega+\mu\varpi i+k)^{2}}}.\frac{1-\frac{\alpha^{2}}{(m\omega+\mu\varpi i+l)^{2}}}{1-\frac{\varphi^{2}\left(\frac{\alpha}{2n+1}\right)}{\varphi^{2}\left(\frac{m\omega+\mu\varpi i+l}{2n+1}\right)}}\right\}. \]
Cela posé, je dis que le second membre de cette équation est pour toute valeur de $m$ et $\mu$ de la forme:
\[ \psi(m,\mu)-\theta(m,\mu)=\frac{v}{(2n+1)^{2}}. \]

Pour démontrer cela, il faut distinguer deux cas, si la limite de $\frac{m\omega+\mu\varpi i}{2n+1}$ est une quantité différente de zéro, et si elle est égale à zéro.

$a)$ Dans le premier cas on aura, en nommant $a$ la limite, dont il s'agit:
\[ \begin{aligned}\varphi^{2}\left(\frac{m\omega+\mu\varpi i+k}{2n+1}\right)&=\Phi^{2}a+v,\\ \varphi^{2}\left(\frac{m\omega+\mu\varpi i+l}{2n+1}\right)&=\Phi^{2}a+v',\\ \varphi^{2}\left(\frac{\alpha}{2n+1}\right)&=\frac{\alpha^{2}}{(2n+1)^{2}}+\frac{v'}{(2n+1)^{2}},\end{aligned} \]
donc:
\[ \begin{aligned}1-\frac{\varphi^{2}\left(\frac{\alpha}{2n+1}\right)}{\varphi^{2}\left(\frac{m\omega+\mu\varpi i+k}{2n+1}\right)}&=1-\frac{\alpha^{2}}{(2n+1)^{2}\varphi^{2}a}+\frac{v}{(2n+1)^{2}},\\ 1-\frac{\varphi^{2}\left(\frac{\alpha}{2n+1}\right)}{\varphi^{2}\left(\frac{m\omega+\mu\varpi i+l}{2n+1}\right)}&=1-\frac{\alpha^{2}}{(2n+1)^{2}.\varphi^{2}a}+\frac{v'}{(2n+1)^{2}}.\end{aligned} \]

\origpage{166}
On a de même
\[ \begin{aligned}1-\frac{\alpha^{2}}{(m\omega+\mu\varpi i+k)^{2}}&=1-\frac{\alpha^{2}}{(2n+1)^{2}\left\{\frac{m\omega+\mu\varpi i+k}{(2n+1)}\right\}}=1-\frac{\alpha^{2}}{(2n+1)^{2}a^{2}}+\frac{v}{(2n+1)^{2}},\\ 1-\frac{\alpha^{2}}{(m\omega+\mu\varpi i+l)^{2}}&=1-\frac{\alpha^{2}}{(2n+1)^{2}a^{2}}+\frac{v'}{(2n+1)^{2}}.\end{aligned} \]
\ednote{In the first line the braced quotient in the middle denominator is printed without a square; a square is expected, as in the neighbouring terms.}

En substituant ces valeurs, l'expression de $\psi(m,\mu)-\theta(m,\mu)$ prendra la forme:
\[ \psi(m,\mu)-\theta(m,\mu)=\log\left\{\frac{1-\frac{v}{(2n+1)^{2}}}{1-\frac{v'}{(2n+1)^{2}}}.\frac{1-\frac{v_{1}}{(2n+1)^{2}}}{1-\frac{v'_{1}}{(2n+1)^{2}}}\right\}, \]
les quantités $v$, $v'$, $v_{1}$, $v'_{1}$ ayant toutes zéro pour limite.

On a donc:
\[ \log\left(1-\frac{v}{(2n+1)^{2}}\right)=\frac{v}{(2n+1)^{2}}\ \text{etc.}, \]
et par conséquent:
\[ \psi(m,\mu)-\theta(m,\mu)=\frac{v}{(2n+1)^{2}}. \]

$b)$ Si la limite de la quantité $\frac{m\omega+\mu\varpi i}{2n+1}$ est égale à zéro, on aura:
\[ \begin{aligned}\varphi^{2}\left(\frac{m\omega+\mu\varpi i+k}{2n+1}\right)&=\frac{(m\omega+\mu\varpi i+k)^{2}}{(2n+1)^{2}}+A.\frac{(m\omega+\mu\varpi i+k)^{4}}{(2n+1)^{4}}+....,\\ \varphi^{2}\left(\frac{\alpha}{2n+1}\right)&=\frac{\alpha^{2}}{(2n+1)^{2}}+A'.\frac{\alpha^{4}}{(2n+1)^{4}}+....,\end{aligned} \]
donc:
\[ 1-\frac{\varphi^{2}\left(\frac{\alpha}{2n+1}\right)}{\varphi^{2}\left(\frac{m\omega+\mu\varpi i+k}{2n+1}\right)}=1-\frac{\alpha^{2}+A'.\frac{\alpha^{4}}{(2n+1)^{2}}+.....}{(m\omega+\mu\varpi i+k)^{2}+A.\frac{(m\omega+\mu\varpi i+k)^{4}}{(2n+1)^{2}}}. \]

Si maintenant $m\omega+\mu\varpi i$ ne va pas indéfiniment en augmentant avec $n$, on aura:
\[ 1-\frac{\varphi^{2}\left(\frac{\alpha}{2n+1}\right)}{\varphi^{2}\left(\frac{m\omega+m\varpi i+k}{2n+1}\right)}=1-\frac{\alpha^{2}}{(m\omega+\mu\varpi i+k)^{2}}+\frac{B}{(2n+1)^{2}}, \]
\ednote{Printed $m\varpi i$ in the argument of the denominator function instead of $\mu\varpi i$; an apparent misprint.}
de même:
\[ 1-\frac{\varphi^{2}\left(\frac{\alpha}{2n+1}\right)}{\varphi^{2}\left(\frac{m\omega+\mu\varpi i+k}{2n+1}\right)}=1-\frac{\alpha^{2}}{(m\omega+\mu\varpi i+l)^{2}}+\frac{C}{(2n+1)^{2}}, \]
\ednote{The argument of the denominator function is printed with $+k$ (the glyph is blotted) where $+l$ is expected, as in the right-hand side.}
donc dans ce cas:

\origpage{167}
\[ \psi(m,\mu)-\theta(m,\mu)=\log\left\{\frac{1-\frac{B'}{(2n+1)^{2}}}{1-\frac{C'}{(2n+1)^{2}}}\right\}, \]
$B'$ et $C'$ ayant des limites finies, ou bien:
\[ \psi(m,\mu)-\theta(m,\mu)=\frac{D}{(2n+1)^{2}}, \]
la limite de $D$ étant également une quantité finie.

Si au contraire la quantité $m\omega+\mu\varpi i$ augmente indéfiniment avec $n$, on a:
\[ \begin{aligned}\frac{1-\frac{\varphi^{2}\left(\frac{\alpha}{2n+1}\right)}{\varphi^{2}\left(\frac{m\omega+\mu\varpi i+k}{2n+1}\right)}}{1-\frac{\alpha^{2}}{(m\omega+\mu\varpi i+k)^{2}}}&=1+\frac{\alpha^{2}}{(2n+1)^{2}}.\left\{\frac{A-A'\frac{\alpha^{2}}{(m\omega+\mu\varpi i+k)^{2}}}{1+A\left(\frac{m\omega+\mu\varpi i+k}{2n+1}\right)^{2}}\right\}\\ &\times\frac{1}{1-\frac{\alpha^{2}}{(m\omega+\mu\varpi i+k)^{2}}};\end{aligned} \]
or les quantités $\frac{1}{m\omega+\mu\varpi i+k}$, $\frac{m\omega+\mu\varpi i+k}{2n+1}$ ont zéro pour limite; donc la quantité précédente sera de la forme:
\[ 1+\frac{\alpha^{2}}{(2n+1)^{2}}.A'', \]
$A''$ ayant une quantité finie pour limite. En changeant $k$ en $l$, et désignant la valeur correspondante de $A''$ par $A''_{1}$, la valeur de $\psi(m,\mu)-\theta(m,\mu)$ deviendra:
\[ \psi(m,\mu)-\theta(m,\mu)=\log\left\{\frac{1+\frac{\alpha^{2}}{(2n+1)^{2}}A''}{1+\frac{\alpha^{2}}{(2n+1)^{2}}A''_{1}}\right\}=\frac{\alpha^{2}(A''-A''_{1})}{(2n+1)^{2}}+\frac{v}{(2n+1)^{2}}. \]
Maintenant la limite de $A''$ est la même que celle de $A$; or il est clair, que cette dernière limite est indépendante de $k$, $m$, $\mu$ (elle est en effet égale au coefficient de $\alpha^{4}$ dans le développement de $\Phi^{2}\alpha$). Donc on aura:
\[ A''=M+v, \]
et en changeant $k$ en $l$:
\[ A''_{1}=M+v', \]
d'où $A''-A''_{1}=v-v'=v$. Donc $A''-A''_{1}$ a zéro pour limite, et par conséquent on a:
\[ \psi(m,\mu)-\theta(m,\mu)=\frac{v}{(2n+1)^{2}}. \]

\origpage{168}

Donc nous avons démontré, qu'en faisant
\[ \psi(m,\mu)-\theta(m,\mu)=\frac{A_{m,\mu}}{(2n+1)^{2}}, \tag{167.} \]
la limite de $A_{m,\mu}$ sera égale à zéro toutes les fois que $m\omega+\mu\varpi i$ augmente indéfiniment avec $n$, et qu'elle sera égale à une quantité finie dans le cas contraire.

\subsection*{29.}

Cela posé, considérons la quantité
\[ \overset{n}{\underset{1}{\Sigma}}_{\mu}\psi(m,\mu). \]
En substituant la valeur de $\psi(m,\mu)$, il viendra:
\[ \overset{n}{\underset{1}{\Sigma}}_{\mu}\psi(m,\mu)=\overset{n}{\underset{1}{\Sigma}}_{\mu}\theta(m,\mu)+\frac{1}{(2n+1)^{2}}.\overset{n}{\underset{1}{\Sigma}}_{\mu}(A_{m,\mu}). \tag{168.} \]
Soit $\nu$ le plus grand nombre entier, contenu dans $\sqrt{n}$, on peut faire:
\[ \begin{aligned}\overset{n}{\underset{1}{\Sigma}}_{\mu}A_{m,\mu}&=A_{m,1}+A_{m,2}+.....+A_{m,\nu}\\ &+A_{m,\nu+1}+A_{m,\nu+2}+.....+A_{m,n}.\end{aligned} \]

Or, d'après la nature des quantités $A_{m,\mu}$, la somme contenue dans la première ligne sera égale à $\nu.A_{m}$, et la seconde égale à $A'_{m}(n-\nu-1)$, où $A_{m}$ est une quantité finie et $A'_{m}$ une quantité, qui a zéro pour limite, donc:
\[ \overset{n}{\underset{1}{\Sigma}}_{\mu}A_{m,\mu}=\nu.A_{m}+(n-\nu-1)A'_{m}=(2n+1).B_{m}, \]
où
\[ B_{m}=\frac{\nu}{2n+1}A_{m}+\frac{n-\nu-1}{2n+1}A'_{m}. \]
Donc la quantité $B_{m}$ a zéro pour limite, en remarquant, que $\nu$ ne surpasse pas $\sqrt{n}$.

Par là l'expression de $\overset{n}{\underset{1}{\Sigma}}_{\mu}\psi(m,\mu)$ se change en:
\[ \overset{n}{\underset{1}{\Sigma}}_{\mu}\psi(m,\mu)=\overset{n}{\underset{1}{\Sigma}}_{\mu}\theta(m,\mu)+\frac{B_{m}}{2n+1}. \tag{169.} \]

Pour avoir la limite de $\overset{n}{\underset{1}{\Sigma}}_{\mu}\theta(m,\mu)$, j'écris
\[ \overset{n}{\underset{1}{\Sigma}}_{\mu}\theta(m,\mu)=\overset{\infty}{\underset{1}{\Sigma}}_{\mu}\theta(m,\mu)-\overset{\infty}{\underset{n+1}{\Sigma}}_{\mu}\theta(m,\mu)=\overset{\infty}{\underset{1}{\Sigma}}_{\mu}\theta(m,\mu)-\overset{\infty}{\underset{1}{\Sigma}}_{\mu}\theta(m,\mu+n). \]
Or on peut trouver la valeur de $\overset{\infty}{\underset{1}{\Sigma}}_{\mu}\theta(m,\mu+n)$ comme suit.

\origpage{169}

On a:
\[ \begin{aligned}\theta(m.\,\mu+n)&=\log\left\{\frac{1-\frac{\alpha^{2}}{[m\omega+(\mu+n)\varpi i+k]^{2}}}{1-\frac{\alpha^{2}}{[m\omega+(\mu+n)\varpi i+l]^{2}}}\right\}\\ &=\alpha^{2}\left\{\frac{1}{[m\omega+(\mu+n)\varpi i+l]^{2}}-\frac{1}{[m\omega+(\mu+n)\varpi i+k]^{2}}\right\}\\ &-\tfrac{1}{2}\alpha^{4}\left\{\frac{1}{[m\omega+(\mu+n)\varpi i+l]^{2}}-\frac{1}{[m\omega+(\mu+n)\varpi i+k]^{2}}\right\}\\ &+\ \text{etc.}\end{aligned} \]
\ednote{Printed $\theta(m.\,\mu+n)$ with a full stop instead of a comma between the arguments; an apparent misprint.}
De là on tire:
\[ \overset{\mu}{\underset{1}{\Sigma}}_{\mu}\theta(m,\mu+n)=\frac{\alpha^{2}}{n}.\Sigma\frac{1}{n}.\theta\left(\frac{\mu}{n}\right)-\frac{\alpha^{4}}{2n^{3}}.\Sigma\frac{1}{n}.\theta_{1}\left(\frac{\mu}{n}\right)+...., \]
où
\[ \begin{gathered}\theta\left(\frac{\mu}{n}\right)=\frac{1}{\left(\frac{m\omega+l}{n}+\varpi i+\frac{\mu}{n}\varpi i\right)^{2}}-\frac{1}{\left(\frac{m\omega+k}{n}+\varpi i+\frac{\mu}{n}\varpi i\right)^{2}},\\ \theta_{1}\left(\frac{\mu}{n}\right)=\frac{1}{\left(\frac{m\omega+l}{n}+\varpi i+\frac{\mu}{n}\varpi i\right)^{4}}-\frac{1}{\left(\frac{m\omega+k}{n}+\varpi i+\frac{\mu}{n}\varpi i\right)^{4}},\\ \text{etc.}\end{gathered} \]
Or on sait que la limite de
\[ \overset{\mu}{\underset{1}{\Sigma}}_{\mu}\frac{1}{n}\theta\left(\frac{\mu}{n}\right)\ \text{est égale à}\ \int_{0}^{x}\theta(x)\,\partial x, \]
donc:
\[ \begin{gathered}\overset{\mu}{\underset{1}{\Sigma}}_{\mu}\frac{1}{n}\theta\left(\frac{\mu}{n}\right)=\int_{0}^{x}\theta(x)\,\partial x+v,\\ \overset{\mu}{\underset{1}{\Sigma}}_{\mu}\frac{1}{n}\theta_{1}\left(\frac{\mu}{n}\right)=\int_{0}^{x}\theta_{1}(x)\,\partial x+v_{1},\\ \text{etc},\end{gathered} \]
et par conséquent, en substituant:
\[ \overset{\mu}{\underset{1}{\Sigma}}_{\mu}\theta(m,\mu+n)=\frac{\alpha^{2}}{n}.\int_{0}^{x}\theta(x)-\frac{\alpha^{4}}{2n^{3}}.\int_{0}^{x}\theta_{1}(x)+....+\frac{v\alpha^{2}}{n}-\frac{v_{1}\alpha^{3}}{2n^{3}}+.... \]
\ednote{The integrals are printed without the differential $\partial x$, and the last fraction has $\alpha^{3}$ where $\alpha^{4}$ is expected; both apparent misprints, reproduced as printed.}
or
\[ \begin{gathered}\theta(x)=\frac{1}{\left(\frac{m\omega+l}{n}+\varpi i+x\varpi i\right)^{2}}-\frac{1}{\left(\frac{m\omega+k}{n}+\varpi i+x\varpi i\right)^{2}}\\ \text{etc.}\end{gathered} \]
donc on aura:

\origpage{170}
\[ \begin{aligned}\int_{0}^{x}\theta(x)\,\partial x&=\frac{1}{\varpi i}.\left\{\frac{1}{\left(\frac{m\omega+l}{n}+\varpi i+x\varpi i\right)}-\frac{1}{\left(\frac{m\omega+k}{n}+\varpi i+x\varpi i\right)}\right\}\\ &-\frac{1}{\varpi i}.\left\{\frac{1}{\left(\frac{m\omega+l}{n}+\varpi i\right)}-\frac{1}{\left(\frac{m\omega+k}{n}+\varpi i\right)}\right\}\\ &=\frac{1}{\varpi i}.\frac{k-l}{n}.\frac{1}{\left(\frac{m\omega+l}{n}+\varpi i+x\varpi i\right)\left(\frac{m\omega+k}{n}+\varpi i+x\varpi i\right)}\\ &-\frac{1}{\varpi i}.\frac{k-l}{n}.\frac{1}{\left(\frac{m\omega+l}{n}+\varpi i\right)\left(\frac{m\omega+k}{n}+\varpi i\right)}.\end{aligned} \]

La limite de cette expression de $\displaystyle\int_{0}^{x}\theta(x)\,\partial x$ est zéro pour une valeur quelconque de $x$. De même on trouvera, que la limite de $\displaystyle\int_{0}^{x}\theta_{1}(x)\,\partial x$ est zéro, donc:
\[ \begin{aligned}\overset{\mu}{\underset{1}{\Sigma}}_{\mu}\theta(m,\mu+n)&=\frac{\alpha^{2}}{n}.v-\frac{\alpha^{4}}{2n^{3}}.v'+\frac{\alpha^{6}}{3n^{5}}.v''-....\\ &=\frac{\alpha^{2}}{2n+1}.\left\{v-\frac{\alpha^{2}}{2n^{2}}v'+\frac{\alpha^{4}}{3n^{4}}v''-....\right\}\frac{2n+1}{n}=\frac{v}{2n+1},\end{aligned} \]
donc aussi, en faisant $\mu=\infty$:
\[ \overset{\infty}{\underset{1}{\Sigma}}_{\mu}\theta(m,\mu+n)=\frac{v}{2n+1}, \]
d'où:
\[ \overset{n}{\underset{1}{\Sigma}}_{\mu}\psi(m,\mu+n)=\overset{\infty}{\underset{1}{\Sigma}}_{\mu}\theta(m,\mu)-\frac{v}{2n+1},\ \text{et} \]
\[ \overset{n}{\underset{1}{\Sigma}}_{\mu}\psi(m,\mu)=\overset{\infty}{\underset{1}{\Sigma}}_{\mu}\theta(m,\mu)+\frac{v_{m}}{2n+1}, \tag{170.} \]
$v_{m}$ ayant zéro pour limite. De là on tire
\[ \overset{n}{\underset{1}{\Sigma}}_{m}\overset{n}{\underset{1}{\Sigma}}_{\mu}\psi(m,\mu)=\overset{n}{\underset{1}{\Sigma}}_{m}\left\{\overset{\infty}{\underset{1}{\Sigma}}_{\mu}\theta(m,\mu)\right\}+\overset{n}{\underset{1}{\Sigma}}\frac{v_{m}}{2n+1}. \]
En prenant la limite des deux membres et remarquant, que
\[ \overset{n}{\underset{1}{\Sigma}}\frac{v_{m}}{2n+1}=\frac{v_{1}+v_{2}+.....v_{m}}{2n+1}=v, \]
on aura:
\[ \text{lim.}\overset{n}{\underset{1}{\Sigma}}_{m}\overset{n}{\underset{1}{\Sigma}}_{\mu}\psi(m,\mu)=\overset{\infty}{\underset{1}{\Sigma}}_{m}\left\{\overset{\infty}{\underset{1}{\Sigma}}_{\mu}\theta(m,\mu)\right\}. \tag{171.} \]
En remettant les valeurs de $\psi(m,\mu)$ et $\theta(m,\mu)$, et passant des logarithmes aux nombres, on en tire:

\origpage{171}
\[ \text{lim.}\overset{n}{\underset{1}{\Pi}}_{m}\overset{n}{\underset{1}{\Pi}}_{\mu}\left\{\frac{1-\frac{\varphi^{2}\left(\frac{\alpha}{2n+1}\right)}{\varphi^{2}\left(\frac{m\omega+\mu\varpi i+k}{2n+1}\right)}}{1-\frac{\varphi^{2}\left(\frac{\alpha}{2n+1}\right)}{\varphi^{2}\left(\frac{m\omega+\mu\varpi i+l}{2n+1}\right)}}\right\}=\overset{\infty}{\underset{1}{\Pi}}_{m}\left\{\overset{\infty}{\underset{1}{\Pi}}_{\mu}\left(\frac{1-\frac{\alpha^{2}}{(m\omega+\mu\varpi i+k)^{2}}}{1-\frac{\alpha^{2}}{(m\omega+\mu\varpi i+l)^{2}}}\right)\right\}. \tag{172.} \]

Par une analyse toute semblable à la précédente, mais plus simple, on trouvera de même:
\[ \text{lim.}\overset{n}{\underset{1}{\Pi}}_{m}\left\{\frac{1-\frac{\varphi^{2}\left(\frac{\alpha}{2n+1}\right)}{\varphi^{2}\left(\frac{m\omega+k}{2n+1}\right)}}{1-\frac{\varphi^{2}\left(\frac{\alpha}{2n+1}\right)}{\varphi^{2}\left(\frac{m\omega+l}{2n+1}\right)}}\right\}=\overset{\infty}{\underset{1}{\Pi}}_{m}\left\{\frac{1-\frac{\alpha^{2}}{(m\omega+k)^{2}}}{1-\frac{\alpha^{2}}{(m\omega+l)^{2}}}\right\}, \tag{173.} \]
\[ \text{lim.}\overset{n}{\underset{1}{\Pi}}_{\mu}\left\{\frac{1-\frac{\varphi^{2}\left(\frac{\alpha}{2n+1}\right)}{\varphi^{2}\left(\frac{\mu\varpi i+k}{2n+1}\right)}}{1-\frac{\varphi^{2}\left(\frac{\alpha}{2n+1}\right)}{\varphi^{2}\left(\frac{\mu\varpi i+k}{2n+1}\right)}}\right\}=\overset{\infty}{\underset{1}{\Pi}}_{\mu}\left\{\frac{1-\frac{\alpha^{2}}{(\mu\varpi i+k)^{2}}}{1-\frac{\alpha^{2}}{(\mu\varpi i+l)^{2}}}\right\}. \tag{174.} \]
\ednote{In the left side of the last equation the lower fraction is printed with $k$ in the argument of the denominator function, as the upper one; presumably $l$ was meant, as in the right side. An apparent misprint.}

\subsection*{30.}

Maintenant rien n'est plus facile, que de trouver les valeurs de $\Phi\alpha$, $f\alpha$, $F\alpha$.

Considérons d'abord la première formule (130.). On a:
\[ \begin{aligned}e^{2}c^{2}\Phi^{2}\left(\frac{m\omega+\mu\varpi i}{2n+1}\right)&=-\frac{1}{\varphi^{2}\left(\frac{m\omega+\mu\varpi i}{2n+1}-\frac{\omega}{2}-\frac{\varpi}{2}i\right)}\\ &=-\frac{1}{\varphi^{2}\left(\frac{(n-m+\tfrac{1}{2})\omega+(n-\mu+\tfrac{1}{2})\varpi i}{2n+1}\right)},\end{aligned} \]
donc:

\origpage{172}
\[ \begin{aligned}\overset{n}{\underset{1}{\Pi}}_{m}\overset{\mu}{\underset{1}{\Pi}}\left\{\frac{1-\frac{\varphi^{2}\beta}{\varphi^{2}\left(\frac{m\omega+\mu\varpi i}{2n+1}\right)}}{1+e^{2}c^{2}.\varphi^{2}\left(\frac{m\omega+\mu\varpi i}{2n+1}\right)\varphi^{2}\beta}\right\}&=\frac{\overset{n}{\underset{1}{\Pi}}_{m}\overset{\mu}{\underset{1}{\Pi}}.\left\{1-\frac{\varphi^{2}\beta}{\varphi^{2}\left(\frac{m\omega+\mu\varpi i}{2n+1}\right)}\right\}}{\overset{n}{\underset{1}{\Pi}}_{m}\overset{\mu}{\underset{1}{\Pi}}.\left\{1-\frac{\varphi^{2}\beta}{\varphi^{2}\left(\frac{(n-m+\tfrac{1}{2})\omega+(n-\mu+\tfrac{1}{2})\varpi i}{2n+1}\right)}\right\}}\\ &=\overset{n}{\underset{1}{\Pi}}_{m}\overset{\mu}{\underset{1}{\Pi}}_{\mu}\left\{\frac{1-\frac{\varphi^{2}\beta}{\varphi^{2}\left(\frac{m\omega+\mu\varpi i}{2n+1}\right)}}{1-\frac{\varphi^{2}\beta}{\varphi^{2}\left(\frac{m\omega+\mu\varpi i+\frac{\omega}{2}+\frac{\varpi}{2}i}{2n+1}\right)}}\right\}.\end{aligned} \]

Cela posé, si l'on fait $\beta=\frac{\alpha}{2n+1}$ et qu'on suppose $n$ infini, il viendra, en faisant usage des formules (172.), (173.), (174.), et remarquant que $\alpha$ est la limite de $(2n+1)\varphi\left(\frac{\alpha}{2n+1}\right)$ et égale à $\alpha$:
\[ \begin{aligned}\Phi\alpha&=\alpha\overset{\infty}{\underset{1}{\Pi}}_{m}\left\{1-\frac{\alpha^{2}}{(m\omega)^{2}}\right\}.\overset{\infty}{\underset{1}{\Pi}}\left\{1+\frac{\alpha^{2}}{(\mu\varpi)^{2}}\right\}\\ &\times\overset{\infty}{\underset{1}{\Pi}}_{m}\left\{\overset{\infty}{\underset{1}{\Pi}}_{\mu}\left(\frac{1-\frac{\alpha^{2}}{(m\omega+\mu\varpi i)^{2}}}{1-\frac{\alpha^{2}}{[(m-\tfrac{1}{2})\omega+(\mu-\tfrac{1}{2})\varpi i]^{2}}}\right).\overset{\infty}{\underset{1}{\Pi}}_{\mu}\left(\frac{1-\frac{\alpha^{2}}{(m\omega-\mu\varpi i)^{2}}}{1-\frac{\alpha^{2}}{[(m-\tfrac{1}{2})\omega-(\mu-\tfrac{1}{2})\varpi i]^{2}}}\right)\right\}.\end{aligned} \tag{175.} \]

Les deux formules (130.) donneront de la même manière, en faisant $\beta=\frac{\alpha}{2n+1}$, et remarquant, que $f(0)=1$, $F(0)=1$:
\[ \begin{aligned}f\alpha&=\overset{\infty}{\underset{1}{\Pi}}_{m}\left\{\overset{\infty}{\underset{1}{\Pi}}_{\mu}\left(\frac{1-\frac{\alpha^{2}}{[(m-\tfrac{1}{2})\omega+\mu\varpi i]^{2}}}{1-\frac{\alpha^{2}}{[(m-\tfrac{1}{2})\omega+(\mu-\tfrac{1}{2})\varpi i]^{2}}}\right)\left(\frac{1-\frac{\alpha^{2}}{[(m-\tfrac{1}{2})\omega-\mu\varpi i]^{2}}}{1-\frac{\alpha^{2}}{[(m-\tfrac{1}{2})\omega-(\mu-\tfrac{1}{2})\varpi i]^{2}}}\right)\right\}\\ &\times\overset{\infty}{\underset{0}{\Pi}}_{\mu}\left\{1-\frac{\alpha^{2}}{(m+\tfrac{1}{2})^{2}\omega^{2}}\right\},\end{aligned} \tag{176.} \]
\[ \begin{aligned}F\alpha&=\overset{\infty}{\underset{0}{\Pi}}_{\mu}\left\{1+\frac{\alpha^{2}}{(\mu+\tfrac{1}{2})\varpi^{2}}\right\}.\overset{\infty}{\underset{1}{\Pi}}_{m}\left\{\overset{\infty}{\underset{1}{\Pi}}_{\mu}\left(\frac{1-\frac{\alpha^{2}}{[m\omega+(\mu-\tfrac{1}{2})\varpi i]^{2}}}{1-\frac{\alpha^{2}}{[(m-\tfrac{1}{2})\omega+(\mu-\tfrac{1}{2})\varpi i]^{2}}}\right)\left(\frac{1-\frac{\alpha^{2}}{[m\omega+(\mu-\tfrac{1}{2})\varpi i]^{2}}}{1-\frac{\alpha^{2}}{[(m-\tfrac{1}{2})\omega+(\mu-\tfrac{1}{2})\varpi i]^{2}}}\right)\right\}.\end{aligned} \tag{177.} \]
\ednote{The two factors of the last product are printed identically, where the earlier formulas have opposite signs; an apparent misprint, reproduced as printed. In the preceding formulas the factor $(m+\tfrac{1}{2})^{2}\omega^{2}$ of (176.) carries its square, and the product index there is printed $\mu$ although $m$ occurs; the factor $(\mu+\tfrac{1}{2})\varpi^{2}$ of (177.) is printed without a square on the bracket. All reproduced as printed.}

On peut aussi donner une forme réelle aux expressions précédentes comme suit:

\origpage{173}
\[
\begin{aligned}
\Phi\alpha &= \alpha.\overset{\infty}{\underset{1}{\Pi}}_{\mu}\left(1+\frac{\alpha^{2}}{\mu^{2}\varpi^{2}}\right).\overset{\infty}{\underset{1}{\Pi}}_{m}\left(1-\frac{\alpha^{2}}{m^{2}\omega^{2}}\right)\\
&\quad\times\overset{\infty}{\underset{1}{\Pi}}_{m}\overset{\infty}{\underset{1}{\Pi}}_{\mu}\frac{1+\frac{(\alpha+m\omega)^{2}}{\mu^{2}\varpi^{2}}}{1+\frac{[\alpha+(m-\tfrac{1}{2})\omega]^{2}}{\mu^{2}\varpi^{2}}}.\frac{1+\frac{(\alpha-m\omega)^{2}}{\mu^{2}\varpi^{2}}}{1+\frac{[\alpha-(m-\tfrac{1}{2})\omega]^{2}}{\mu^{2}\varpi^{2}}}.\left(\frac{1+\frac{(m-\tfrac{1}{2})^{2}\omega^{2}}{\mu^{2}\varpi^{2}}}{1+\frac{m^{2}\omega^{2}}{\mu^{2}\varpi^{2}}}\right)^{2},
\end{aligned}
\tag{178.}
\]
\[
\begin{aligned}
f\alpha &= \overset{\infty}{\underset{1}{\Pi}}_{m}\left(1-\frac{\alpha^{2}}{(m-\tfrac{1}{2})^{2}\omega^{2}}\right)\\
&\quad\times\overset{\infty}{\underset{1}{\Pi}}_{m}\overset{\infty}{\underset{1}{\Pi}}_{\mu}\frac{1+\frac{[\alpha+(m-\tfrac{1}{2})\omega]^{2}}{\mu^{2}\varpi^{2}}}{1+\frac{[\alpha+(m-\tfrac{1}{2})\omega]^{2}}{(\mu-\tfrac{1}{2})^{2}\varpi^{2}}}.\frac{1+\frac{[\alpha-(m-\tfrac{1}{2})\omega]^{2}}{\mu^{2}\varpi^{2}}}{1+\frac{[\alpha-(m-\tfrac{1}{2})\omega]^{2}}{(\mu-\tfrac{1}{2})^{2}\varpi^{2}}}.\left(\frac{1+\frac{(m-\tfrac{1}{2})^{2}\omega^{2}}{(\mu-\tfrac{1}{2})^{2}\varpi^{2}}}{1+\frac{(m-\tfrac{1}{2})^{2}\omega^{2}}{\mu^{2}\varpi^{2}}}\right)^{2},
\end{aligned}
\tag{179.}
\]
\[
\begin{aligned}
F\alpha &= \overset{\infty}{\underset{0}{\Pi}}_{\mu}\left(1+\frac{\alpha^{2}}{(\mu+\tfrac{1}{2})^{2}\varpi^{2}}\right)\\
&\quad\times\overset{\infty}{\underset{1}{\Pi}}_{m}\overset{\infty}{\underset{1}{\Pi}}_{\mu}\frac{1+\frac{(\alpha+m\omega)^{2}}{(\mu-\tfrac{1}{2})^{2}\varpi^{2}}}{1+\frac{[\alpha+(m-\tfrac{1}{2})\omega]^{2}}{(\mu-\tfrac{1}{2})^{2}\varpi^{2}}}.\frac{1+\frac{(\alpha-m\omega)^{2}}{(\mu-\tfrac{1}{2})^{2}\varpi^{2}}}{1+\frac{[\alpha-(m-\tfrac{1}{2})\omega]^{2}}{(\mu-\tfrac{1}{2})^{2}\varpi^{2}}}.\left(\frac{1+\frac{(m-\tfrac{1}{2})^{2}\omega^{2}}{(\mu-\tfrac{1}{2})^{2}\varpi^{2}}}{1+\frac{m^{2}\omega^{2}}{(\mu-\tfrac{1}{2})^{2}\varpi^{2}}}\right)^{2}.
\end{aligned}
\tag{180.}
\]

Ces transformations s'opèrent aisément au moyen de la formule:
\[
\begin{aligned}
\left(1-\frac{\alpha^{2}}{(a+bi)^{2}}\right)\left(1-\frac{\alpha^{2}}{(a-bi)^{2}}\right) &= \left(1+\frac{\alpha}{a+bi}\right)\left(1+\frac{\alpha}{a-bi}\right)\left(1-\frac{\alpha}{a+bi}\right)\left(1-\frac{\alpha}{a-bi}\right)\\
&= \frac{(\alpha+a)^{2}+b^{2}}{a^{2}+b^{2}}.\frac{(\alpha-a)^{2}+b^{2}}{a^{2}+b^{2}} = \left(1+\frac{(\alpha+a)^{2}}{b^{2}}\right)\left(1+\frac{(\alpha-a)^{2}}{b^{2}}\right).\frac{1}{\left(1+\frac{a^{2}}{b^{2}}\right)^{2}}.
\end{aligned}
\]

\subsection*{31.}

Dans ce qui précède nous sommes parvenus à deux espèces d'expressions des fonctions $\Phi\alpha$, $f\alpha$, $F\alpha$: les unes donnent ces fonctions décomposées en fractions partielles, dont la totalité forme des séries infinies doubles, les autres donnent ces mêmes fonctions décomposées en un nombre infini de facteurs, dont chacun est à son tour composé d'une infinité de facteurs.

Or on peut beaucoup simplifier les formules précédentes au moyen des fonctions exponentielles et circulaires. C'est ce que nous allons voir par ce qui suit:

Considérons d'abord les équations (178.), (179.), (180.). En vertu des formules connues, on a:

\origpage{174}
\[
\frac{\sin y}{y} = \overset{\infty}{\underset{1}{\Pi}}_{\mu}\left(1-\frac{y^{2}}{\mu^{2}\pi^{2}}\right);\ \cos y = \overset{\infty}{\underset{1}{\Pi}}_{\mu}\left(1-\frac{y^{2}}{(\mu-\tfrac{1}{2})^{2}\pi^{2}}\right);
\]
donc:
\[
\overset{\infty}{\underset{1}{\Pi}}_{\mu}\frac{\left(1-\frac{z^{2}}{\mu^{2}\pi^{2}}\right)}{\left(1-\frac{y^{2}}{(\mu-\tfrac{1}{2})\pi^{2}}\right)} = \frac{\sin z}{z\cos y};\ \overset{\infty}{\underset{1}{\Pi}}_{\mu}\left(\frac{1-\frac{z^{2}}{(\mu-\tfrac{1}{2})\pi^{2}}}{1-\frac{y^{2}}{(\mu-\tfrac{1}{2})\pi^{2}}}\right) = \frac{\cos z}{\cos y}.
\]
\ednote{In both products the bracket $(\mu-\tfrac{1}{2})$ under $\pi^{2}$ is printed without the square (the denominator of the first product, and the numerator and denominator of the second); an apparent misprint for $(\mu-\tfrac{1}{2})^{2}\pi^{2}$, reproduced as printed.}

En vertu de cette formule il est clair, que les expressions de $\Phi\alpha$, $f\alpha$, $F\alpha$ peuvent être mises sous la forme:
\[
\begin{aligned}
\Phi\alpha &= \frac{\varpi}{\pi}.\frac{\sin\left(\alpha\frac{\pi}{\varpi}i\right)}{i}.\overset{\infty}{\underset{1}{\Pi}}_{m}\left(1-\frac{\alpha^{2}}{m^{2}\omega^{2}}\right)\\
&\quad\times\overset{\infty}{\underset{1}{\Pi}}_{m}\left(\frac{\sin(\alpha+m\omega)\frac{\pi}{\varpi}i.\sin(\alpha-m\omega)\frac{\pi}{\varpi}i.\cos^{2}(m-\tfrac{1}{2})\omega\frac{\pi}{\varpi}i}{\cos[\alpha+(m-\tfrac{1}{2})\omega]\frac{\pi}{\varpi}i.\cos[\alpha-(m-\tfrac{1}{2})\omega]\frac{\pi}{\varpi}i.\sin^{2}m\omega\frac{\pi}{\varpi}i}.\frac{\left(m\omega\frac{\pi}{\varpi}i\right)^{2}}{(\alpha+m\omega)(\alpha-m\omega).\frac{\pi^{2}}{\omega^{2}}i^{2}}\right),
\end{aligned}
\]
\ednote{In the denominator of the second fraction the factor $\frac{\pi^{2}}{\omega^{2}}$ is printed with $\omega$, where $\varpi$ is expected (compare the factor $\frac{\pi}{\varpi}$ elsewhere in the line); an apparent misprint, reproduced as printed.}
\[
\begin{aligned}
f\alpha &= \overset{\alpha}{\underset{1}{\Pi}}_{m}\left(1-\frac{\alpha^{2}}{(m-\tfrac{1}{2})^{2}\omega^{2}}\right)\\
&\quad\times\overset{\infty}{\underset{1}{\Pi}}_{m}\left(\text{tang}[\alpha+(m-\tfrac{1}{2})\omega]\frac{\pi}{\varpi}i.\text{tang}[\alpha-(m-\tfrac{1}{2})\omega]\frac{\pi}{\varpi}i.\cot^{2}(m-\tfrac{1}{2})\omega\frac{\pi}{\varpi}i.\frac{(m-\tfrac{1}{2})^{2}\omega^{2}}{\alpha^{2}-(m-\tfrac{1}{2})^{2}\omega^{2}}\right),
\end{aligned}
\]
\ednote{In the first line of the expression for $f\alpha$ the upper limit of the product is printed as $\alpha$, where $\infty$ is expected; an apparent misprint, reproduced as printed.}
\[
F\alpha = \cos\left(\alpha\frac{\pi}{\varpi}i\right).\overset{\infty}{\underset{1}{\Pi}}_{m}\left(\frac{\cos(\alpha+m\omega)\frac{\pi}{\varpi}i.\cos(\alpha-m\omega)\frac{\pi}{\varpi}i.\cos^{2}(m-\tfrac{1}{2})\omega\frac{\pi}{\varpi}i}{\cos[\alpha+(m-\tfrac{1}{2})\omega]\frac{\pi}{\varpi}i.\cos[\alpha-(m-\tfrac{1}{2})\omega]\frac{\pi}{\varpi}i.\cos^{2}m\omega.\frac{\pi}{\varpi}i}\right).
\]

On trouvera des expressions réelles, en substituant au lieu des fonctions circulaires leurs expressions en fonctions exponentielles.

On a:
\[
\begin{gathered}
\sin(a-b).\sin(a+b) = \sin^{2}a-\sin^{2}b,\\
\cos(a+b).\cos(a-b) = \cos^{2}a-\sin^{2}b,
\end{gathered}
\]
donc:
\[
\frac{\sin(\alpha+m\omega)\frac{\pi}{\varpi}i.\sin(\alpha-m\omega)\frac{\pi}{\varpi}i}{\sin^{2}m\omega\frac{\pi}{\varpi}i} = -\left\{1-\frac{\sin^{2}\alpha\frac{\pi}{\varpi}i}{\sin^{2}m\omega\frac{\pi}{\varpi}i}\right\},
\]
\[
\frac{\cos[\alpha+(m-\tfrac{1}{2})\omega]\frac{\pi}{\varpi}i.\cos[\alpha-(m-\tfrac{1}{2})\omega]\frac{\pi}{\varpi}i}{\cos^{2}(m-\tfrac{1}{2})\omega\frac{\pi}{\varpi}i} = 1-\frac{\sin^{2}\alpha\frac{\pi}{\varpi}i}{\cos^{2}(m-\tfrac{1}{2})\omega\frac{\pi}{\varpi}i},
\]

\origpage{175}
\[
\text{tang}[\alpha+(m-\tfrac{1}{2})\omega]\frac{\pi}{\varpi}i.\text{tang}[\alpha-(m-\tfrac{1}{2})\omega]\frac{\pi}{\varpi}i.\cot^{2}(m-\tfrac{1}{2})\omega\frac{\pi}{\varpi}i = -\frac{1-\frac{\sin^{2}\alpha\frac{\pi}{\varpi}i}{\sin^{2}(m-\tfrac{1}{2})\omega\frac{\pi}{\varpi}i}}{1-\frac{\sin^{2}\alpha\frac{\pi}{\varpi}i}{\cos^{2}(m-\tfrac{1}{2})\omega\frac{\pi}{\varpi}i}}.
\]

D'après cela et en remarquant, que
\[
\frac{m^{2}\omega^{2}}{\alpha^{2}-m^{2}\omega^{2}} = -\frac{1}{1-\frac{\alpha^{2}}{m^{2}\omega^{2}}}\ \text{et}
\]
\[
\frac{(m-\tfrac{1}{2})^{2}\omega^{2}}{\alpha^{2}-(m-\tfrac{1}{2})^{2}\omega^{2}} = -\frac{1}{1-\frac{\alpha^{2}}{(m-\tfrac{1}{2})^{2}\omega^{2}}},
\]
il est clair, qu'on aura:
\[
\Phi\alpha = \frac{\pi}{\varpi}.\frac{\sin\left(\frac{\alpha}{\varpi}\pi i\right)}{i}.\overset{\infty}{\underset{1}{\Pi}}_{m}.\frac{1-\frac{\sin^{2}\alpha\frac{\pi}{\varpi}i}{\sin^{2}m\omega\frac{\pi}{\varpi}i}}{1-\frac{\sin^{2}\alpha\frac{\pi}{\varpi}i}{\cos^{2}(m-\tfrac{1}{2})\omega\frac{\pi}{\varpi}i}},
\tag{181.}
\]
\ednote{The factor before the sine is printed as the fraction with $\pi$ over $\varpi$, whereas the preceding expression has $\varpi$ over $\pi$; an apparent misprint, reproduced as printed.}
\[
f\alpha = \overset{\infty}{\underset{0}{\Pi}}_{m}.\frac{1-\frac{\sin^{2}\alpha\frac{\pi}{\varpi}i}{\sin^{2}(m+\tfrac{1}{2})\omega\frac{\pi}{\varpi}i}}{1-\frac{\sin^{2}\alpha\frac{\pi}{\varpi}i}{\cos^{2}(m+\tfrac{1}{2})\omega\frac{\pi}{\varpi}i}},
\tag{182.}
\]
\[
F\alpha = \cos\left(\frac{\alpha}{\varpi}\pi i\right).\overset{\infty}{\underset{1}{\Pi}}_{m}.\frac{1-\frac{\sin^{2}\frac{\alpha}{\varpi}\pi i}{\cos^{2}m\frac{\omega}{\varpi}\pi i}}{1-\frac{\sin^{2}\frac{\alpha}{\varpi}\pi i}{\cos^{2}(m-\tfrac{1}{2})\frac{\omega}{\varpi}\pi i}}.
\tag{183.}
\]

\origpage{176}

En substituant au lieu des cosinus et sinus d'arcs imaginaires leurs valeurs en quantités exponentielles, ces formules deviendront:
\[
\Phi\alpha = \tfrac{1}{2}.\frac{\varpi}{\pi}\left(h^{\frac{\alpha}{\varpi}\pi}-h^{-\frac{\alpha}{\varpi}\pi}\right).\overset{\infty}{\underset{1}{\Pi}}_{m}.\frac{1-\left(\frac{h^{\frac{\alpha}{\varpi}\pi}-h^{-\frac{\alpha}{\varpi}\pi}}{h^{m\frac{\omega}{\varpi}\pi}-h^{-m\frac{\omega}{\varpi}\pi}}\right)^{2}}{1+\left(\frac{h^{\frac{\alpha}{\varpi}\pi}-h^{-\frac{\alpha}{\varpi}\pi}}{h^{(m-\frac{1}{2})\frac{\omega}{\varpi}\pi}+h^{-(m-\frac{1}{2})\frac{\omega}{\varpi}\pi}}\right)^{2}},
\tag{184.}
\]
\[
f\alpha = \overset{\infty}{\underset{0}{\Pi}}_{m}.\frac{1-\left(\frac{h^{\frac{\alpha}{\varpi}\pi}-h^{-\frac{\alpha}{\varpi}\pi}}{h^{(m+\frac{1}{2})\frac{\omega}{\varpi}\pi}-h^{-(m+\frac{1}{2})\frac{\omega}{\varpi}\pi}}\right)^{2}}{1+\left(\frac{h^{\frac{\alpha}{\varpi}\pi}-h^{-\frac{\alpha}{\varpi}\pi}}{h^{(m+\frac{1}{2})\frac{\omega}{\varpi}\pi}+h^{-(m+\frac{1}{2})\frac{\omega}{\varpi}\pi}}\right)^{2}},
\tag{185.}
\]
\[
F\alpha = \tfrac{1}{2}\left(e^{\frac{\alpha}{\varpi}\pi}+e^{-\frac{\alpha}{\varpi}\pi}\right).\overset{\infty}{\underset{1}{\Pi}}_{m}\frac{1+\left(\frac{h^{\frac{\alpha}{\varpi}\pi}-h^{-\frac{\alpha}{\varpi}\pi}}{h^{m\frac{\omega}{\varpi}\pi}+h^{-m\frac{\omega}{\varpi}\pi}}\right)^{2}}{1+\left(\frac{h^{\frac{\alpha}{\varpi}\pi}-h^{-\frac{\alpha}{\varpi}\pi}}{h^{(m-\frac{1}{2})\frac{\omega}{\varpi}\pi}+h^{-(m-\frac{1}{2})\frac{\omega}{\varpi}\pi}}\right)^{2}},
\tag{186.}
\]
où $h$ est le nombre 2,1718218.....
\ednote{The value of $h$ is printed as 2,1718218; the base of the natural logarithms is 2,7182818; an apparent misprint, reproduced as printed.} \ednote{In the first factor of the last formula the base is printed as $e$ in the sum, where the other exponentials of the page use $h$; an apparent misprint, reproduced as printed.}

On peut encore transformer cette formule de la manière suivante:

Si l'on remplace $\alpha$ par $\alpha i$, on aura les valeurs de $\Phi(\alpha i)$, $f(\alpha i)$, $F(\alpha i)$. En changeant maintenant $c$ en $e$ et $e$ en $c$, les quantités:
\[
\omega;\ \varpi;\ \Phi(\alpha i);\ f(\alpha i);\ F(\alpha i)
\]
se changeront respectivement en:
\[
\varpi;\ \omega;\ i\Phi\alpha;\ F\alpha;\ f\alpha,
\]
donc les formules précédentes donneront:
\[
\Phi\alpha = \frac{\omega}{\pi}.\sin\frac{\alpha\pi}{\omega}.\overset{\infty}{\underset{1}{\Pi}}_{m}.\frac{1+\frac{4\sin^{2}\left(\frac{\alpha\pi}{\omega}\right)}{\left(h^{\frac{m\varpi\pi}{\omega}}-h^{-\frac{m\varpi\pi}{\omega}}\right)^{2}}}{1-\frac{4\sin^{2}\left(\frac{\alpha\pi}{\omega}\right)}{\left(h^{\frac{(2m-1)\varpi\pi}{2\omega}}+h^{-\frac{(2m-1)\varpi\pi}{2\omega}}\right)^{2}}},
\tag{187.}
\]

\origpage{177}
\[
F\alpha = \overset{\infty}{\underset{0}{\Pi}}_{m}\frac{1+\frac{4\sin^{2}\left(\frac{\alpha\pi}{\omega}\right)}{\left(h^{\frac{(2m+1)\varpi\pi}{2\omega}}-h^{-\frac{(2m+1)\varpi\pi}{2\omega}}\right)^{2}}}{1-\frac{4\sin^{2}\left(\frac{\alpha\pi}{\omega}\right)}{\left(h^{\frac{(2m+1)\varpi\pi}{2\omega}}+h^{-\frac{(2m+1)\varpi\pi}{2\omega}}\right)^{2}}},
\tag{188.}
\]
\[
f\alpha = \cos\left(\frac{\alpha\pi}{\omega}\right).\overset{\infty}{\underset{1}{\Pi}}_{m}.\frac{1-\frac{4\sin^{2}\left(\frac{\alpha\pi}{\omega}\right)}{\left(h^{\frac{m\varpi\pi}{\omega}}+h^{-\frac{m\varpi\pi}{\omega}}\right)^{2}}}{1-\frac{4\sin^{2}\left(\frac{\alpha\pi}{\omega}\right)}{\left(h^{\frac{(2m+1)\varpi\pi}{2\omega}}+h^{-\frac{(2m+1)\varpi\pi}{2\omega}}\right)^{2}}},
\tag{189.}
\]

\subsection*{32.}

Considérons maintenant les formules (160.), (161.), (162.).

On a:
\[
\overset{\infty}{\underset{0}{\Sigma}}_{\mu}(-1)^{\mu}.\frac{(2\mu+1)\pi}{y^{2}+(\mu+\tfrac{1}{2})^{2}\pi^{2}} = \frac{2}{h^{y}+h^{-y}},
\]
donc, en faisant
\[
y = [\alpha\pm(m+\tfrac{1}{2})\omega]\frac{\pi}{\varpi}:
\]
\[
\overset{\infty}{\underset{0}{\Sigma}}_{\mu}(-1)^{\mu}.\frac{(2\mu+1)\varpi}{[\alpha\pm(m+\tfrac{1}{2})\omega]^{2}+(\mu+\tfrac{1}{2})^{2}\varpi^{2}} = \frac{2\pi}{\varpi}.\frac{1}{h^{[\alpha\pm(m+\frac{1}{2})\omega]\frac{\pi}{\varpi}}+h^{-[\alpha\pm(m+\frac{1}{2})\omega]\frac{\pi}{\varpi}}}.
\]
En vertu de cette formule il est aisé de voir, que les expressions (153.), (162.) de $\Phi\alpha$ et $F\alpha$ deviendront:
\[
\Phi\alpha = \frac{2}{ec}.\frac{\pi}{\varpi}.\overset{\infty}{\underset{0}{\Sigma}}_{m}(-1)^{m}.\left\{\frac{1}{h^{[\alpha-(m+\frac{1}{2})\omega]\frac{\pi}{\varpi}}+h^{-[\alpha-(m-\frac{1}{2})\omega]\frac{\pi}{\varpi}}}-\frac{1}{h^{[\alpha+(m+\frac{1}{2})\omega]\frac{\pi}{\varpi}}+h^{-[\alpha+(m+\frac{1}{2})\omega]\frac{\pi}{\varpi}}}\right\},
\tag{190.}
\]
\ednote{In the first denominator the second exponent is printed with $m-\tfrac{1}{2}$ where the first has $m+\tfrac{1}{2}$; an apparent misprint, reproduced as printed.}
\[
F\alpha = \frac{2}{c}.\frac{\pi}{\varpi}.\overset{\infty}{\underset{0}{\Sigma}}_{m}\left\{\frac{1}{h^{[\alpha-(m+\frac{1}{2})\omega]\frac{\pi}{\varpi}}+h^{-[\alpha-(m+\frac{1}{2})\omega]\frac{\pi}{\varpi}}}+\frac{1}{h^{[\alpha+(m+\frac{1}{2})\omega]\frac{\pi}{\varpi}}+h^{-[\alpha+(m+\frac{1}{2})\omega]\frac{\pi}{\varpi}}}\right\}.
\tag{191.}
\]

Les expressions précédentes de $\Phi\alpha$, $F\alpha$, peuvent être mises encore sous beaucoup d'autres formes: je vais rappeler les plus remarquables. D'abord en réunissant les termes du second membre, on trouvera:

\origpage{178}
\[
\Phi\alpha = \frac{2}{ec}.\frac{\pi}{\varpi}.\overset{\infty}{\underset{0}{\Sigma}}_{m}(-1)^{m}.\left\{\frac{\left(h^{\frac{\alpha\pi}{\varpi}}-h^{-\frac{\alpha\pi}{\varpi}}\right)\left(h^{(m+\frac{1}{2})\frac{\omega\pi}{\varpi}}-h^{-(m+\frac{1}{2})\frac{\omega\pi}{\varpi}}\right)}{h^{\frac{2\alpha\pi}{\varpi}}+h^{-\frac{2\alpha\pi}{\varpi}}+h^{(2m+1)\frac{\omega\pi}{\varpi}}+h^{-(2m+1)\frac{\omega\pi}{\varpi}}}\right\},
\tag{192.}
\]
\[
F\alpha = \frac{2}{c}.\frac{\pi}{\varpi}.\overset{\infty}{\underset{0}{\Sigma}}_{m}.\left\{\frac{\left(h^{\frac{\alpha\pi}{\varpi}}+h^{-\frac{\alpha\pi}{\varpi}}\right)\left(h^{(m+\frac{1}{2})\frac{\omega\pi}{\varpi}}+h^{-(m+\frac{1}{2})\frac{\omega\pi}{\varpi}}\right)}{h^{\frac{2\alpha\pi}{\varpi}}+h^{-\frac{2\alpha\pi}{\varpi}}+h^{(2m+1)\frac{\omega\pi}{\varpi}}+h^{-(2m+1)\frac{\omega\pi}{\varpi}}}\right\}.
\tag{193.}
\]
Si, pour abréger, on suppose:
\[
h^{\frac{\alpha\pi}{\varpi}} = \varepsilon\ \text{et}\ h^{\frac{\omega\pi}{\varpi}} = r^{2},
\tag{194.}
\]
ces formules, en développant le second membre, deviendront:
\[
\Phi\alpha = \frac{2}{ec}.\frac{\pi}{\varpi}\left(\varepsilon-\frac{1}{\varepsilon}\right).\left\{\frac{r-\frac{1}{r}}{r^{2}+\varepsilon^{2}+\frac{1}{\varepsilon^{2}}+\frac{1}{r^{2}}}-\frac{r^{3}-\frac{1}{r^{3}}}{r^{6}+\varepsilon^{2}+\frac{1}{\varepsilon^{2}}+\frac{1}{r^{6}}}+\frac{r^{5}-\frac{1}{r^{5}}}{r^{10}+\varepsilon^{2}+\frac{1}{\varepsilon^{2}}+\frac{1}{r^{10}}}-\ldots\right\},
\tag{195.}
\]
\[
F\alpha = \frac{2}{c}.\frac{\pi}{\varpi}.\left(\varepsilon+\frac{1}{\varepsilon}\right).\left\{\frac{r+\frac{1}{r}}{r^{2}+\varepsilon^{2}+\frac{1}{\varepsilon^{2}}+\frac{1}{r^{2}}}+\frac{r^{3}+\frac{1}{r^{3}}}{r^{6}+\varepsilon^{2}+\frac{1}{\varepsilon^{2}}+\frac{1}{r^{6}}}+\frac{r^{5}+\frac{1}{r^{5}}}{r^{10}+\varepsilon^{2}+\frac{1}{\varepsilon^{2}}+\frac{1}{r^{10}}}-\ldots\right\}.
\tag{196.}
\]

En mettant $\alpha i$ au lieu de $\alpha$ dans les formules (192.), (193.), changeant ensuite $c$ en $e$ et $e$ en $c$, et remarquant, que les quantités
\[
\omega,\ \varpi,\ \Phi(\alpha i),\ F(\alpha i),\ h^{\alpha i\frac{\pi}{\varpi}}-h^{-\alpha i\frac{\pi}{\varpi}},\ h^{\alpha i\frac{\pi}{\varpi}}+h^{\alpha i\frac{\pi}{\varpi}},
\]
\ednote{In the last quantity the sign of the second exponent is printed without a minus sign; an apparent misprint for $h^{-\alpha i\frac{\pi}{\varpi}}$, reproduced as printed.}
se changeront respectivement en:
\[
\varpi,\ \omega,\ i\Phi(\alpha),\ f\alpha,\ 2i.\sin\alpha\frac{\pi}{\omega},\ 2\cos\alpha\frac{\pi}{\omega},
\]
il viendra:
\[
\Phi\alpha = \frac{4}{ec}.\frac{\pi}{\omega}.\overset{\infty}{\underset{0}{\Sigma}}_{m}(-1)^{m}.\left\{\frac{\sin\frac{\alpha\pi}{\omega}.\left(h^{(m+\frac{1}{2})\frac{\varpi\pi}{\omega}}-h^{-(m+\frac{1}{2})\frac{\varpi\pi}{\omega}}\right)}{h^{(2m+1)\frac{\varpi\pi}{\omega}}+2\cos2\alpha\frac{\pi}{\omega}+h^{-(2m+1)\frac{\varpi\pi}{\omega}}}\right\},
\tag{197.}
\]
\[
f\alpha = \frac{4}{e}.\frac{\pi}{\omega}.\overset{\infty}{\underset{0}{\Sigma}}_{m}.\left\{\frac{\cos\frac{\alpha\pi}{\omega}.\left(h^{(m+\frac{1}{2})\frac{\varpi\pi}{\omega}}+h^{-(m+\frac{1}{2})\frac{\varpi\pi}{\omega}}\right)}{h^{(2m+1)\frac{\varpi\pi}{\omega}}+2\cos2\alpha\frac{\pi}{\omega}+h^{-(2m+1)\frac{\varpi\pi}{\omega}}}\right\}.
\tag{198.}
\]
En faisant pour abréger
\[
h^{\frac{\varpi\pi}{2\omega}} = \varrho,
\tag{199.}
\]

\origpage{179}
et développant, on obtiendra:
\[
\varphi\left(\alpha\frac{\omega}{2}\right) = \frac{4}{ec}.\frac{\pi}{\omega}.\sin\left(\alpha\frac{\pi}{2}\right).\left\{\frac{\varrho-\frac{1}{\varrho}}{\varrho^{2}+2\cos(\alpha\pi)+\frac{1}{\varrho^{2}}}-\frac{\varrho^{3}-\frac{1}{\varrho^{3}}}{\varrho^{6}+2\cos(\alpha\pi)+\frac{1}{\varrho^{6}}}+\frac{\varrho^{5}-\frac{1}{\varrho^{5}}}{\varrho^{10}+2\cos(\alpha\pi)+\frac{1}{\varrho^{10}}}-....\right\},
\tag{200.}
\]
\[
f\left(\alpha\frac{\omega}{2}\right) = \frac{4}{e}.\frac{\pi}{\omega}.\cos\left(\alpha\frac{\pi}{2}\right).\left\{\frac{\varrho+\frac{1}{\varrho}}{\varrho^{2}+2\cos(\alpha\pi)+\frac{1}{\varrho^{2}}}+\frac{\varrho^{3}+\frac{1}{\varrho^{3}}}{\varrho^{6}+2\cos(\alpha\pi)+\frac{1}{\varrho^{6}}}+\frac{\varrho^{5}+\frac{1}{\varrho^{5}}}{\varrho^{10}+2\cos(\alpha\pi)+\frac{1}{\varrho^{10}}}+....\right\}.
\tag{201.}
\]

En substituant dans les formules (190.), (191.) au lieu de $h^{\alpha\frac{\pi}{\varpi}}$ et $h^{\frac{\omega\pi}{2\varpi}}$ leurs valeurs $\varepsilon$ et $r$, il viendra:
\[
\Phi\alpha = \frac{2}{ec}.\frac{\pi}{\varpi}.\overset{\infty}{\underset{0}{\Sigma}}_{m}(-m).\left\{\frac{1}{\varepsilon.r^{-(2m+1)}+\varepsilon^{-1}.r^{2m+1}}-\frac{1}{\varepsilon.r^{2m+1}+\varepsilon^{-1}.r^{-(2m+1)}}\right\},
\tag{202.}
\]
\ednote{The factor is printed as $(-m)$ where $(-1)^{m}$ is expected; an apparent misprint, reproduced as printed.}
\[
F\alpha = \frac{2}{c}.\frac{\pi}{\varpi}.\overset{\infty}{\underset{0}{\Sigma}}_{m}\left\{\frac{1}{\varepsilon.r^{-2m-1}+\varepsilon^{-1}.r^{2m+1}}+\frac{1}{\varepsilon.r^{2m+1}+\varepsilon^{-1}.r^{-2m-1}}\right\}.
\tag{203.}
\]

En supposant maintenant $\alpha<\frac{\omega}{2}$, on aura:
\[
\frac{1}{\varepsilon.r^{-2m-1}+\varepsilon^{-1}.r^{2m+1}} = \frac{\varepsilon^{-1}.r^{-2m-1}}{1+\varepsilon^{2}.r^{-4m-2}} = \varepsilon.r^{-2m-1}-\varepsilon^{3}.r^{-6m-3}+\varepsilon^{5}.r^{-10m-5}-....,
\]
\ednote{In the first fraction the numerator is printed $\varepsilon^{-1}.r^{-2m-1}$ where $\varepsilon.r^{-2m-1}$ is expected (the expansion beside it begins with $\varepsilon.r^{-2m-1}$); an apparent misprint, reproduced as printed.}
\[
\frac{1}{\varepsilon.r^{2m+1}+\varepsilon^{-1}.r^{-2m-1}} = \frac{\varepsilon^{-1}.r^{-2m-1}}{1+\varepsilon^{-2}.r^{-4m-2}} = \varepsilon^{-1}.r^{-2m-1}-\varepsilon^{-3}.r^{-6m-3}+\varepsilon^{-5}.r^{-10m-5}-....,
\]
donc:
\[
\begin{aligned}
\Phi\alpha &= \frac{2}{ec}.\frac{\pi}{\varpi}.\left\{\overset{\infty}{\underset{0}{\Sigma}}_{m}(-1)^{m}\left\{(\varepsilon-\varepsilon^{-1}).r^{-2m-1}-(\varepsilon^{3}-\varepsilon^{-3}).r^{-6m-3}+(\varepsilon^{5}-\varepsilon^{-5})\,r^{-10m-5}-....\right\}\right\}\\
&= \frac{2}{ec}.\frac{\pi}{\varpi}.\left\{(\varepsilon-\varepsilon^{-1}).\overset{\infty}{\underset{0}{\Sigma}}_{m}(-1)^{m}.r^{-2m-1}-(\varepsilon^{3}-\varepsilon^{-3}).\overset{\infty}{\underset{0}{\Sigma}}_{m}(-1)^{m}.r^{-6m-3}+....\right\}.
\end{aligned}
\]
Or:
\[
\overset{\infty}{\underset{0}{\Sigma}}_{m}(-1)^{m}.r^{-2m-1} = r^{-1}-r^{-3}+r^{-5}-.... = \frac{r^{-1}}{1+r^{-2}} = \frac{r}{r^{2}+1},
\]
\[
\overset{\infty}{\underset{0}{\Sigma}}_{m}(-1)^{m}.r^{3m-3} = r^{-3}-r^{-6}+r^{-9}-.... = \frac{r^{-3}}{1+r^{-6}} = \frac{r^{3}}{r^{6}+1}
\]
\ednote{The exponent in the summand is printed $3m-3$ and the series exponents as $-3$, $-6$, $-9$; the sum is evidently meant to be that of $(-1)^{m}.r^{-6m-3}$, with series $r^{-3}-r^{-9}+r^{-15}-\ldots$; apparent misprints, reproduced as printed.}
etc., donc:
\[
\Phi\alpha = \frac{2}{ec}.\frac{\pi}{\varpi}.\left\{\frac{\varepsilon-\varepsilon^{-1}}{r+r^{-1}}-\frac{\varepsilon^{3}-\varepsilon^{-3}}{r^{3}+r^{-3}}+\frac{\varepsilon^{5}-\varepsilon^{-5}}{r^{5}+r^{-5}}-.....\right\}.
\tag{204.}
\]
De la même manière on trouvera:
\[
F\alpha = \frac{2}{c}.\frac{\pi}{\varpi}.\left\{\frac{\varepsilon+\varepsilon^{-1}}{r-r^{-1}}-\frac{\varepsilon^{3}+\varepsilon^{-3}}{r^{3}-r^{-3}}+\frac{\varepsilon^{5}+\varepsilon^{-5}}{r^{5}-r^{-5}}-.....\right\}.
\tag{205.}
\]

\origpage{180}
En mettant $\alpha\frac{\varpi}{2}i$ au lieu de $\alpha$, et changeant ensuite $e$ en $c$ et $c$ en $e$,
\[
\varepsilon;\ \varpi;\ \Phi\left(\alpha\frac{\varpi}{2}i\right);\ F\left(\alpha\frac{\varpi}{2}i\right);\ r;\ \varepsilon^{m}+\varepsilon^{-m};\ \varepsilon^{m}-\varepsilon^{-m}
\]
se changent en:
\[
\varpi;\ \omega;\ i\varphi\left(\alpha\frac{\omega}{2}\right);\ f\left(\alpha\frac{\omega}{2}\right);\ \varrho;\ 2\cos\left(m\alpha\frac{\pi}{2}\right);\ 2i\sin\left(m\alpha\frac{\pi}{2}\right);
\]
donc:
\[
\varphi\left(\alpha\frac{\omega}{2}\right) = \frac{4}{ec}.\frac{\pi}{\omega}.\left\{\frac{\sin\left(\alpha\frac{\pi}{2}\right)}{\varrho+\frac{1}{\varrho}}-\frac{\sin\left(3\alpha\frac{\pi}{2}\right)}{\varrho^{3}+\frac{1}{\varrho^{3}}}+\frac{\sin\left(5\alpha\frac{\pi}{2}\right)}{\varrho^{5}+\frac{1}{\varrho^{5}}}-....\right\},
\tag{206.}
\]
\[
f\left(\alpha\frac{\omega}{2}\right) = \frac{4}{e}.\frac{\pi}{\omega}.\left\{\frac{\cos\left(\alpha\frac{\pi}{2}\right)}{\varrho-\frac{1}{\varrho}}-\frac{\cos\left(3\alpha\frac{\pi}{2}\right)}{\varrho^{3}-\frac{1}{\varrho^{3}}}+\frac{\cos\left(5\alpha\frac{\pi}{2}\right)}{\varrho^{5}-\frac{1}{\varrho^{5}}}-....\right\}.
\tag{207.}
\]

Ces quatre dernières formules offrent d'expressions très simples des fonctions $\Phi\alpha$, $f\alpha$, $F\alpha$. Par les différentiations et intégrations on peut en déduire une foule d'autres plus ou moins remarquables.

\subsection*{33.}

Dans le cas, $e=c$, les formules précédentes prennent une forme plus simple, à cause de la relation $\omega=\varpi$, qui a lieu dans ce cas. Soit pour plus de simplicité $e=c=1$. On a:
\[
r = h^{\frac{\omega\pi}{2\varpi}} = h^{\frac{\pi}{2}},\ \varrho = h^{\frac{\varpi\pi}{2\omega}} = h^{\frac{\pi}{2}},
\]
donc en substituant, et faisant dans (204.), (205.) $\alpha=\alpha\frac{\varpi}{2}$, il vient:
\[
\varphi\left(\alpha\frac{\omega}{2}\right) = 2\frac{\pi}{\omega}\left\{\frac{h^{\frac{\alpha\pi}{2}}-h^{-\frac{\alpha\pi}{2}}}{h^{\frac{\pi}{2}}+h^{-\frac{\pi}{2}}}-\frac{h^{\frac{3\alpha\pi}{2}}-h^{-\frac{3\alpha\pi}{2}}}{h^{\frac{3\pi}{2}}+h^{-\frac{3\pi}{2}}}+\frac{h^{\frac{5\alpha\pi}{2}}-h^{-\frac{5\alpha\pi}{2}}}{h^{\frac{5\pi}{2}}+h^{-\frac{5\pi}{2}}}-.....\right\},
\]
\[
F\left(\alpha\frac{\omega}{2}\right) = 2\frac{\pi}{\omega}\left\{\frac{h^{\frac{\alpha\pi}{2}}+h^{-\frac{\alpha\pi}{2}}}{h^{\frac{\pi}{2}}-h^{-\frac{\pi}{2}}}-\frac{h^{\frac{3\alpha\pi}{2}}+h^{-\frac{3\alpha\pi}{2}}}{h^{\frac{3\pi}{2}}-h^{-\frac{3\pi}{2}}}+\frac{h^{\frac{5\alpha\pi}{2}}+h^{-\frac{5\alpha\pi}{2}}}{h^{\frac{5\pi}{2}}-h^{-\frac{5\pi}{2}}}-.....\right\},
\]
\[
\varphi\left(\alpha\frac{\omega}{2}\right) = \frac{4\pi}{\omega}\left\{\sin\left(\alpha\frac{\pi}{2}\right).\frac{h^{\frac{\pi}{2}}}{1+h^{\pi}}-\sin\left(3\alpha\frac{\pi}{2}\right).\frac{h^{\frac{3\pi}{2}}}{1+h^{3\pi}}+\sin\left(5\alpha\frac{\pi}{2}\right).\frac{h^{\frac{5\pi}{2}}}{1+h^{5\pi}}-....\right\},
\]
\[
f\left(\alpha\frac{\omega}{2}\right) = \frac{4\pi}{\omega}\left\{\cos\left(\alpha\frac{\pi}{2}\right).\frac{h^{\frac{\pi}{2}}}{h^{\pi}-1}-\cos\left(3\alpha\frac{\pi}{2}\right).\frac{h^{\frac{3\pi}{2}}}{h^{3\pi}-1}+\sin\left(5\alpha\frac{\pi}{2}\right).\frac{h^{\frac{5\pi}{2}}}{h^{5\pi}-1}-....\right\}.
\]
\ednote{In the last display the third term has a sine, where the pattern of the first two terms requires a cosine; an apparent misprint, reproduced as printed.}

\origpage{181}
Les fonctions $\varphi$, $f$, $F$ sont déterminées par les équations
\[
\begin{gathered}
\alpha.\frac{\omega}{2} = \int_{0}^{x}\frac{\partial x}{\sqrt{(1-x^{4})}};\ \frac{\omega}{2} = \int_{0}^{1}\frac{\partial x}{\sqrt{(1-x^{4})}};\\
x = \varphi\left(\alpha\frac{\omega}{2}\right) = \sqrt{(1-x^{2})} = f\left(\alpha\frac{\omega}{2}\right) = \sqrt{(1+x^{2})} = F\left(\alpha\frac{\omega}{2}\right).
\end{gathered}
\]

Si dans les deux dernières formules on fait $\alpha=0$, et qu'on remarque, qu'alors la valeur de $\dfrac{\varphi\left(\alpha\frac{\omega}{2}\right)}{\sin\left(\alpha\frac{\pi}{2}\right)}$ est $=\frac{\omega}{\pi}$, et celle de $\dfrac{\sin\left(m\alpha\frac{\pi}{2}\right)}{\sin\left(\alpha\frac{\pi}{2}\right)}=m$, on trouvera:
\[
\frac{\omega}{2} = \frac{\pi}{2}.\left\{\frac{h^{\frac{\pi}{2}}}{h^{\pi}-1}-\frac{h^{\frac{3\pi}{2}}}{h^{3\pi}-1}+\frac{h^{\frac{5\pi}{2}}}{h^{5\pi}-1}-\ldots\right\} = \int_{0}^{1}\frac{\partial x}{\sqrt{(1-x^{4})}},
\]
\[
\frac{\omega^{2}}{4} = \pi^{2}.\left\{\frac{h^{\frac{\pi}{2}}}{h^{\pi}+1}-3.\frac{h^{\frac{3\pi}{2}}}{h^{3\pi}+1}+5.\frac{h^{\frac{5\pi}{2}}}{h^{5\pi}+1}-\ldots\right\} = \left(\int_{0}^{1}\frac{\partial x}{\sqrt{(1-x^{4})}}\right)^{2}.
\]

(La continuation dans le cahier suivant.)

\end{document}
